% Shared preamble for the paper. Section files are \input into main.tex and must compile with exactly this.
\documentclass[11pt]{amsart}
\usepackage[utf8]{inputenc}
% T1 encoding omitted: without Latin Modern installed it selects bitmap fonts; the default encoding gives Type 1 Computer Modern.
\usepackage{amsmath,amssymb,amsthm,mathtools}
\usepackage{booktabs,array,longtable}
\usepackage[margin=1.15in]{geometry}
\usepackage[colorlinks=true,linkcolor=black,citecolor=black,urlcolor=black]{hyperref}
\numberwithin{equation}{section}
\newtheorem{theorem}{Theorem}[section]
\newtheorem{lemma}[theorem]{Lemma}
\newtheorem{proposition}[theorem]{Proposition}
\newtheorem{corollary}[theorem]{Corollary}
\theoremstyle{definition}
\newtheorem{definition}[theorem]{Definition}
\newtheorem{hypothesis}[theorem]{Hypothesis}
\newtheorem{observation}[theorem]{Observation}
\theoremstyle{remark}
\newtheorem{remark}[theorem]{Remark}
\newcommand{\R}{\mathbb{R}}
\newcommand{\Z}{\mathbb{Z}}
\newcommand{\N}{\mathbb{N}}
\newcommand{\SU}{\mathrm{SU}}
\newcommand{\U}{\mathrm{U}}
\newcommand{\Tr}{\operatorname{Tr}}
\newcommand{\Cov}{\operatorname{Cov}}
\newcommand{\Var}{\operatorname{Var}}
\newcommand{\cL}{\mathcal{L}}
\newcommand{\cH}{\mathcal{H}}
\newcommand{\cG}{\mathcal{G}}
\newcommand{\E}{\mathbb{E}}
\renewcommand{\P}{\mathbb{P}}
\newcommand{\repo}[1]{\texttt{#1}}

\title[An upper bound on the lattice mass gap at weak coupling]{An upper bound on the lattice mass gap\\ at weak coupling}
\author{Kirtus Justus}
\address{Theoretical Physicist, Port Coquitlam, British Columbia, Canada}
\date{23 November 2026. Working draft; not reviewed by a mathematician. See Section~\ref{sec:provenance}.}
\begin{document}
\begin{abstract}
For $\SU(2)$ lattice gauge theory in four dimensions with the Wilson action, we prove that the mass gap of the transfer matrix, in lattice units, tends to zero or is zero as the coupling tends to its weak limit, for every infinite-volume limit of periodic states along boxes of side $2^k$. This is an upper bound and does not address whether the gap is positive. We also record a conditional reduction of the continuum problem to five lattice hypotheses, of which two are proved, and five attempted routes to the principal one, each of which fails.
\end{abstract}
\maketitle
\tableofcontents

\section{Introduction}\label{sec:intro}

Lattice gauge theory with the Wilson action is reflection positive and has a self-adjoint transfer matrix at every coupling~\cite{OS78,Lu77}. At strong coupling the transfer matrix has a gap above the vacuum, for every compact gauge group~\cite{OS78}. A continuum limit requires the opposite at weak coupling: the gap, measured in lattice units, must close. We prove that it does.

Let $\beta>0$. On the periodic box $\Lambda_L=(\Z/L\Z)^4$ let $\mu_{\beta,L}$ be the probability measure on $\SU(2)$-valued link variables with density proportional to $\exp\bigl(-\beta\sum_p(1-\cos\alpha_p)\bigr)$ with respect to Haar measure, where the sum runs over plaquettes and $\tfrac12\Tr U_p=\cos\alpha_p$, $\alpha_p\in[0,\pi]$. Let $\cL_\beta$ be the set of weak limit points of $\mu_{\beta,2^k}$ as $k\to\infty$. Each $\mu\in\cL_\beta$ is translation invariant and reflection positive, and has a self-adjoint transfer matrix $T_\mu$ with $\|T_\mu\|\le 1$ and invariant vector $\Omega$ (Section~\ref{sec:proof}). Put $m_\mu(\beta)=-\log\|T_\mu|_{\Omega^\perp}\|\in[0,\infty]$.

\begin{theorem}\label{thm:main}
Let $O_p=\beta(1-\cos\alpha_p)$, let $p_0$ be a plaquette in a spatial coordinate plane at time $0$ and $p_n$ its translate by $n$ lattice steps in time. For every $n\ge 1$,
\[
\liminf_{\beta\to\infty}\ \inf_{\mu\in\cL_\beta}\ \frac{\Cov_\mu(O_{p_0},O_{p_n})}{\Var_\mu(O_{p_0})}\ \ge\ \Bigl(\frac{G(n)}{G(0)}\Bigr)^{2},
\]
where $G(n)$ is the correlation of one field-strength component of the free lattice Maxwell field in four dimensions at time separation $n$; $G(0)=\tfrac12$, $G(n)>0$, and $G(n)=\pi^{-2}n^{-4}+O(n^{-6})$.
\end{theorem}

Consequently $\sup_{\mu\in\cL_\beta}m_\mu(\beta)\to 0$ as $\beta\to\infty$.

The limits are taken in this order: the box, then the coupling, then the separation. The statement allows $m_\mu(\beta)=0$.

\subsection*{Relation to the mass gap problem}
The problem posed by Jaffe and Witten~\cite{JW} asks, for every compact simple gauge group, for a quantum Yang--Mills theory on $\R^4$ satisfying axioms at least as strong as those of Osterwalder and Schrader, with a mass gap $\Delta>0$: no spectrum of the Hamiltonian in $(0,\Delta)$. Theorem~\ref{thm:main} is not that statement and is not a step of a proof of it. It bounds a lattice gap from above. It holds, with the same proof, for the gauge group $\U(1)$, whose four-dimensional lattice theory has a massless phase at weak coupling~\cite{Guth80,FS82}. What the mass gap problem needs on the lattice is the opposite inequality at every coupling, Hypothesis~\ref{hyp:A} below, and this paper does not prove it.

\subsection*{What else is in the paper}
The work reported here began as an attempt at the whole problem, written with the expectation that it would fail, in order to record where it fails. It fails at Hypothesis~\ref{hyp:A}. Section~\ref{sec:reduction} states a conditional reduction: five hypotheses on the lattice theories imply limit Schwinger functions that satisfy the Osterwalder--Schrader axioms with a gap. Two of the five are proved for $\SU(2)$: the lattice axioms, and the closing of the gap, which is Theorem~\ref{thm:main}. Both hold for $\U(1)$ as well. Section~\ref{sec:proof} proves Theorem~\ref{thm:main}. Section~\ref{sec:routes} records five routes to Hypothesis~\ref{hyp:A} and the step at which each fails. Section~\ref{sec:shared} lists six properties, each proved or cited, that $\SU(2)$, $\U(1)$ and $\Z_2$ lattice gauge theories share in four dimensions; since $\U(1)$ has no gap at weak coupling, no argument using only these properties can prove Hypothesis~\ref{hyp:A}. Section~\ref{sec:CE} discusses two further hypotheses. Section~\ref{sec:notclaimed} lists what is not claimed. Section~\ref{sec:provenance} states how the paper was produced and checked.

\subsection*{A stronger result by a different method}
During the preparation of this paper, three other working groups of the same research program, using a different method (a Ward identity obtained by varying Haar measure along rooted paths, with heat-kernel estimates), obtained the bound $m_\mu(\beta)\le C\beta^{-1/8}$ for even-periodic limit states, and families of linearly independent gauge-invariant vectors at energies of order $\beta^{-1/8}$~\cite{RepoC4,RepoB4}. Those results imply the last sentence of Theorem~\ref{thm:main} with a rate. This paper does not use them, and its author has not read their proofs. One of those groups also reports that the random constant background of Lemma~\ref{pf:lem-R} vanishes, which turns the inequality of Theorem~\ref{thm:main} into an equality~\cite{RepoThread}; that is not proved here.

\subsection*{Disclosure}
Every statement and proof in this paper was written by sessions of one artificial-intelligence model family (Claude, Anthropic), working as a director with subordinate workers, at the direction of the author, who is accountable for the paper. Each proof was then attacked by other sessions of the same family. Agreement among them is not independent review. Lemma~\ref{pf:lem-R} and the asymptotics of $G$ were checked, at the scope stated in Section~\ref{sec:provenance}, by sessions of a second model lineage (OpenAI Codex). No mathematician has read the paper. It should be read as a claim awaiting review.

\section{What is not claimed}\label{sec:notclaimed}

\begin{enumerate}
\item No statement is made about the existence of a mass gap, on the lattice or in the continuum. Hypothesis~\ref{hyp:A} is open.
\item No continuum theory is constructed. Theorem~\ref{thm:conditional} is conditional on hypotheses that contain the whole difficulty.
\item Theorem~\ref{thm:main} is restricted to limits of periodic states along boxes of side $2^k$. Other boundary conditions, other sequences of boxes and other Gibbs states are not covered.
\item Theorem~\ref{thm:main} gives no rate.
\item The proof of Theorem~\ref{thm:main} uses ten standard facts of measure theory without proof or citation. They are listed in Section~\ref{pf:sec-standard}.
\item Nothing is proved about gauge groups other than $\SU(2)$. The proof of Theorem~\ref{thm:main} uses the group only through error terms of order $\beta^{-1/2}$, and we expect it to extend; we have not written the extension.
\item The numerical statements of Sections~\ref{sec:routes}--\ref{sec:CE} are measurements on small lattices. None is used in a proof.
\item The observations of Sections~\ref{sec:routes} and~\ref{sec:shared} say that particular arguments fail. They do not say that Hypothesis~\ref{hyp:A} is false, or that no argument of a similar kind can succeed.
\item The paper has not been independently reviewed (Section~\ref{sec:provenance}).
\end{enumerate}

\section{Setting and two reductions}\label{sec:reduction}

This section states the lattice setting and two conditional theorems. Theorem~\ref{thm:reduction} treats one two-point function. Theorem~\ref{thm:conditional} carries a family of local fields to the Osterwalder--Schrader axioms. Both theorems are proved here, with one exception: Step~9 of the proof of Theorem~\ref{thm:conditional} uses the reconstruction theorem as cited and a description of the reconstructed space that is recalled, not read. Neither theorem proves a property of four-dimensional Yang--Mills theory: each moves the whole difficulty into its hypotheses. Table~\ref{red:status} at the end of the section gives the status of each hypothesis for $\SU(2)$.

\subsection{The lattice theory}\label{red:setting}

Let $\Lambda_k=(\Z/2^k\Z)^4$ be the periodic box of side $2^k$ and let $\mathcal{P}_k$ be its set of plaquettes. A configuration assigns $U_b\in\SU(2)$ to each oriented edge $b$, with $U_{-b}=U_b^{-1}$. For a plaquette $p$ let $U_p$ be the ordered product of the four edge variables around $p$. The Wilson action at coupling $\beta>0$ is
\begin{equation}\label{red:action}
  A_{\beta,k}(U)=\beta\sum_{p\in\mathcal{P}_k}\Bigl(1-\tfrac12\Tr U_p\Bigr),
\end{equation}
and the periodic Wilson state is the probability measure $Z^{-1}e^{-A_{\beta,k}(U)}\prod_b dU_b$, with $dU_b$ the Haar measure. Let $\cL_\beta$ be the set of weak limit points, as $k\to\infty$, of these states, regarded as measures on $\SU(2)^{E(\Z^4)}$, where $E(\Z^4)$ is the set of edges of $\Z^4$. The space $\SU(2)^{E(\Z^4)}$ is compact and metrizable, so $\cL_\beta$ is not empty. For each $\beta$ we choose one $\mu_\beta\in\cL_\beta$ and write $\langle\cdot\rangle_\beta$ for its expectation. Nothing below depends on the choice beyond the hypotheses stated.

Coordinates are $x=(x_0,x_1,x_2,x_3)$, with $x_0=t$ the time. Let $W_4$ be the hypercubic group: the $384$ signed permutation matrices acting on $\R^4$ and on $\Z^4$. For $c\in\Z$ the reflection in the site hyperplane $\{x_\mu=c\}$ acts on configurations, and $\Theta$ denotes the induced map on functions of the configuration. We write $\theta$ for the reflection $(t,\vec x)\mapsto(-t,\vec x)$ in the plane $\{t=0\}$.

\emph{Reflection positivity and the transfer matrix.} Hypothesis~\ref{hyp:H0} below states reflection positivity of $\mu_\beta$ in site hyperplanes. Given it, the standard construction \cite{OS78} produces the following objects; we cite it and do not repeat it, and \cite{OS78} was not read for this section. Let $\mathcal{A}_+$ be the bounded functions of the edge variables of edges with both endpoints in $\{t\ge0\}$. The form $(F,G)\mapsto\langle\overline{\Theta F}\,G\rangle_\beta$ is positive semidefinite on $\mathcal{A}_+$. Let $\cH_\beta$ be the completion of $\mathcal{A}_+$ modulo its null space, $[F]$ the class of $F$, and $\Omega_\beta=[1]$. Let $\tau_n$ be the translation by $n$ lattice steps in the time direction. Then $T_\beta[F]:=[\tau_1F]$ defines an operator with the following properties: $T_\beta$ is self-adjoint by translation invariance; $\langle v,T_\beta^2v\rangle\ge0$ by reflection positivity in the site plane $\{t=1\}$; $\|T_\beta\|\le1$; and $T_\beta\Omega_\beta=\Omega_\beta$. Site-plane reflection positivity alone does not give $T_\beta\ge0$. For the Wilson action $T_\beta\ge0$ is a theorem of \cite{Lu77}, recalled and not read for this section. This paper does not use it.

\begin{definition}\label{red:def-gap}
The \emph{gap in lattice units} is $m(\beta):=-\log\|T_\beta|_{\Omega_\beta^\perp}\|\in[0,\infty]$. When $0<m(\beta)<\infty$, the \emph{correlation length} is $\xi(\beta):=1/m(\beta)$ and the \emph{unit of length} is
\[
  a(\beta):=m(\beta).
\]
A lattice point $n\in\Z^4$ sits at the physical point $a(\beta)n\in a(\beta)\Z^4$.
\end{definition}

\emph{What this choice of unit does.} The unit of length is the correlation length. In this unit the gap equals $1$ by definition, at every $\beta$. The reduction therefore moves nothing. It does not convert a hard statement into an easier one. It fixes the form in which the difficulty must be met, and its content is its hypotheses. Hypothesis~\ref{hyp:A} at weak coupling, $m(\beta)>0$ for large $\beta$ in infinite volume, is the lattice form of the mass gap problem. Hypothesis~\ref{hyp:B} says that the unit $a(\beta)$ shrinks to zero, so that a limit $\beta\to\infty$ is a continuum limit. The other hypotheses ask that correlations stay regular when lengths are measured in this unit.

\subsection{One two-point function}\label{red:R}

Theorem~\ref{thm:reduction} is stated for abstract data, as in the source. For each $\beta\ge\beta_1$ we are given:
\begin{itemize}
\item a Hilbert space $\cH_\beta$ with a unit vector $\Omega_\beta$;
\item a self-adjoint operator $T_\beta$ with $0<T_\beta\le1$ (that is, $T_\beta\ge0$ and $T_\beta$ injective) and $T_\beta\Omega_\beta=\Omega_\beta$; we write $T_\beta=e^{-H_\beta}$ with $H_\beta\ge0$ self-adjoint;
\item a vector $\psi_\beta\ne0$ with $\psi_\beta\perp\Omega_\beta$.
\end{itemize}
The gap is $m(\beta):=\inf\operatorname{spec}(H_\beta|_{\Omega_\beta^\perp})$, which agrees with Definition~\ref{red:def-gap}. Let $\nu_\beta$ be the spectral measure of $\psi_\beta$ for $H_\beta$. It is a finite positive measure on $[m(\beta),\infty)$ with $\nu_\beta([m(\beta),\infty))=\|\psi_\beta\|^2>0$. For real $n\ge0$ put
\begin{equation}\label{red:C}
  C_\beta(n):=\langle\psi_\beta,e^{-nH_\beta}\psi_\beta\rangle=\int e^{-n\varepsilon}\,d\nu_\beta(\varepsilon).
\end{equation}
In the application $\psi_\beta=(O-\langle O\rangle_\beta)\Omega_\beta$ for a bounded gauge-invariant function $O$ of the edges of one time slice near the origin, and $C_\beta(n)$ is its connected correlation at time separation $n$. It is a lattice correlation only at integer $n$; at other $n$ it is defined by the functional calculus. To apply the theorem to the Wilson transfer matrix one needs $0<T_\beta$, which is the theorem of \cite{Lu77}. Theorem~\ref{thm:conditional} does not need it.

\begin{hypothesis}[(A): gap at every coupling]\label{hyp:A}
$m(\beta)>0$ for all $\beta\ge\beta_1$.
\end{hypothesis}

\begin{hypothesis}[(B): criticality at weak coupling]\label{hyp:B}
$m(\beta)\to0$ as $\beta\to\infty$.
\end{hypothesis}

Under (A) and (B), after enlarging $\beta_1$, $\xi(\beta)=1/m(\beta)\in(1,\infty)$ for $\beta\ge\beta_1$.

\begin{hypothesis}[(C): bounded energy at the correlation time]\label{hyp:C}
There are $c\in(0,1)$ and $D<\infty$ such that $C_\beta(c\,\xi)\le D\,C_\beta(\xi)$ for all $\beta\ge\beta_1$.
\end{hypothesis}

The source first stated (C) with $c=\tfrac12$. Every $c\in(0,1)$ gives the same proof, and we state that form. In Theorem~\ref{thm:conditional}, Hypothesis~(A) is called H1 and Hypothesis~(B) is called H2. There $m(\beta)$ is the gap of Definition~\ref{red:def-gap}.

Let $E:=\varepsilon\,\xi$ be the energy in units of the gap. Let $\tilde\nu_\beta$ be the image of $\nu_\beta$ under $\varepsilon\mapsto\varepsilon\xi$; it is carried by $[1,\infty)$. Define the probability measure $\mu_\beta$ on $[1,\infty)$ by
\begin{equation}\label{red:mu}
  d\mu_\beta(E):=\frac{e^{-E}\,d\tilde\nu_\beta(E)}{\int e^{-E'}\,d\tilde\nu_\beta(E')},
  \qquad
  G_\beta(t):=\frac{C_\beta(t\xi)}{C_\beta(\xi)}\quad(t\ge0).
\end{equation}
Hypothesis~(C) says exactly that $\int e^{(1-c)E}\,d\mu_\beta(E)\le D$.

\begin{theorem}[Theorem R$'$]\label{thm:reduction}
Assume Hypotheses~\ref{hyp:A}, \ref{hyp:B} and~\ref{hyp:C}, with constants $c\in(0,1)$ and $D$. Then:
\begin{enumerate}
\item[(R1)] For every $\beta\ge\beta_1$, $\mu_\beta$ is a probability measure on $[1,\infty)$ and $G_\beta(t)=\int e^{-(t-1)E}\,d\mu_\beta(E)$ for every $t\ge0$.
\item[(R2)] $\mu_\beta((K,\infty))\le D\,e^{-(1-c)K}$ for every $K\ge0$, and $\int E\,d\mu_\beta\le \log D/(1-c)$, for every $\beta\ge\beta_1$.
\item[(R3)] Every sequence $\beta_k\to\infty$ has a subsequence along which $\mu_{\beta}$ converges weakly to a probability measure $\mu$ on $[1,\infty)$, and $G_\beta(t)\to G(t):=\int e^{-(t-1)E}\,d\mu(E)$ for every $t>c$.
\item[(R4)] $G(1)=1$; $0<G(t)\le e^{-(t-1)}$ for $t\ge1$; $\int E\,d\mu\le\log D/(1-c)$; and $G(t)\ge D^{-(t-1)/(1-c)}$ for $t\ge1$.
\item[(R5)] The unit of length $a(\beta)=m(\beta)$ tends to $0$.
\end{enumerate}
\end{theorem}

\begin{proof}
(R1). By (A), $\xi<\infty$ and $\tilde\nu_\beta$ is carried by $[1,\infty)$. The denominator in \eqref{red:mu} equals $C_\beta(\xi)=\int e^{-\xi\varepsilon}\,d\nu_\beta>0$, because $\nu_\beta\ne0$ and the integrand is positive. So $\mu_\beta$ is a probability measure on the closed set $[1,\infty)$. For $t\ge0$,
\[
  \frac{C_\beta(t\xi)}{C_\beta(\xi)}=\frac{\int e^{-tE}\,d\tilde\nu_\beta}{\int e^{-E}\,d\tilde\nu_\beta}=\int e^{-(t-1)E}\,d\mu_\beta(E).
\]
Both integrals are finite at fixed $\beta$ because $\tilde\nu_\beta$ is finite and $e^{-tE}\le1$. Hypothesis~(C) is used only for bounds uniform in $\beta$.

(R2). For $E>K$ we have $e^{-E}\le e^{-(1-c)K}e^{-cE}$. Hence
\[
  \int_{(K,\infty)}e^{-E}\,d\tilde\nu_\beta\le e^{-(1-c)K}\int e^{-cE}\,d\tilde\nu_\beta=e^{-(1-c)K}C_\beta(c\xi)\le D\,e^{-(1-c)K}C_\beta(\xi).
\]
Divide by $C_\beta(\xi)$. For the mean, Jensen's inequality for the convex function $\exp$ gives $\exp\bigl((1-c)\int E\,d\mu_\beta\bigr)\le\int e^{(1-c)E}\,d\mu_\beta=C_\beta(c\xi)/C_\beta(\xi)\le D$.

(R3). By (R2) the family $(\mu_\beta)$ is tight on $[1,\infty)$. Prokhorov's theorem gives a weakly convergent subsequence. Its limit $\mu$ has mass $1$, and it is carried by $[1,\infty)$ because that set is closed. For $t\ge1$ the function $E\mapsto e^{-(t-1)E}$ is bounded and continuous on $[1,\infty)$, so the integrals converge. For $c<t<1$ the integrand $e^{(1-t)E}$ is continuous and uniformly integrable:
\[
  \int_{(K,\infty)}e^{(1-t)E}\,d\mu_\beta\le e^{-(t-c)K}\int e^{(1-c)E}\,d\mu_\beta\le D\,e^{-(t-c)K}.
\]
Truncation at $K$ and weak convergence then give $G_\beta(t)\to G(t)$.

(R4). $G(1)=\mu([1,\infty))=1$. For $t\ge1$ and $E\ge1$ we have $0<e^{-(t-1)E}\le e^{-(t-1)}$, and the integral of a positive function against a probability measure is positive. The function $E\mapsto E$ is continuous and bounded below, so $\int E\,d\mu\le\liminf\int E\,d\mu_\beta\le\log D/(1-c)$ along the subsequence. Jensen's inequality gives $G(t)\ge\exp\bigl(-(t-1)\int E\,d\mu\bigr)\ge D^{-(t-1)/(1-c)}$ for $t\ge1$.

(R5). This is Definition~\ref{red:def-gap} and Hypothesis~\ref{hyp:B}.
\end{proof}

Hypothesis~(B) enters only in (R5). Theorem~\ref{thm:reduction} is Prokhorov's theorem and properties of Laplace transforms. The conclusions ``gap at least $1$'' and ``not zero'' hold by construction: the unit is the gap, and the correlation is normalized at one correlation length.

Every admissible constant in Hypothesis~\ref{hyp:C} satisfies $D\ge e^{1-c}$, so $D\ge\sqrt e$ when $c=\tfrac12$. Indeed $\mu_\beta$ is a probability measure carried by $E\ge1$, so $D\ge\int e^{(1-c)E}\,d\mu_\beta\ge e^{1-c}$.

\emph{What Theorem~\ref{thm:reduction} does not give.} It gives one two-point function of one observable along one time axis, along subsequences. It gives no Euclidean invariance and no interaction: a free massive lattice field satisfies its hypotheses, and the theorem returns a free field. The lowest energy that $\psi_\beta$ sees in the limit can exceed $1$, so the theorem shows a gap of at least $1$, not equal to $1$. The limit can be $\mu=\delta_1$, a pure exponential, which a local relativistic field does not produce at a fixed spatial point. The normalization $C_\beta(\xi)$ tends to zero, and nothing in the theorem makes the normalizations of different observables consistent.

Hypothesis~\ref{hyp:C} is not a statement about locality alone. It also asks that the observable overlap the lightest states, as the following example shows.

\begin{observation}\label{red:twomass}
Let the data be two decoupled free massive lattice scalar fields on $\Z^4$ with masses $m_1$ fixed and $m_2=m_2(\beta)\to0$, and let $\psi_\beta=\varphi_1(0)\Omega_\beta$, the first field at one site. Then Hypotheses~\ref{hyp:A} and~\ref{hyp:B} hold, and Hypothesis~\ref{hyp:C} fails for every $c\in(0,1)$ and every $D<\infty$.
\end{observation}

\begin{proof}
For one free field of mass $m$ the one-particle energies are $2\operatorname{asinh}\bigl(\tfrac12(m^2+\sum_{i=1}^3 4\sin^2(p_i/2))^{1/2}\bigr)\ge E_0(m):=2\operatorname{asinh}(m/2)$, and $\varphi(0)\Omega$ lies in the one-particle space. For the decoupled pair the gap is $m(\beta)=\min(E_0(m_1),E_0(m_2))=E_0(m_2)$, which is positive and tends to $0$. The spectral measure of $\psi_\beta$ is carried by $[E_0(m_1),\infty)$. Hence
\[
  \frac{C_\beta(c\xi)}{C_\beta(\xi)}=\frac{\int e^{-c\xi\varepsilon}\,d\nu_\beta}{\int e^{-\xi\varepsilon}\,d\nu_\beta}\ge e^{(1-c)\xi E_0(m_1)}\to\infty
\]
as $\xi=1/E_0(m_2)\to\infty$.
\end{proof}

The free-field dispersion relation used here is standard; Section~\ref{red:free} records its numerical check.

\subsection{Fields and hypotheses for the full chain}\label{red:S-setting}

Fix a finite family $O_1,\dots,O_r$ of \emph{site-scalar} observables. Each $O_i(n)$ is a bounded, real, gauge-invariant function of the edges with both endpoints within lattice distance $1$ of the site $n$, and it is invariant under $W_4$ acting about $n$. An example is $O(n)=\sum_p\tfrac12\Tr U_p$ over the $24$ plaquettes having $n$ as a corner. Write $\hat O_i(n):=O_i(n)-\langle O_i\rangle_\beta$. A reflection $\theta'$ in any site hyperplane maps $\hat O_i(n)$ to $\hat O_i(\theta'n)$ exactly, with no shift of the base point. That is the reason for site-scalar fields.

Let $N_\beta:=2\lfloor\xi/2\rfloor$ and define the normalization
\[
  Z_i(\beta)^2:=\langle\hat O_i(0)\,\hat O_i(N_\beta e_0)\rangle_\beta .
\]
For $x=(x_1,\dots,x_k)\in(a\Z^4)^k$ and indices $\iota=(i_1,\dots,i_k)$, the \emph{lattice Schwinger functions} are
\begin{equation}\label{red:Sbeta}
  S^\beta_\iota(x):=\prod_{j=1}^k Z_{i_j}(\beta)^{-1}\,\Bigl\langle\prod_{j=1}^k\hat O_{i_j}(x_j/a)\Bigr\rangle_\beta .
\end{equation}
Let ${}^0\mathcal S(\R^{4k})$ be the space of Schwartz functions that vanish, with all their derivatives, wherever $x_i=x_j$ for some $i\ne j$, and let ${}^0\mathcal S_+(\R^{4k})$ be the subspace of functions whose support lies in $\{0<t_1<\dots<t_k\}$ \cite[eq.~(2.25)]{KQR21}. For $f\in{}^0\mathcal S(\R^{4k})$ put
\[
  S^\beta_\iota(f):=a^{4k}\sum_{x\in(a\Z^4)^k}S^\beta_\iota(x)\,f(x).
\]
A configuration with two coinciding points contributes $0$, because $f$ vanishes there. We use the Schwartz norms $|f|_N:=\max_{|\alpha|\le N}\sup_y(1+|y|)^N|\partial^\alpha f(y)|$. The norms in \cite[eq.~(9.18)]{KQR21} were not compared with these; E0$'$ below is stated with ours. For $x\in(\R^4)^k$ let $r(x):=\min_{i<j}|x_i-x_j|$ and $\rho_j(x):=\min_{l\ne j}|x_j-x_l|$.

There is no smoothing flow. Every point carries a local composite field, renormalized multiplicatively by $Z_i(\beta)$ and evaluated only at non-coinciding points, which is what ${}^0\mathcal S$ allows. All short-distance control is placed in Hypothesis~\ref{hyp:H3}.

\begin{hypothesis}[H0: lattice axioms]\label{hyp:H0}
$\mu_\beta$ is invariant under the translations $\Z^4$ and under $W_4$. It is reflection positive for reflections in every site hyperplane $\{x_\mu=c\}$, $c\in\Z$: $\langle\overline{\Theta F}\,F\rangle_\beta\ge0$ for every bounded $F$ depending only on edges in $\{x_\mu\ge c\}$.
\end{hypothesis}

H1 is Hypothesis~\ref{hyp:A} and H2 is Hypothesis~\ref{hyp:B}, both for the gap $m(\beta)$ of Definition~\ref{red:def-gap} of the state $\mu_\beta$.

\begin{hypothesis}[H3: uniform regularity in units of the gap]\label{hyp:H3}
There are constants $\alpha,\gamma,b,D<\infty$, and for each $k$ and $\delta>0$ a constant $L_{k,\delta}$, such that for all $\beta\ge\beta_1$, all $k$ and all $\iota$:
\begin{enumerate}
\item[(a)] $Z_i(\beta)>0$;
\item[(b)] (growth) $|S^\beta_\iota(x)|\le\alpha\gamma^k(k!)^b\prod_{j=1}^k\max(1,\rho_j(x)^{-D})$ for $x\in(a\Z^4)^k$ with distinct points;
\item[(c)] (equicontinuity) $|S^\beta_\iota(x)-S^\beta_\iota(y)|\le L_{k,\delta}|x-y|$ for $x,y\in(a\Z^4)^k$ with all pairwise distances at least $\delta$.
\end{enumerate}
\end{hypothesis}

Increasing $D$ weakens (b), so we take $D\ge1$.

\begin{hypothesis}[H4: asymptotic rotation invariance, one rotation]\label{hyp:H4}
For one angle $\vartheta$ with $\vartheta/\pi$ irrational, let $R_\vartheta$ be the rotation by $\vartheta$ in the $(x_1,x_2)$-plane, and let $[y]_a$ be the nearest point of $a\Z^4$ to $y$ (ties broken lexicographically). For every $k$, $\iota$, $\delta>0$ and $R<\infty$,
\[
  \sup\bigl|S^\beta_\iota([R_\vartheta x]_a)-S^\beta_\iota([x]_a)\bigr|\to0\qquad(\beta\to\infty),
\]
the supremum taken over configurations in the ball of radius $R$ with pairwise distances at least $\delta$.
\end{hypothesis}

Here $[x]_a$ is taken pointwise. Each point moves by at most $a$, so $|[x]_a-x|\le a\sqrt k$ in $(\R^4)^k$.

\emph{The Osterwalder--Schrader axioms.} We read the axioms in \cite{KQR21}, Section~2.2 and Section~9.1--9.2. We did not read \cite{OS73} or \cite{OS75}. For $f$ on $\R^{4k}$ write $(\Theta f)^*(x):=\overline{f(\theta x_1,\dots,\theta x_k)}$, and for $s\in\R$ write $f_s(x):=f(x_1-se_0,\dots,x_k-se_0)$. A family of distributions $S_\iota$ on ${}^0\mathcal S(\R^{4k})$ satisfies:
\begin{itemize}
\item[E0$'$] (linear growth) if there is an integer $s$ with $|S_\iota(f)|\le\sigma_k|f|_{ks}$ for all $f\in{}^0\mathcal S(\R^{4k})$, where $\sigma_k\le\alpha'(k!)^{b'}$ for constants $\alpha',b'$ \cite[eq.~(9.18)]{KQR21};
\item[E1] (Euclidean covariance) if $S_\iota(f\circ g)=S_\iota(f)$ for every Euclidean motion $g$ of $\R^4$, acting on each point;
\item[E2] (reflection positivity) if $\sum_{\iota,\iota'}S_{(\iota',\iota)}\bigl((\Theta f_{\iota'})^*\otimes f_\iota\bigr)\ge0$ for every finite family $(f_\iota)$ with $f_\iota\in{}^0\mathcal S_+$, where the empty index tuple carries a constant and $S_\emptyset=1$;
\item[E3] (symmetry) if $S_\iota$ is invariant under simultaneous permutation of the points and the indices;
\item[E4] (cluster) if $S_{(\iota',\iota)}\bigl((\Theta f)^*\otimes g_{\lambda\vec v}\bigr)\to S_{\iota'}((\Theta f)^*)\,S_\iota(g)$ as $\lambda\to\infty$, for $f,g\in{}^0\mathcal S_+$ and every spatial vector $\vec v\ne0$; here $g_{\lambda\vec v}$ is $g$ translated by $\lambda\vec v$.
\end{itemize}
In \cite{KQR21}, E4 is named as the Osterwalder--Schrader cluster property for spatial translations (footnote~20 there); the formula displayed is the standard form, written by us. \cite[Section~9.2.1]{KQR21} records that Lemma~8.8 of \cite{OS73} is wrong; the 1975 form with E0$'$ is the one used here.

\begin{theorem}[Reconstruction theorem \cite{OS75}, as read in \cite{KQR21}]\label{red:OS}
Distributions $S_\iota$ satisfying E0$'$ and E1--E3 are the Schwinger functions of a Wightman theory with fields $\varphi_1,\dots,\varphi_r$.
\end{theorem}

We also use the form of the reconstructed objects: the Hilbert space is the completion of finite families in ${}^0\mathcal S_+$ modulo the null space of the E2 form, the vacuum $\Omega$ is the class of the constant $1$, and $e^{-sH}$ acts by $[f]\mapsto[f_s]$ for $s\ge0$. This description is taken from \cite[Section~4]{OS73} as recalled. It was not read.

\begin{theorem}[Theorem S]\label{thm:conditional}
Assume Hypotheses~\ref{hyp:H0}, \ref{hyp:A} (H1), \ref{hyp:B} (H2), \ref{hyp:H3} and~\ref{hyp:H4}. Then there is a sequence $\beta_j\to\infty$ such that, for every $k$ and $\iota$:
\begin{enumerate}
\item $S^{\beta_j}_\iota\to S_\iota$ uniformly on compact sets of configurations with pairwise distances at least $\delta>0$ (through the lattice approximants of Step~1 below), with $S_\iota$ continuous on the set of configurations of pairwise distinct points; and $S^{\beta_j}_\iota(f)\to S_\iota(f):=\int S_\iota f$ for every $f\in{}^0\mathcal S(\R^{4k})$.
\item The $S_\iota$ satisfy the Osterwalder--Schrader axioms in the 1975 form: E0$'$, E1, E2, E3 and E4.
\item $S_{ii}(0,e_0)=1$ for every $i$, so the theory is not empty.
\item By Theorem~\ref{red:OS} there is a Wightman theory with fields $\varphi_1,\dots,\varphi_r$. Its vacuum is the unique translation-invariant vector, and the joint spectrum of $(H,\vec P)$ on $\Omega^\perp$ lies in $\{p^2\ge1,\ p^0>0\}$: the mass gap is at least $1$ in these units.
\end{enumerate}
The theorem asserts nothing about uniqueness of the limit, interaction, short-distance behaviour, or other gauge groups.
\end{theorem}

\subsection{Proof of Theorem~\ref{thm:conditional}}\label{red:S-proof}

By H2, after enlarging $\beta_1$, $a(\beta)\le1$ for $\beta\ge\beta_1$; by H1, $a(\beta)>0$. Fix any sequence $\beta\to\infty$; all subsequences below are taken from it, and $a\to0$ along it.

\emph{Step 1: compactness.} For $x\in(\R^4)^k$ whose rounded points $[x_j]_a$ are distinct, put $\tilde S^\beta_\iota(x):=S^\beta_\iota([x]_a)$. For $\ell\in\N$ let
\[
  K_\ell:=\{x\in(\R^4)^k:\ |x_j|\le\ell\text{ for all }j,\ r(x)\ge1/\ell\}.
\]
If $a\le1/(4\ell)$, then for $x\in K_\ell$ the rounded configuration has pairwise distances at least $1/\ell-2a\ge1/(2\ell)$. By H3(b), $|\tilde S^\beta_\iota|\le\alpha\gamma^k(k!)^b(2\ell)^{Dk}$ on $K_\ell$. By H3(c), for $x,y\in K_\ell$,
\begin{equation}\label{red:asym-eq}
  |\tilde S^\beta_\iota(x)-\tilde S^\beta_\iota(y)|\le L_{k,1/(2\ell)}\bigl(|x-y|+2a\sqrt k\bigr).
\end{equation}

\begin{lemma}\label{red:AA}
Let $K$ be a compact metric space and $f_j:K\to\R$ functions with $\sup|f_j|\le M$ and $|f_j(x)-f_j(y)|\le L\,d(x,y)+\varepsilon_j$, where $\varepsilon_j\to0$. Then a subsequence converges uniformly on $K$ to an $L$-Lipschitz function.
\end{lemma}

\begin{proof}
Let $Q\subset K$ be countable and dense. A diagonal argument gives a subsequence, still written $f_j$, that converges at every point of $Q$. Given $\eta>0$, choose a finite set $Q_\eta\subset Q$ such that every $x\in K$ lies within $\eta$ of some $q\in Q_\eta$. Then $|f_i(x)-f_j(x)|\le2L\eta+\varepsilon_i+\varepsilon_j+\max_{q\in Q_\eta}|f_i(q)-f_j(q)|$, which is less than $(2L+1)\eta$ for $i,j$ large. So the sequence is uniformly Cauchy. Letting $j\to\infty$ in the hypothesis shows that the limit is $L$-Lipschitz.
\end{proof}

Apply Lemma~\ref{red:AA} on each $K_\ell$ with $\varepsilon=2aL_{k,1/(2\ell)}\sqrt k$, and extract diagonally over the countably many triples $(k,\iota,\ell)$. This gives $\beta_j\to\infty$ and functions $S_\iota$ on the set of configurations of distinct points, Lipschitz on each $K_\ell$, with $\tilde S^{\beta_j}_\iota\to S_\iota$ uniformly on each $K_\ell$. Every compact set of configurations with pairwise distances at least $\delta$ lies in some $K_\ell$. Since $\rho_j([x]_a)\ge\rho_j(x)-2a$, H3(b) passes to the limit:
\begin{equation}\label{red:growth-limit}
  |S_\iota(x)|\le\alpha\gamma^k(k!)^b\prod_{j}\max(1,\rho_j(x)^{-D})\quad\text{for distinct points.}
\end{equation}

\emph{Step 2: the distribution property and E0$'$.}

\begin{lemma}[Taylor]\label{red:taylor}
Let $f\in{}^0\mathcal S(\R^{4k})$, $M\ge1$ and $P\ge0$, and let $N\ge\max(M,P)$. If $r(x)\le1$, then
\[
  |f(x)|\le\frac{(4k)^M}{M!}\,2^P\,|f|_N\;r(x)^M\,(1+|x|)^{-P}.
\]
\end{lemma}

\begin{proof}
Choose $i<j$ with $|x_i-x_j|=r(x)$. Let $z$ agree with $x$ except that $z_i=z_j=\tfrac12(x_i+x_j)$. Then $z$ lies in $\{x_i=x_j\}$ and $|x-z|=r(x)/\sqrt2<1$. Put $g(s):=f(z+s(x-z))$ for $s\in[0,1]$. All derivatives of $f$ of order less than $M$ vanish at $z$, so Taylor's formula with integral remainder gives
\[
  f(x)=g(1)=\int_0^1\frac{(1-s)^{M-1}}{(M-1)!}\,g^{(M)}(s)\,ds,
  \qquad
  g^{(M)}(s)=\sum_{|\alpha|=M}\frac{M!}{\alpha!}\,\partial^\alpha f(z+s(x-z))\,(x-z)^\alpha .
\]
Now $|(x-z)^\alpha|\le|x-z|^M\le r(x)^M$ and $\sum_{|\alpha|=M}M!/\alpha!=(4k)^M$. Every point $y$ of the segment satisfies $|y-x|<1$, so $1+|x|\le2(1+|y|)$ and $|\partial^\alpha f(y)|\le|f|_N(1+|y|)^{-P}\le2^P|f|_N(1+|x|)^{-P}$. Combining, $|f(x)|\le\frac{1}{M!}(4k)^M r(x)^M2^P|f|_N(1+|x|)^{-P}$.
\end{proof}

Take $M:=Dk+1$ and $P:=4k+2$. Then $N:=(D+6)k\ge\max(M,P)$. Put
\[
  A_k:=\alpha\gamma^k(k!)^b\,2^{4k+2}\,(4e)^{Dk+1}.
\]
Since $M!\ge(M/e)^M$ and $D\ge1$, we have $(4k)^M/M!\le(4ek/M)^M\le(4e)^{M}$.

\emph{Claim.} Let $x$ have distinct points, and let $S$ be either $S^\beta_\iota$ (for $x$ in the lattice) or $S_\iota$. Then
\begin{equation}\label{red:pointwise}
  |S(x)f(x)|\le A_k\,|f|_{(D+6)k}\,\min(1,r(x))\,(1+|x|)^{-(4k+2)} .
\end{equation}
If $r(x)\ge1$, every $\rho_j\ge1$, so $|S(x)|\le\alpha\gamma^k(k!)^b$ by H3(b) or \eqref{red:growth-limit}, and $|f(x)|\le|f|_N(1+|x|)^{-P}$. If $r(x)<1$, then $\prod_j\max(1,\rho_j^{-D})\le r(x)^{-Dk}$, and Lemma~\ref{red:taylor} gives $|f(x)|\le(4e)^M2^P|f|_Nr(x)^{Dk+1}(1+|x|)^{-P}$. The product is at most the right side of \eqref{red:pointwise}.

Let $I_k:=\int_{\R^{4k}}(1+|y|)^{-(4k+2)}\,dy=|S^{4k-1}|/\bigl(4k(4k+1)\bigr)$ and $I'_k:=\int_{\R^{4k}}(1+|y|)^{-(4k+1)}\,dy$; both are finite. If $a\sqrt{4k}\le1$, then each lattice point $x$ is within distance $1$ of every point $y$ of the cube $x+[0,a)^{4k}$, so $1+|y|\le2(1+|x|)$ and
\begin{equation}\label{red:riemann}
  a^{4k}\sum_{x\in(a\Z^4)^k}(1+|x|)^{-P}\le2^P\int_{\R^{4k}}(1+|y|)^{-P}\,dy .
\end{equation}

\emph{E0$'$.} Integrating \eqref{red:pointwise} gives $|S_\iota(f)|\le A_kI_k|f|_{(D+6)k}$. With $|S^{4k-1}|=2\pi^{2k}/\Gamma(2k)\le3\pi^{2k}$ and $4k(4k+1)\ge20$,
\[
  A_kI_k\le7\alpha\,\Gamma^k(k!)^b,\qquad \Gamma:=16\gamma(4e)^D\pi^2,
\]
and $\Gamma^k\le e^\Gamma k!$. This is E0$'$ with $s=D+6$ and $\sigma_k\le7\alpha e^{\Gamma}(k!)^{b+1}$. For the lattice functions, \eqref{red:pointwise} and \eqref{red:riemann} give
\begin{equation}\label{red:uniform}
  |S^\beta_\iota(f)|\le2^{4k+2}A_kI_k\,|f|_{(D+6)k}\qquad\text{whenever }a\le1/(2\sqrt k).
\end{equation}

\emph{Convergence of $S^{\beta_j}_\iota(f)$.} Fix $f$, $\delta\in(0,1)$ and $R\ge1$. Let $\mathsf A:=\{x:|x|\le R,\ r(x)\ge\delta\}$ and let $\mathsf B$ be its complement. On $\mathsf B$, either $r(x)<\delta$, or $|x|>R$ and then $(1+|x|)^{-(4k+2)}\le(1+R)^{-1}(1+|x|)^{-(4k+1)}$. By \eqref{red:pointwise} and \eqref{red:riemann}, the contributions of $\mathsf B$ to $S^{\beta_j}_\iota(f)$ and to $\int S_\iota f$ are each at most $2^{4k+2}A_k|f|_{(D+6)k}(\delta I_k+(1+R)^{-1}I'_k)$, uniformly in $j$ large. The set $\mathsf A$ is compact and lies in some $K_\ell$, on which $\tilde S^{\beta_j}_\iota\to S_\iota$ uniformly; at lattice points $\tilde S^{\beta_j}_\iota=S^{\beta_j}_\iota$. So $a^{4k}\sum_{x\in\mathsf A}|S^{\beta_j}_\iota-S_\iota||f|\to0$. Finally $S_\iota f$ is continuous on a neighbourhood of $\mathsf A$, and the boundary of $\mathsf A$ has Lebesgue measure zero, so the Riemann sums $a^{4k}\sum_{x\in\mathsf A}S_\iota(x)f(x)$ converge to $\int_{\mathsf A}S_\iota f$. Letting $j\to\infty$, then $\delta\to0$ and $R\to\infty$, gives $S^{\beta_j}_\iota(f)\to S_\iota(f)$. This proves part~1 of the theorem and E0$'$.

\emph{Step 3: translations.} By H0, $S^\beta_\iota(x+an)=S^\beta_\iota(x)$ for $n\in\Z^4$, and $[x+an]_a=[x]_a+an$. Let $y\in\R^4$ and choose $n_j\in\Z^4$ with $|a_jn_j-y|\le a_j$. By \eqref{red:asym-eq}, $|\tilde S^{\beta_j}_\iota(x+y)-\tilde S^{\beta_j}_\iota(x+a_jn_j)|\le L(a_j\sqrt k+2a_j\sqrt k)\to0$ on each $K_\ell$, and $\tilde S^{\beta_j}_\iota(x+a_jn_j)=\tilde S^{\beta_j}_\iota(x)$. Hence $S_\iota(x+y)=S_\iota(x)$.

\emph{Step 4: rotations.} Let $\mathcal G$ be the set of $g\in O(4)$ with $S_\iota(gx)=S_\iota(x)$ for all $k$, $\iota$ and all configurations $x$ of distinct points. It is a subgroup, and it is closed because each $S_\iota$ is continuous on the $O(4)$-invariant open set of distinct configurations. $W_4\subset\mathcal G$: by H0, $\mu_\beta$ is $W_4$-invariant, $a\Z^4$ is $W_4$-invariant and each $O_i$ is a $W_4$-scalar, so $S^\beta_\iota(wx)=S^\beta_\iota(x)$ for lattice $x$ and $w\in W_4$; the rounding $[wx]_a$ differs from $w[x]_a$ by at most $2a\sqrt k$, which costs $O(a)$ by \eqref{red:asym-eq}. $R_\vartheta\in\mathcal G$: $R_\vartheta$ is an isometry, so it maps the ball of radius $R$ with separation $\delta$ to itself, and H4 together with the uniform convergence of Step~1 gives $S_\iota(R_\vartheta x)=\lim\tilde S^{\beta_j}_\iota(R_\vartheta x)=\lim\tilde S^{\beta_j}_\iota(x)=S_\iota(x)$. By Lemma~\ref{red:dense} below, $\mathcal G=O(4)$. With Step~3 and the change of variables $\int S_\iota(x)f(g^{-1}x)\,dx=\int S_\iota(gx)f(x)\,dx$, this is E1.

\begin{lemma}\label{red:dense}
Let $\vartheta/\pi$ be irrational. The closed subgroup $\mathcal K$ of $O(4)$ generated by $W_4$ and $R_\vartheta$ is $O(4)$.
\end{lemma}

\begin{proof}
For $0\le\mu<\nu\le3$ let $L_{\mu\nu}:=E_{\mu\nu}-E_{\nu\mu}$, where $E_{\mu\nu}$ is the matrix unit, and let $R^{\mu\nu}_\varphi:=\exp(\varphi L_{\mu\nu})$, the rotation by $\varphi$ in the $(x_\mu,x_\nu)$-plane. So $R_\vartheta=R^{12}_\vartheta$.

(i) Kronecker. Since $\vartheta/(2\pi)$ is irrational, the multiples $n\vartheta$, $n\in\N$, are dense modulo $2\pi$. So the powers $R_\vartheta^n=R^{12}_{n\vartheta}$ are dense in $\{R^{12}_\varphi:\varphi\in\R\}$, and $\mathcal K$, being closed, contains every $R^{12}_\varphi$.

(ii) Conjugation. For $\mu<\nu$ let $\sigma\in W_4$ be a permutation matrix with $\sigma e_1=e_\mu$ and $\sigma e_2=e_\nu$. Then $\sigma R^{12}_\varphi\sigma^{-1}=R^{\mu\nu}_\varphi$, so $\mathcal K$ contains the six one-parameter groups $\{R^{\mu\nu}_\varphi\}$.

(iii) The six matrices $L_{\mu\nu}$ form a basis of the antisymmetric $4\times4$ matrices, which is the Lie algebra $\mathfrak{so}(4)$, of dimension $6$. Let $\Phi(\varphi_1,\dots,\varphi_6)$ be the product of the six $\exp(\varphi_pL_p)$ in a fixed order. Then $\Phi:\R^6\to SO(4)$ is smooth with $\Phi(0)=1$, and its differential at $0$ is $(\varphi_p)\mapsto\sum_p\varphi_pL_p$, an isomorphism onto $\mathfrak{so}(4)$. By the inverse function theorem $\Phi$ maps a neighbourhood of $0$ onto a neighbourhood $U$ of $1$ in $SO(4)$, and $U\subset\mathcal K$ by (ii).

(iv) Connectedness. The subgroup generated by $U$ is a union of open sets, hence open; an open subgroup is also closed, since its complement is a union of cosets. $SO(4)$ is connected, so this subgroup is $SO(4)$, and $SO(4)\subset\mathcal K$.

(v) Reflection. $W_4$ contains $\operatorname{diag}(-1,1,1,1)$, of determinant $-1$. So $\mathcal K\supset SO(4)\cup\operatorname{diag}(-1,1,1,1)SO(4)=O(4)$.
\end{proof}

\emph{Step 5: symmetry (E3) and reality.} The lattice fields are commuting real random variables. So $S^\beta_\iota(x)$ is real and invariant under simultaneous permutation of points and indices. Both properties pass to pointwise limits.

\emph{Step 6: reflection positivity (E2).} Let $\mathcal T(\R^{4k})$ be the set of $f\in{}^0\mathcal S(\R^{4k})$ whose (closed) support lies in $\{x: t(x_j)>0\text{ for all }j\}$. Then ${}^0\mathcal S_+\subset\mathcal T$. We prove positivity on $\mathcal T$, which we need in Step~7. Let $f=(f_\iota)$ be a finite family with $f_\iota\in\mathcal T$ (and a constant for the empty tuple), and set
\[
  F_\beta:=\sum_\iota a^{4k}\sum_{x\in(a\Z^4)^k}f_\iota(x)\prod_{j}Z_{i_j}(\beta)^{-1}\hat O_{i_j}(x_j/a).
\]
A lattice point with some $t(x_j)=0$ lies outside the closed support of $f_\iota$, so only sites with time coordinate at least $1$ contribute. Each such $\hat O_i(n)$ depends on edges in $\{t\ge0\}$. The series converges absolutely in the supremum norm, because the $\hat O_i$ are bounded and $f_\iota$ is a Schwartz function. So $F_\beta$ is bounded and depends only on edges in $\{t\ge0\}$. Since $\Theta\hat O_i(n)=\hat O_i(\theta n)$ and the $\hat O_i$ are real,
\[
  0\le\langle\overline{\Theta F_\beta}\,F_\beta\rangle_\beta=\sum_{\iota,\iota'}S^\beta_{(\iota',\iota)}\bigl((\Theta f_{\iota'})^*\otimes f_\iota\bigr)
\]
by H0. Each $(\Theta f_{\iota'})^*\otimes f_\iota$ lies in ${}^0\mathcal S$: inside each factor this holds by assumption, and at a point where a point $\theta$-reflected from the first group coincides with a point of the second group, the common time is either $\le0$ or $>0$, so one of the two factors vanishes on a neighbourhood. By Step~2 each term converges along $\beta_j$, and a finite sum of limits of non-negative numbers is non-negative. This gives E2, and positivity of the same form on families in $\mathcal T$.

\emph{Step 7: the gap inequality and cluster (E4).} For finite families $f,g$ in $\mathcal T$, with no constant component, put $S(f):=\sum_\iota S_\iota(f_\iota)$ and, for $s\ge0$,
\[
  B_{f,g}(s):=\sum_{\iota,\iota'}S_{(\iota',\iota)}\bigl((\Theta f_{\iota'})^*\otimes(g_\iota)_s\bigr)-\overline{S(f)}\,S(g),\qquad F_f(s):=B_{f,f}(s).
\]
By Step~6, $F_f(0)\ge0$ (apply E2 to the family consisting of $f$ and the constant $-S(f)$).

\begin{lemma}[gap inequality]\label{red:gap}
Assume H0, H1 with the unit $a=m$, H2 and H3. For finite families $f,g$ in $\mathcal T$ and $s\ge0$,
\[
  |B_{f,g}(s)|\le e^{-s}F_f(0)^{1/2}F_g(0)^{1/2};\qquad\text{in particular }|F_f(s)|\le e^{-s}F_f(0).
\]
\end{lemma}

\begin{proof}
On the lattice let $F^f_\beta,F^g_\beta$ be the observables of Step~6 and $v_f:=[F^f_\beta]-\langle F^f_\beta\rangle\Omega_\beta\in\Omega_\beta^\perp$, and similarly $v_g$. For $n\in\N$, $\tau_nF^g_\beta$ is the observable of the family $(g_\iota)_{na}$. Since $T_\beta$ is self-adjoint and $T_\beta\Omega_\beta=\Omega_\beta$, it maps $\Omega_\beta^\perp$ into itself, and
\[
  \langle v_f,T_\beta^nv_g\rangle=\langle\overline{\Theta F^f_\beta}\,\tau_nF^g_\beta\rangle_\beta-\overline{\langle F^f_\beta\rangle}\langle F^g_\beta\rangle=:B^\beta_{f,g}(na),
\]
the lattice-smeared analogue of $B_{f,g}$. Hence
\begin{equation}\label{red:latticegap}
  |B^\beta_{f,g}(na)|\le\|T_\beta|_{\Omega_\beta^\perp}\|^n\|v_f\|\|v_g\|=e^{-mn}\bigl(B^\beta_{f,f}(0)B^\beta_{g,g}(0)\bigr)^{1/2},
\end{equation}
and $mn=na$ because $a=m$. This uses only that $T_\beta$ is self-adjoint. Now let $n_j:=\lfloor s/a_j\rfloor$ and $s_j:=a_jn_j\to s$. Translation is continuous in each Schwartz norm, so $|(\Theta f_{\iota'})^*\otimes((g_\iota)_{s_j}-(g_\iota)_s)|_{N}\to0$, and the differences lie in ${}^0\mathcal S$. By the uniform bound \eqref{red:uniform} and Step~2, $B^{\beta_j}_{f,g}(s_j)\to B_{f,g}(s)$, and likewise at $s=0$. Pass to the limit in \eqref{red:latticegap}.
\end{proof}

The repaired step is this lemma. The source stated instead the monotone form $F(t+s)\le e^{-s}F(t)$ for all $t,s\ge0$. That form needs $T_\beta\ge0$, because for odd $n$ the number $\langle v,T_\beta^nv\rangle$ can be negative when $T_\beta$ has negative spectrum; H0 gives only $T_\beta$ self-adjoint with $T_\beta^2\ge0$. For the Wilson action $T_\beta\ge0$ is a theorem of \cite{Lu77}, which this paper does not use.

\emph{Cluster.} Let $f\in{}^0\mathcal S_+(\R^{4k})$, $g\in{}^0\mathcal S_+(\R^{4l})$ and let $\vec v\ne0$ be spatial. Translations by spatial vectors map ${}^0\mathcal S_+$ to itself and commute with $\Theta$. For $u,w$ in ${}^0\mathcal S_+$ write $\langle u,w\rangle:=S((\Theta u)^*\otimes w)$. By E2 this form is positive semidefinite, so it satisfies the Cauchy--Schwarz inequality, and $\langle w_{\lambda\vec v},w_{\lambda\vec v}\rangle=\langle w,w\rangle$ by E1.

(a) \emph{Reduction to compact support.} Let $\chi\in C_c^\infty(\R^{4k})$ equal $1$ near $0$ and put $f^R:=f\chi(\cdot/R)$. Then $f^R\to f$ in every Schwartz norm, $f^R$ has compact support contained in that of $f$, and $f^R\in{}^0\mathcal S_+$. Its support is a compact subset of the open set $\{0<t_1<\dots<t_k\}$, so it has positive distance from $\{t_1=0\}$ and from every coincidence set. By Cauchy--Schwarz, $|\langle f-f^R,g_{\lambda\vec v}\rangle|\le\langle f-f^R,f-f^R\rangle^{1/2}\langle g,g\rangle^{1/2}$, which tends to $0$ as $R\to\infty$ uniformly in $\lambda$, by E0$'$. The same holds for $g$, and $S(f^R)\to S(f)$ by E0$'$. So it suffices to treat $f,g$ with compact support of this kind.

(b) \emph{Rotation.} Let $\hat v:=\vec v/|\vec v|$ and let $\varrho\in SO(4)$ be the rotation by $\pi/2$ in the plane of $e_0$ and $\hat v$ that maps $\hat v$ to $e_0$. Put $u:=(\Theta f)^*\otimes g_{\lambda\vec v}$. By E1, $S(u)=S(u\circ\varrho^{-1})$, and $u\circ\varrho^{-1}=h\otimes(g')_{\lambda|\vec v|}$ with $h:=(\Theta f)^*\circ\varrho^{-1}$ and $g':=g\circ\varrho^{-1}$, both compactly supported, in a ball of radius $c_0$ say, and with supports away from all coincidence sets. Translating the whole configuration by $-c_0e_0$ (E1) gives $S(u)=S(h_{-c_0}\otimes(g'_{c_0})_{s})$ with $s:=\lambda|\vec v|-2c_0$. For $\lambda|\vec v|\ge2c_0$, $h_{-c_0}$ is supported in $\{$all $t<0\}$ and $G:=g'_{c_0}$ in $\{$all $t>0\}$. Put $\Phi:=(\Theta h_{-c_0})^*$. Then $\Phi,G\in\mathcal T$ and $h_{-c_0}=(\Theta\Phi)^*$.

(c) By Lemma~\ref{red:gap}, $S(u)=\overline{S(\Phi)}S(G)+B_{\Phi,G}(s)$ with $|B_{\Phi,G}(s)|\le e^{-s}F_\Phi(0)^{1/2}F_G(0)^{1/2}\to0$ as $\lambda\to\infty$. By E1 and reality, $\overline{S(\Phi)}=S((\Theta\Phi)^*)=S(h)=S((\Theta f)^*)$ and $S(G)=S(g)$. This is E4, with exponential rate at least $1$ for compactly supported $f,g$.

\emph{Step 8: not empty.} $S^\beta_{ii}(0,aN_\beta e_0)=1$ by the definition of $Z_i$, and $|aN_\beta-1|\le2a$ because $\xi-2<N_\beta\le\xi$ and $a\xi=1$. By \eqref{red:asym-eq} on $K_2$, $S_{ii}(0,e_0)=\lim_j\tilde S^{\beta_j}_{ii}(0,a_jN_{\beta_j}e_0)=1$. This is part~3. For later use, let $f\in C_c^\infty(\R^4)$ be real, non-negative, supported within distance $r_0<\tfrac14$ of $\tfrac12e_0$, with $\int f=1$. Then $f\in{}^0\mathcal S_+(\R^4)$, and as $r_0\to0$, $\langle f,f\rangle=\int\!\!\int S_{ii}(\theta x,y)f(x)f(y)\,dx\,dy\to S_{ii}(-\tfrac12e_0,\tfrac12e_0)=1$ by continuity and Step~3. The one-point functions vanish because the fields are centred, so $S_i(f)=0$.

\emph{Step 9: reconstruction and the gap.} By Steps~2--6, the $S_\iota$ satisfy E0$'$ and E1--E3, and Theorem~\ref{red:OS} gives a Wightman theory with fields $\varphi_1,\dots,\varphi_r$. In the description of the reconstructed space recalled after Theorem~\ref{red:OS}, a vector $\Psi=[f]$ with $f$ a finite family in ${}^0\mathcal S_+$ satisfies $\langle\Psi,e^{-sH}\Psi\rangle-|\langle\Omega,\Psi\rangle|^2=F_f(s)$. Let $\nu_\Psi$ be the spectral measure of $H$ for the component of $\Psi$ orthogonal to $\Omega$, so $F_f(s)=\int e^{-s\varepsilon}\,d\nu_\Psi(\varepsilon)$. If $\nu_\Psi([0,1-\eta])=c_1>0$ for some $\eta>0$, then $F_f(s)\ge c_1e^{-(1-\eta)s}$, which contradicts $F_f(s)\le e^{-s}F_f(0)$ (Lemma~\ref{red:gap}) for large $s$. So the spectral projection $E_H([0,1))$ annihilates the orthogonal component of a dense set of vectors, hence of every vector: $\operatorname{spec}(H|_{\Omega^\perp})\subset[1,\infty)$. By Step~8, the vector $[f]$ there is non-zero and orthogonal to $\Omega$, so $\Omega^\perp\ne\{0\}$.

\emph{Unique vacuum.} A translation-invariant vector is invariant under time translations, so $H\Psi=0$; by the gap, $\Psi$ is a multiple of $\Omega$.

\emph{Mass hyperboloid.} The Wightman theory carries a unitary representation of the Poincar\'e group, so the joint spectrum of $(H,\vec P)$ on $\Omega^\perp$ is a closed Lorentz-invariant subset of the closed forward cone. If it contained a point $p$ with $0<p^2<1$, a boost would carry $p$ to a point with energy $\sqrt{p^2}<1$. If it contained $p\ne0$ with $p^2=0$, boosts would give points of arbitrarily small energy. If it contained $p=0$, the energy $0$ would be in $\operatorname{spec}(H|_{\Omega^\perp})$. Each contradicts the gap. So the joint spectrum on $\Omega^\perp$ lies in $\{p^2\ge1,\ p^0>0\}$. This proves part~4 and completes the proof of Theorem~\ref{thm:conditional}. \qed

\emph{Where each hypothesis enters.} Step~1: H3(b), H3(c), H2. Step~2: H3(b), Step~1. Step~3: H0, H2, H3(c). Step~4: H0, H4, H3(c), Lemma~\ref{red:dense}. Step~5: reality and commutativity of the fields. Step~6: H0, site-scalarity, boundedness of the fields, Step~2. Step~7: H0 (self-adjointness of $T_\beta$), H1 with $a=m$, Steps~2 and~6, E1. Step~8: H3(a), H3(c), the normalization. Step~9: Theorem~\ref{red:OS} and the recalled description of the reconstructed space.

\subsection{What Theorem~\ref{thm:conditional} does not give}\label{red:not}

\begin{itemize}
\item \emph{Uniqueness.} The limit is taken along a subsequence. Different subsequences, or different choices of $\mu_\beta\in\cL_\beta$, need not give the same limit.
\item \emph{Interaction.} Nothing in H0--H4 excludes a free theory; Section~\ref{red:free} shows that the free lattice field satisfies the analogues of all of them. A test on composite fields does not detect interaction in the form ``the limit is not a generalized free field'': for the free field the Wick square $:\!\varphi^2\!:$ has non-zero truncated four-point function. A correct form would be a statement about the reconstructed theory, such as a non-trivial scattering matrix. It is not formulated here.
\item \emph{Short-distance behaviour.} The theorem says nothing about agreement with asymptotic freedom at short distances.
\item \emph{Other groups.} The proof uses no property of $\SU(2)$. For another compact simple group it requires H0--H4 for that group, and nothing is claimed about them here.
\end{itemize}
The official problem statement \cite{JW} remarks, in substance, that existence by compactness alone is not accepted unless other techniques establish the properties of the limit. Theorem~\ref{thm:conditional} is an argument of exactly that shape: Arzel\`a--Ascoli, limits of quadratic forms, and the reconstruction theorem. It assumes, as H1, H3 and H4, exactly the properties that such a remark asks to be established.

\subsection{A test of the chain on a free field}\label{red:free}

The lattice free massive scalar field on $\Z^4$, with mass parameter $m\to0$, gives an analogue of the setting of Theorem~\ref{thm:conditional}. The unit is the exact gap $a=E_0(m)=2\operatorname{asinh}(m/2)$, and the fields are $\varphi$ and $:\!\varphi^2\!:$. It is an analogue and not an instance: the fields are unbounded, there is no gauge group, and boundedness is used in Step~6. Moment bounds replace boundedness. We do not prove a version of Theorem~\ref{thm:conditional} for this analogue. The test checks that no hypothesis of the chain fails on a model whose answer is known, and that the chain then returns the free field of mass $1$, with two-point function $S_2(x)=K_1(|x|)/(|x|K_1(1))$.

Two separately written programs computed the numbers in Table~\ref{red:freetab}. Program~II uses a Bessel--Schwinger representation of the lattice Green function. Mass parameters are $m=\tfrac14,\tfrac18,\tfrac1{16},\tfrac1{32}$ unless a column says otherwise. These are numerical results, not proofs.

\begin{table}[ht]
\centering
\small
\begin{tabular}{@{}p{0.17\linewidth}p{0.47\linewidth}p{0.29\linewidth}@{}}
\toprule
Item & Program I & Program II \\
\midrule
H0 & RP matrix on 40 random points (seed 12345): least eigenvalue $+3.1\cdot10^{-10}$ & -- \\
H1 & zero-momentum correlator equals $e^{-E_0n}/(2\sinh E_0)$ to $10^{-13}$ ($n\ge1$), 4 of 4 masses & gap $2\operatorname{asinh}(m/2)$ to $10^{-15}$, 12 of 12 \\
H3(b), $\varphi$ & $\sup|x|^2S_2=1.64,\ 2.00,\ 2.11,\ 2.15$ & $1.636,\ 2.001,\ 2.114,\ 2.149$ \\
H3(b), $:\!\varphi^2\!:$ & -- & $\sup|x|^4S_{\varphi^2}\le4.62$ \\
H3(c) & Lipschitz constant on the axis, $|x|\ge\tfrac12$: $5.1,\ 10.0,\ 15.0,\ 19.1$ (continuum $25.1$) & -- \\
H4 & $|S(0,\cos1,\sin1,0)-S(0,1,0,0)|=3.9\cdot10^{-2},\ 2.7\cdot10^{-2},\ 6.7\cdot10^{-3}$; $|S(1,0,0,0)-S(\tfrac12,\tfrac12,\tfrac12,\tfrac12)|=3.7\cdot10^{-2},\ 9.0\cdot10^{-3},\ 2.2\cdot10^{-3}$ ($m=\tfrac18,\tfrac1{16},\tfrac1{32}$) & rotation differences agree with program I \\
Limit & relative error against $K_1(|x|)/(|x|K_1(1))$ at $m=\tfrac1{32}$: $\le2.4\cdot10^{-3}$ at 8 points, on and off axis & $\le2.3\cdot10^{-3}$ at 6 points; $a^{-2}G$ against $K_1(|x|)/(4\pi^2|x|)$: $\le1.4\cdot10^{-3}$ \\
$:\!\varphi^2\!:$, truncated 4-point & -- & Monte Carlo $1.54\cdot10^{-5}\pm2.1\cdot10^{-6}$; Wick $1.60\cdot10^{-5}$ \\
\bottomrule
\end{tabular}
\caption{Free-field analogue of the chain: numerical values.}\label{red:freetab}
\end{table}

\emph{Status of each line.} Wick's theorem for Gaussian measures is a theorem; with it, the truncated four-point function of $:\!\varphi^2\!:$ is $16$ times a sum of three cycle products of the positive two-point function, hence positive. The dispersion relation behind H1 and the gap $2\operatorname{asinh}(m/2)$ is a standard computation; here it is checked numerically. Reflection positivity of the nearest-neighbour Gaussian measure in site planes is standard and is not proved here; the eigenvalue in the table illustrates it. H3(b) for $\varphi$ follows from Wick's theorem, the inequality $|x_i-x_j|^{-2}\le\rho_i^{-1}\rho_j^{-1}$ and the bound $G_m(n)\le C/|n|^2$ on the lattice Green function, uniform in $m$; that bound is recalled from the literature on Green-function asymptotics, and we do not prove it here. Isotropy of the limit (H4) is supported by numerics only; its proof by a local central limit theorem is recalled and not written out. The convergence to $K_1(|x|)/(|x|K_1(1))$ is numerical. With these qualifications the free field satisfies the analogues of H0--H4, and the chain returns the free field of mass $1$. Interaction, in the sense of Section~\ref{red:not}, fails for $\varphi$ (its truncated four-point function vanishes).

\subsection{Status of the hypotheses for \texorpdfstring{$\SU(2)$}{SU(2)}}\label{red:statussec}

\begin{table}[ht]
\centering
\small
\begin{tabular}{@{}lll@{}}
\toprule
Hypothesis & Content & Status for $\SU(2)$ \\
\midrule
H0 (Hypothesis~\ref{hyp:H0}) & lattice axioms & proved \\
H1 = (A) (Hypothesis~\ref{hyp:A}) & gap at every coupling & open: it is the problem \\
H2 = (B) (Hypothesis~\ref{hyp:B}) & $m(\beta)\to0$ & proved (Theorem~\ref{thm:main}) \\
H3 (Hypothesis~\ref{hyp:H3}) & uniform regularity & open \\
H4 (Hypothesis~\ref{hyp:H4}) & one irrational rotation & open \\
\bottomrule
\end{tabular}
\caption{Status of the hypotheses of Theorem~\ref{thm:conditional} for $\SU(2)$.}\label{red:status}
\end{table}

For H0, translation and $W_4$ invariance of the periodic states pass to weak limits, and site-plane reflection positivity of the periodic Wilson states \cite{OS78} passes to weak limits. Hypothesis~(C) of Theorem~\ref{thm:reduction} is open for $\SU(2)$ as well.

% sec_proof.tex -- Section: proof of Theorem 1.1 (Theorem B').
% To be \input into a document using preamble.tex. All labels carry the prefix pf:.
\section{Proof of Theorem~\ref{thm:main}}\label{sec:proof}

We prove Theorem~\ref{thm:main} with three restrictions, which are part of its statement.
The periodic boxes have side $L=2^k$.
The states are weak limit points of periodic finite-volume states, not arbitrary Gibbs states.
The conclusion ``$m_\mu(\beta)\to0$'' includes the case $m_\mu(\beta)=0$.
We recall the statement:
\begin{equation}\label{pf:eq-main}
\begin{gathered}
\liminf_{\beta\to\infty}\ \inf_{\mu\in\cL_\beta}\ \frac{\Cov_\mu(O_{p_0},O_{p_n})}{\Var_\mu(O_{p_0})}\ \ge\ \Big(\frac{G(n)}{G(0)}\Big)^2\quad\text{for every } n\ge1,\\
\text{and consequently }\ \sup_{\mu\in\cL_\beta}m_\mu(\beta)\to0\ \text{ as }\beta\to\infty,
\end{gathered}
\end{equation}
where $G$ is the free lattice Maxwell field-strength correlation of Definition~\ref{pf:def-G}, $G(0)=\tfrac12$, $G(n)>0$, and $G(n)=1/(\pi^2n^4)+O(n^{-6})$.
The order of limits is: the box first, then the coupling, then the separation.

\subsection{Setting and notation}\label{pf:sec-setting}

\emph{Lattice.} Directions are $0,1,2,3$; direction $0$ is time. $e_\mu$ is the unit vector. $\Lambda_L=(\Z/L\Z)^4$ with $L=2^k$, $|\Lambda_L|=L^4$. A link is an ordered pair $(x,x+e_\mu)$; we write $U_{x,\mu}\in\SU(2)$ for its variable and $U_{x+e_\mu,x}:=U_{x,\mu}^{-1}$ for the reversed link. A plaquette is $p=(x;\mu,\nu)$ with $\mu<\nu$; its corner is $x_p:=x$, and its holonomy, based at $x_p$, is
\[
U_p:=U_{x,\mu}\,U_{x+e_\mu,\nu}\,U_{x+e_\nu,\mu}^{-1}\,U_{x,\nu}^{-1}.
\]
For a path $\gamma=(y_0,\dots,y_m)$ of nearest neighbours, $U_\gamma:=U_{y_0y_1}\cdots U_{y_{m-1}y_m}$.
A gauge transformation $g=(g_x)$ acts by $(U^g)_{x,\mu}=g_xU_{x,\mu}g_{x+e_\mu}^{-1}$; then $U_\gamma\mapsto g_{y_0}U_\gamma g_{y_m}^{-1}$ and $U_p\mapsto g_{x_p}U_pg_{x_p}^{-1}$.
Each link of $\Z^4$ lies in exactly $6$ plaquettes: for each of the three other directions, one on each side.

\emph{Class angle.} For $U\in\SU(2)$, $\alpha(U)\in[0,\pi]$ is defined by $\tfrac12\Tr U=\cos\alpha(U)$.
Identify $\SU(2)$ with the unit quaternions, that is, with the unit sphere $S^3\subset\R^4$, so that $\tfrac12\Tr U$ is the Euclidean inner product of $U$ with $1$. Then $\alpha(U)$ is the geodesic distance from $U$ to $1$ on $S^3$, and left and right multiplications are isometries. Hence $(U,V)\mapsto\alpha(UV^{-1})$ is a bi-invariant metric, and
\begin{equation}\label{pf:eq-alpha}
\begin{gathered}
\alpha(U^{-1})=\alpha(U),\qquad \alpha(VUV^{-1})=\alpha(U),\\
\alpha(UV)\le\alpha(U)+\alpha(V),\qquad |\alpha(AVB)-\alpha(AB)|\le\alpha(AVA^{-1})=\alpha(V),
\end{gathered}
\end{equation}
the last by the triangle inequality for the metric.
We write $\alpha_p:=\alpha(U_p)$; it does not depend on the base point or orientation of the loop.

\emph{Action and states.} $S:=\sum_p(1-\cos\alpha_p)$ and
\[
\mu_{\beta,L}(dU):=Z_L(\beta)^{-1}e^{-\beta S(U)}\prod_{\text{links}}dU_{x,\mu},\qquad Z_L(b):=\int e^{-bS}\,dU,
\]
with normalised Haar measure $dU$. By Weyl's formula, the law of $\alpha(U)$ under Haar measure has density $\tfrac2\pi\sin^2\alpha$ on $[0,\pi]$.
We embed $\Lambda_L$ in $\Z^4$ around the origin and regard $\mu_{\beta,L}$ as a measure on $\SU(2)^{E(\Z^4)}$ for local questions.
$\cL_\beta$ is the set of weak limit points of $(\mu_{\beta,2^k})_{k\ge1}$. It is nonempty because $\SU(2)^{E(\Z^4)}$ is compact.
$O_p:=\beta(1-\cos\alpha_p)\in[0,2\beta]$. $p_0:=(0;1,2)$ and $p_n:=p_0+ne_0$.

\emph{Lie algebra.} We identify $\R^3$ with $\mathfrak{su}(2)$ through $x\mapsto\sigma\cdot x$ (Pauli matrices). For $|x|\le\pi$, $U=e^{i\sigma\cdot x}=\cos|x|+i\sin|x|\,\hat x\cdot\sigma$ and $\alpha(U)=|x|$. For $\alpha(U)<\pi$ we write $\log U$ for the unique $x$ with $|x|<\pi$ and $U=e^{i\sigma\cdot x}$; for $\alpha(U)=\pi$ we set $\log U:=(\pi,0,0)$. In all cases $|\log U|=\alpha(U)$. $\mathrm{Ad}(U)\in\mathrm{SO}(3)$ is defined by $\sigma\cdot(\mathrm{Ad}(U)x)=U(\sigma\cdot x)U^{-1}$; it is the rotation by the angle $2\alpha(U)$, so $\|\mathrm{Ad}(U)-I\|\le2|\sin\alpha(U)|\le2\alpha(U)$.
Norms of vectors in $(\R^3)^E$ are Euclidean over all links and colours.

\emph{Forms.} Real $r$-forms on $\Z^4$ are functions on oriented $r$-cells. For a $1$-form $A$, $(dA)_{(x;\mu,\nu)}=A_{x+e_\mu,\nu}-A_{x,\nu}-A_{x+e_\nu,\mu}+A_{x,\mu}$, and $d^*$ is the $\ell^2$-adjoint of $d$. On a $3$-cube $c=(x;\lambda,\mu,\nu)$, $\lambda<\mu<\nu$,
$(d\omega)_c=\omega_{\mu\nu}(x+e_\lambda)-\omega_{\mu\nu}(x)-\omega_{\lambda\nu}(x+e_\mu)+\omega_{\lambda\nu}(x)+\omega_{\lambda\mu}(x+e_\nu)-\omega_{\lambda\mu}(x)$.
For $\R^3$-valued forms, $d$ acts on each colour.

\emph{Boxes.} $R=z+[-N,N]^4$ with $N\ge1$. The shell $\partial R$ is the set of sites with $|x-z|_\infty=N$; the other sites of $R$ are interior. $\Lambda'_R$ is the set of links with at least one interior endpoint. $P_R$ is the set of plaquettes containing a link of $\Lambda'_R$. Shell links and shell plaquettes are those with all sites in $\partial R$. $T_{\rm int}$ is the set of links $(x,x+e_0)$ with $x$ interior. $T_s$ is a spanning tree of the shell graph, rooted at a shell site $b$. $\Lambda''_R:=\Lambda'_R\setminus T_{\rm int}$, $\nu_R:=|\Lambda''_R|$, and $\nu_s$ is the number of shell links not in $T_s$. We write $R_N:=[-N,N]^4$.

\subsection{Outline}\label{pf:sec-outline}

The proof has four parts.

\emph{Part I: reduction to one correlation} (\S\ref{pf:sec-rp}--\S\ref{pf:sec-transfer}). Reflection positivity in site planes holds on every periodic box of even side and passes to every $\mu\in\cL_\beta$. It gives a Hilbert space and a self-adjoint contraction $T_\mu$, the transfer matrix, with $m_\mu=-\log\|T_\mu|_{\Omega^\perp}\|$. For the vector $v$ of $O_{p_0}-\mu(O_{p_0})$ we have $\Cov_\mu(O_{p_0},O_{p_n})=\langle v,T_\mu^nv\rangle\le e^{-nm_\mu}\Var_\mu(O_{p_0})$ for every $n\ge1$. So a lower bound $c(n)$ on the ratio in \eqref{pf:eq-main}, with $-\log c(n)/n\to0$, forces $\sup_\mu m_\mu\to0$.

\emph{Part II: uniform tails} (\S\ref{pf:sec-chess}--\S\ref{pf:sec-free}). A chessboard estimate, which we prove for sides $2^k$ by repeated reflection, bounds $\mu_{\beta,L}(O_p>t)$ by $\exp(f_L(\beta/2)-f_L(\beta)-t/8)$ with $f_L$ the free energy per site. Elementary upper and lower bounds on $f_L$ give $\mu(O_p>t)\le e^{60.2-t/8}$ for all $\mu\in\cL_\beta$, $\beta\ge1$, $t\ge0$. So the rescaled curvature $\sqrt\beta\,\alpha_p$ is tight uniformly in $\beta$ and $\mu$, with all moments bounded.

\emph{Part III: the Gaussian limit} (\S\ref{pf:sec-box}--\S\ref{pf:sec-L4}). We condition on the links outside a box $R$ and fix the gauge by a tree in the shell and a forest of interior time links; this gauge uses no interior information, so the conditioning is not circular. The shell of a $4$-box is simply connected, so bounded shell curvatures give bounded shell data. The only zero of the action with trivial shell data is the trivial configuration, and the Dirichlet Maxwell form is nondegenerate. A Laplace argument (Theorem~\ref{pf:thm-L3}) shows that the rescaled conditional law converges, uniformly over bounded shell data, to the Gaussian with the Dirichlet Maxwell form. A lasso field $F^\beta_p=\sqrt\beta\,\mathrm{Ad}(U_{\gamma_p})\log U_p$, defined without a gauge tree, then has subsequential limits along every sequence $\mu_k\in\cL_{\beta_k}$, $\beta_k\to\infty$ (Proposition~\ref{pf:prop-L4}). Every such limit $F$ is closed, colour and translation invariant, has all moments, is the limit of the observables $O_p$ in law and in moments, and satisfies the Gaussian DLR equations of the free Maxwell field in every box.

\emph{Part IV: the free comparison} (\S\ref{pf:sec-D}--\S\ref{pf:sec-assembly}). The Dirichlet covariances increase to the projection $P$ onto exact $2$-forms. By the downward martingale theorem every limit field is $F=\xi+f$, with $\xi$ the free field of covariance $P$ in three independent colours and $f$ an independent background. The background is closed, co-closed and second-order stationary; a two-line computation in $L^2$ shows that it is a random constant. Wick's theorem then gives
$\Cov(X_0,X_n)=\tfrac32G(n)^2+G(n)\E|f|^2+\tfrac14\Var|f|^2$ with $X_m=\tfrac12|F_{p_m}|^2$, and each term has ratio at least $(G(n)/G(0))^2$ to its counterpart at $n=0$. Finally $G(n)>0$ and $G(n)=1/(\pi^2n^4)+O(n^{-6})$, so $c(n)=4G(n)^2=4/(\pi^4n^8)\,(1+O(n^{-2}))$ and $-\log c(n)/n\to0$.

Two constants are obtained by compactness and are not explicit: $\eta(\delta)$ in Theorem~\ref{pf:thm-L3} (hence the constant $c$ there) and the combinatorial constant $K_N$ of Lemma~\ref{pf:lem-shell}. Neither enters the final bound. The measure-theoretic facts used without proof are listed in \S\ref{pf:sec-standard}.

\subsection{Reflection positivity and limit states}\label{pf:sec-rp}

Reflection positivity of the Wilson action is due to Osterwalder and Seiler~\cite{OS78}, and the transfer matrix of lattice gauge theory to L\"uscher~\cite{Lu77}; we re-prove the case we use. Fix $j\in\{0,1,2,3\}$ and even $L$. Let $\vartheta$ be the reflection $x_j\mapsto-x_j$ of $\Lambda_L$. It fixes the two site planes $\pi=\{x_j=0\}\cup\{x_j=L/2\}$. On configurations it acts by $(\vartheta U)_{b}=U_{\vartheta b}$, where a link in direction $j$ is mapped to a reversed link, whose variable is the inverse. For a function $F$ put $\theta F:=\overline{F\circ\vartheta}$.
Let $E_0$ be the set of links in $\pi$; $E_+$ (resp. $E_-$) the set of links not in $E_0$ with both endpoints in $\{0\le x_j\le L/2\}$ (resp. $\{L/2\le x_j\le L\}$). These three sets partition the links, and $\vartheta$ fixes $E_0$ pointwise and exchanges $E_+$ and $E_-$. Let $\cH_+$ be the bounded measurable functions of the links in $E_0\cup E_+$.

\begin{lemma}[site-plane reflection positivity]\label{pf:lem-rp}
For even $L$, every $j$ and all $F,G\in\cH_+$: $\mu_{\beta,L}(\theta F\cdot F)\ge0$, $\mu_{\beta,L}(\theta F\cdot G)=\overline{\mu_{\beta,L}(\theta G\cdot F)}$, and $|\mu_{\beta,L}(\theta F\cdot G)|^2\le\mu_{\beta,L}(\theta F\cdot F)\,\mu_{\beta,L}(\theta G\cdot G)$. The same holds for the reflection $x_j\mapsto2a-x_j$ for every $a\in\Z/L\Z$, with the corresponding half-spaces.
\end{lemma}

\begin{proof}
A plaquette has extent at most one in direction $j$, so all its sites lie in $\{0\le x_j\le L/2\}$ or all in $\{L/2\le x_j\le L\}$. Let $W_+$ be the sum of $\cos\alpha_p$ over plaquettes with all sites in the first set, with weight $\tfrac12$ for plaquettes contained in $\pi$. Since $\vartheta$ maps $U_p$ to a conjugate of $U_{\vartheta p}^{\pm1}$ and $\tfrac12\Tr U^{-1}=\tfrac12\Tr U$ on $\SU(2)$, we get $\sum_p\cos\alpha_p=W_++W_+\circ\vartheta$. Haar measure is invariant under $U\mapsto U^{-1}$, so $\vartheta$ maps $\prod_{E_-}dU$ to $\prod_{E_+}dU$. For $F\in\cH_+$, with $I(U_{E_0}):=\int F e^{\beta W_+}\prod_{E_+}dU$,
\[
Z_L(\beta)\,\mu_{\beta,L}(\theta F\cdot F)=e^{-6\beta|\Lambda_L|}\int\prod_{E_0}dU\ \overline{I(U_{E_0})}\,I(U_{E_0})\ \ge0 .
\]
So $(F,G)\mapsto\mu_{\beta,L}(\theta F\cdot G)$ is a positive semidefinite Hermitian form on $\cH_+$, which gives the other two claims. Translation invariance of $\mu_{\beta,L}$ gives the version at $a$.
\end{proof}

\begin{lemma}[limit states]\label{pf:lem-limit}
Let $\beta>0$ and $\mu\in\cL_\beta$. Then:
\begin{enumerate}
\item $\mu$ is invariant under the translations, rotations and reflections of $\Z^4$, and under every gauge transformation $g$ (each local function depends on finitely many $g_x$).
\item For the reflection $\vartheta$ in $\{x_0=0\}$ and every bounded continuous $F$ depending on finitely many links in $\{x_0\ge0\}$: $\mu(\theta F\cdot F)\ge0$.
\item (DLR) For every finite set $\Lambda$ of links and every bounded measurable $g$, $\mu(g)=\mu(\gamma_\Lambda g)$, where $\gamma_\Lambda g$ integrates $U_\Lambda$ against $Z^{-1}\exp(-\beta\sum_{p\cap\Lambda\ne\emptyset}(1-\cos\alpha_p))\prod_\Lambda dU$, the other links fixed.
\end{enumerate}
\end{lemma}

\begin{proof}
Let $\mu=\lim_i\mu_{\beta,L_i}$ weakly. Each $\mu_{\beta,L}$ has the symmetries in (1) for the torus, and for a local continuous function only finitely many of them matter once $L$ is large; weak limits preserve the identities $\mu_L(h\circ s)=\mu_L(h)$ for continuous local $h$. For (2), once $L_i$ is large, $F\in\cH_+$ on the torus, and $\nu\mapsto\nu(\theta F\cdot F)$ is weakly continuous. For (3), let first $g$ be continuous and local. Then $\gamma_\Lambda g$ is continuous and local, and $\mu_{L_i}(\gamma_\Lambda g)=\mu_{L_i}(g)$ for large $L_i$; pass to the limit. The two finite measures $\mu$ and $\mu\gamma_\Lambda$ agree on continuous local functions, hence they are equal, which is the claim for bounded measurable $g$.
\end{proof}

\subsection{The transfer matrix and the spectral bound}\label{pf:sec-transfer}

Fix $\mu\in\cL_\beta$. Let $\mathcal A_+$ be the bounded continuous functions of finitely many links in $\{x_0\ge0\}$, $\langle F,G\rangle:=\mu(\theta F\cdot G)$, $\cH$ the completion of $\mathcal A_+$ modulo null vectors, $[F]$ the class of $F$, and $\Omega:=[1]$. Let $\tau$ be the unit time translation, $(\tau F)(U)=F(U_{\cdot-e_0})$, so that $\tau\mathcal A_+\subset\mathcal A_+$ and $\theta\tau=\tau^{-1}\theta$.

\begin{lemma}[transfer matrix]\label{pf:lem-transfer}
$T[F]:=[\tau F]$ defines a self-adjoint contraction $T=T_\mu$ on $\cH$ with $T\Omega=\Omega$ and $T\Omega^\perp\subset\Omega^\perp$.
\end{lemma}

\begin{proof}
By translation invariance, $\langle G,\tau F\rangle=\mu(\theta G\cdot\tau F)=\mu(\tau^{-1}\theta G\cdot F)=\mu(\theta(\tau G)\cdot F)=\langle\tau G,F\rangle$; likewise $\langle\tau^aF,\tau^bF\rangle=\langle F,\tau^{a+b}F\rangle$. Put $a_m:=\|\tau^mF\|$. Cauchy--Schwarz gives $a_{2^i}^2=\langle F,\tau^{2^{i+1}}F\rangle\le a_0a_{2^{i+1}}$, hence $a_1\le a_0^{1-2^{-j}}a_{2^j}^{2^{-j}}$ for every $j$. Since $a_m\le\|F\|_\infty$, letting $j\to\infty$ gives $a_1\le a_0$ (if $a_0=0$ then $a_1=0$). So $\tau$ maps null vectors to null vectors and defines a symmetric contraction on a dense subspace, which extends to a self-adjoint contraction. $T\Omega=\Omega$ is clear, and $\langle\Omega,TF\rangle=\mu(\tau F)=\mu(F)=\langle\Omega,F\rangle$ gives invariance of $\Omega^\perp$.
\end{proof}

We define $m_\mu(\beta):=-\log\|T_\mu|_{\Omega^\perp}\|\in[0,\infty]$.

\begin{proposition}[spectral bound and reduction]\label{pf:prop-spectral}
Let $\mu\in\cL_\beta$. Then $\Var_\mu(O_{p_0})>0$, and for every $n\ge1$
\begin{equation}\label{pf:eq-S1}
\Cov_\mu(O_{p_0},O_{p_n})\le e^{-n\,m_\mu(\beta)}\Var_\mu(O_{p_0}).
\end{equation}
Consequently, if for every $n\ge1$ there are $c(n)>0$ with
\[\liminf_{\beta\to\infty}\inf_{\mu\in\cL_\beta}\Cov_\mu(O_{p_0},O_{p_n})/\Var_\mu(O_{p_0})\ge c(n)\]
and $\inf_n(-\log c(n))/n\le0$, then $\sup_{\mu\in\cL_\beta}m_\mu(\beta)\to0$ as $\beta\to\infty$.
\end{proposition}

\begin{proof}
$O_{p_0}$ is real, bounded, continuous and depends on links in $\{x_0=0\}$, which $\vartheta$ fixes; so $\theta O_{p_0}=O_{p_0}$. Let $v:=[O_{p_0}-\mu(O_{p_0})]\in\Omega^\perp$. Then $\|v\|^2=\Var_\mu(O_{p_0})$ and $\langle v,T^nv\rangle=\mu\big((O_{p_0}-\mu O_{p_0})(O_{p_n}-\mu O_{p_0})\big)=\Cov_\mu(O_{p_0},O_{p_n})$, using $\mu(O_{p_n})=\mu(O_{p_0})$. Since $\Omega^\perp$ is invariant, $|\langle v,T^nv\rangle|\le\|T|_{\Omega^\perp}\|^n\|v\|^2$, which is \eqref{pf:eq-S1} for every $n\ge1$, odd or even.
By Lemma~\ref{pf:lem-limit}(3) with $\Lambda=\{(0,e_1)\}$, given the other links the law of $U:=U_{0,1}$ has a positive density with respect to Haar measure, and $O_{p_0}=\beta(1-\tfrac12\Tr(UV))$ with $V$ fixed by the other links. Each level set $\{\tfrac12\Tr(UV)=s\}$ is a translate of a conjugacy class and has Haar measure zero. So the conditional law of $O_{p_0}$ has no atoms, and $\Var_\mu(O_{p_0})>0$.
For the reduction, fix $n$ and $\varepsilon\in(0,c(n))$. For $\beta$ large and every $\mu\in\cL_\beta$ the ratio is $\ge c(n)-\varepsilon>0$, so \eqref{pf:eq-S1} gives $m_\mu(\beta)\le-\log(c(n)-\varepsilon)/n$. Hence $\limsup_\beta\sup_\mu m_\mu\le-\log c(n)/n$ for every $n$, and the infimum over $n$ is $\le0$. Since $m_\mu\ge0$, the claim follows.
\end{proof}

\subsection{The chessboard estimate for dyadic sides}\label{pf:sec-chess}

The chessboard estimate is due to Fr\"ohlich, Israel, Lieb and Simon~\cite{FILS78}; we re-prove the case we use. Let $L=2^k$. A block is a closed unit cube $[t,t+1]$, $t\in\Lambda_L$, with all its links and faces. For $t\in\Lambda_L$ let $\vartheta_t$ be the lattice map that is, in each direction $j$, the translation $x_j\mapsto x_j+t_j$ if $t_j$ is even and the reflection $x_j\mapsto t_j+1-x_j$ if $t_j$ is odd. It maps block $0$ onto block $t$. For a function $f$ of the links of block $0$ put $(f\circ\vartheta_t)(U):=f(\vartheta_t^*U)$, $(\vartheta_t^*U)_b:=U_{\vartheta_tb}$. Let $\theta_a$ denote the reflection $x_j\mapsto2a-x_j$ in one direction $j$.

\begin{lemma}[chessboard estimate]\label{pf:lem-chess}
Let $L=2^k$ and let $f\ge0$ be a bounded measurable function of the links of block $0$. Then
\[
\mu_{\beta,L}(f)\le\mu_{\beta,L}\Big(\prod_{t\in\Lambda_L}f\circ\vartheta_t\Big)^{1/|\Lambda_L|}.
\]
\end{lemma}

\begin{proof}
For a set $S$ of blocks put $Z(S):=\mu_{\beta,L}(\prod_{t\in S}f\circ\vartheta_t)$, so $Z(\emptyset)=1$.
In direction $j$, $\theta_a\circ\vartheta_t=\vartheta_{t'}$ with $t'_j=2a-1-t_j$ and $t'_i=t_i$ for $i\ne j$. Indeed, if $t_j$ is even, $x\mapsto2a-x-t_j$ is the reflection $x\mapsto t'_j+1-x$ with $t'_j$ odd; if $t_j$ is odd, $x\mapsto2a-(t_j+1-x)$ is the translation by the even number $t'_j$. Hence $(f\circ\vartheta_t)\circ\theta_a=f\circ\vartheta_{t'}$, and $\theta(f\circ\vartheta_t)=f\circ\vartheta_{t'}$ since $f$ is real.

\emph{Doubling step.} Let $S\subset H_+:=\{t:a\le t_j\le a+L/2-1\}$. Blocks in $H_+$ have all their sites in $\{a\le x_j\le a+L/2\}$, so $F:=\prod_{t\in S}f\circ\vartheta_t$ lies in $\cH_+$ for the reflection $\theta_a$. The set $\theta_aS$ lies in $\{a-L/2\le t_j\le a-1\}$ and is disjoint from $S$. Lemma~\ref{pf:lem-rp} with $G=1$ gives
\[
Z(S)^2=|\mu(\theta1\cdot F)|^2\le\mu(\theta F\cdot F)\,\mu(1)=Z(S\cup\theta_aS).
\]

\emph{Iteration.} In direction $1$ take $a=-2^i+1$ for $i=0,\dots,k-1$, starting from $S=\{0\}$. If $S=\{t:t_1\in[-2^i+1,0],\,t_2=t_3=t_0=0\}$, then $S\subset H_+$ because $0\le a+L/2-1$ is $2^i\le L/2$, and $\theta_aS$ has $t_1\in[-2^{i+1}+1,-2^i]$. So $S\cup\theta_aS$ has $t_1\in[-2^{i+1}+1,0]$. After $k$ steps $t_1$ runs over $\Z/L\Z$. Repeat in directions $2,3,0$ with slabs. After $4k$ doublings, $Z(\{0\})^{2^{4k}}\le Z(\Lambda_L)$, and $2^{4k}=|\Lambda_L|$.
\end{proof}

\subsection{Free-energy bounds and the uniform tail}\label{pf:sec-free}

Let $f_L(b):=|\Lambda_L|^{-1}\log Z_L(b)$ and $\hat z(b):=\int e^{-b(1-\cos\alpha(U))}dU$.

\begin{lemma}[forest gauge]\label{pf:lem-forest}
Let $\mathcal T$ be a forest of links of a finite graph, $B$ a set of sites containing exactly one site of each tree of $\mathcal T$, and $F$ an integrable function of the links that is invariant under gauge transformations equal to $1$ on $B$. Then $\int F\,dU=\int F|_{U_{\mathcal T}=1}\,dU_{\mathcal T^c}$.
\end{lemma}

\begin{proof}
Define $g$ by $g=1$ on $B$ and, along each tree from its site in $B$, $g_y:=U_{yx}^{-1}g_x$ when $(x,y)$ is a tree link with $x$ already treated; then $U^g=1$ on $\mathcal T$, and $g$ depends only on $U_{\mathcal T}$. Thus $F(U)=F(U^g)$. At fixed $U_{\mathcal T}$, the map $U_{\mathcal T^c}\mapsto(U^g)_{\mathcal T^c}$ is a product of left and right multiplications, which preserve Haar measure. Integrate $U_{\mathcal T^c}$ first, then $U_{\mathcal T}$.
\end{proof}

\begin{lemma}[free-energy bounds]\label{pf:lem-free}
For every $b>0$ and every $\beta\ge1$,
\begin{align*}
f_L(b)&\le3(1-1/L)\log\hat z(b),\qquad \hat z(b)\le(\pi/2)^{5/2}b^{-3/2},\\
f_L(\beta)&\ge-48+(3+1/L)\log\big(c_H\beta^{-3/2}\big),\qquad c_H:=2\sin^21/(3\pi).
\end{align*}
Consequently $\limsup_{L=2^k\to\infty}\big(f_L(\beta/2)-f_L(\beta)\big)\le D^*:=48+3\log\big((\pi/2)^{5/2}2^{3/2}\big)-3\log c_H=60.1922\ldots$ for every $\beta\ge1$.
\end{lemma}

\begin{proof}
Let $\mathcal T$ be the set of time links $(x,x+e_0)$ with $x_0\in\{0,\dots,L-2\}$; it is a forest of $L^3$ paths, each containing $L$ sites. Apply Lemma~\ref{pf:lem-forest} to $e^{-bS}$ with $B=\{x_0=0\}$.

\emph{Upper bound.} Drop every plaquette factor except those of the $(0,i)$ plaquettes with $x_0\le L-2$; each dropped factor is $\le1$. With $U_{\mathcal T}=1$ such a plaquette has holonomy $V_{x_0}V_{x_0+1}^{-1}$, up to conjugation and inversion, where $V_s$ is the variable of the spatial link $\ell$ at time $s$. So for each of the $3L^3$ spatial links $\ell$ we integrate $\prod_{s=0}^{L-2}w(V_sV_{s+1}^{-1})$ with $w(U)=e^{-b(1-\cos\alpha(U))}$. Integrating $V_{L-1},V_{L-2},\dots,V_1$ in turn gives $\hat z(b)^{L-1}$ by Haar invariance; the remaining variables integrate to $1$. Hence $Z_L(b)\le\hat z(b)^{3L^3(L-1)}$. Next, $\sin\alpha\le\alpha$ and $1-\cos\alpha\ge2\alpha^2/\pi^2$ on $[0,\pi]$ give
$\hat z(b)\le\frac2\pi\int_0^\infty\alpha^2e^{-2b\alpha^2/\pi^2}d\alpha=(\pi/2)^{5/2}b^{-3/2}$.

\emph{Lower bound.} After gauge fixing there are $4L^4-L^3(L-1)=3L^4+L^3$ free links. Restrict each to $B_s:=\{\alpha(U)\le s\}$, $s:=\beta^{-1/2}\le1$. By \eqref{pf:eq-alpha}, $\alpha_p\le4s$ and $1-\cos\alpha_p\le\alpha_p^2/2\le8s^2$, so $\beta S\le6|\Lambda_L|\cdot8=48|\Lambda_L|$ there. Also $\mathrm{Haar}(B_s)=\frac2\pi\int_0^s\sin^2\alpha\,d\alpha\ge\frac2\pi\sin^21\int_0^s\alpha^2d\alpha=c_Hs^3$, since $\sin\alpha\ge\alpha\sin1$ on $[0,1]$. Hence $Z_L(\beta)\ge e^{-48|\Lambda_L|}(c_H\beta^{-3/2})^{3L^4+L^3}$.

\emph{Combination.} Inserting $b=\beta/2$,
\[
f_L(\beta/2)-f_L(\beta)\le48+3(1-\tfrac1L)\log\big((\pi/2)^{5/2}2^{3/2}\big)-(3+\tfrac1L)\log c_H+\tfrac6L\log\beta .
\]
The terms $\tfrac92\log\beta$ cancel and the remaining $\beta$-dependence is $\tfrac6L\log\beta$, which vanishes as $L\to\infty$ at fixed $\beta$. The numerical value of $D^*$ is a direct evaluation ($c_H=0.150258\ldots$).
\end{proof}

\begin{corollary}[uniform tail]\label{pf:cor-tail}
For every $\beta\ge1$, every $\mu\in\cL_\beta$, every plaquette $p$ and every $t\ge0$,
\[
\mu(O_p>t)\le e^{60.2-t/8}.
\]
Consequently $\sqrt\beta\,\alpha_p\le(\pi/\sqrt2)\sqrt{O_p}$ has Gaussian tails uniformly in $\beta\ge1$ and $\mu\in\cL_\beta$, and $\sup_{\beta\ge1}\sup_{\mu\in\cL_\beta}\mu(O_p^m)\le\int_0^\infty mt^{m-1}\min(1,e^{60.2-t/8})\,dt<\infty$ for every $m\ge1$.
\end{corollary}

\begin{proof}
By the lattice symmetries of $\mu_{\beta,L}$ it suffices to take $p=p_0$, which is a face of block $0$. Apply Lemma~\ref{pf:lem-chess} with $f=\mathbf 1\{O_{p_0}>t\}$. The image $\vartheta_tp_0$ is the $(1,2)$ plaquette spanning $[t_1,t_1+1]\times[t_2,t_2+1]$ at $x_3=t_3$ or $t_3+1$ and $x_0=t_0$ or $t_0+1$, whichever is even. So the images are the $|\Lambda_L|/4$ distinct $(1,2)$ plaquettes with $x_3,x_0$ even, each hit $4$ times, and since indicators are idempotent the product is the indicator that all these plaquettes have $O_p>t$. On that event $\beta S\ge t|\Lambda_L|/4$, so $e^{-\beta S}\le e^{-\beta S/2}e^{-t|\Lambda_L|/8}$ and
\[
\mu_{\beta,L}(O_{p_0}>t)\le\Big(e^{-t|\Lambda_L|/8}\,Z_L(\beta/2)/Z_L(\beta)\Big)^{1/|\Lambda_L|}=\exp\big(f_L(\beta/2)-f_L(\beta)-t/8\big).
\]
The event $\{O_p>t\}$ is open, so for $\mu=\lim\mu_{\beta,L_i}$ we get $\mu(O_p>t)\le\liminf_i\mu_{\beta,L_i}(O_p>t)\le e^{D^*-t/8}$ by Lemma~\ref{pf:lem-free}, and $D^*<60.2$. The inequality $\sqrt\beta\alpha_p\le(\pi/\sqrt2)\sqrt{O_p}$ is $1-\cos\alpha\ge2\alpha^2/\pi^2$. The moment bound is $\mu(O^m)=\int_0^\infty mt^{m-1}\mu(O>t)dt$.
\end{proof}

\subsection{Algebraic tools and the geometry of a box}\label{pf:sec-box}

\begin{lemma}[three estimates]\label{pf:lem-E}
\leavevmode
\begin{enumerate}
\item[(E1)] For $U\in\SU(2)$, $\|U-1\|_{HS}^2=4(1-\cos\alpha(U))$.
\item[(E2)] For Hermitian $2\times2$ matrices $A_1,\dots,A_k$, signs $\varepsilon_i=\pm1$ and $s:=\sum_i\|A_i\|_{op}$,
\[\big\|e^{i\varepsilon_1A_1}\cdots e^{i\varepsilon_kA_k}-1-i\textstyle\sum_i\varepsilon_iA_i\big\|_{op}\le e^s-1-s\le\tfrac{s^2}2e^s.\]
\item[(E3)] Let $x_1,\dots,x_4\in\R^3$, $U_p=e^{i\sigma\cdot x_1}e^{i\sigma\cdot x_2}e^{-i\sigma\cdot x_3}e^{-i\sigma\cdot x_4}$, $dx_p:=x_1+x_2-x_3-x_4$, $s_p:=\sum_i|x_i|<\pi$. Then $\alpha(U_p)\le s_p$ and $y_p:=\log U_p$ satisfies $|y_p-dx_p|\le s_p^2e^{s_p}$.
\end{enumerate}
\end{lemma}

\begin{proof}
(E1): $U-1=(\cos\alpha-1)+i\sin\alpha\,\hat n\cdot\sigma$, and the cross term is traceless, so $\|U-1\|_{HS}^2=2(1-\cos\alpha)^2+2\sin^2\alpha=4(1-\cos\alpha)$.
(E2): expand each exponential. The terms of total degree $0$ and $1$ are $1$ and $i\sum\varepsilon_iA_i$; the terms of degree $\ge2$ are bounded in norm, term by term, by the corresponding terms of $\prod_ie^{\|A_i\|}=e^s$. Finally $e^s-1-s=\sum_{m\ge2}s^m/m!\le\frac{s^2}2\sum_{m\ge2}s^{m-2}/(m-2)!$ because $m!\ge2(m-2)!$.
(E3): $\alpha(e^{\pm i\sigma\cdot x_i})=|x_i|$ and \eqref{pf:eq-alpha} give $|y_p|=\alpha(U_p)\le s_p<\pi$, so $U_p=e^{i\sigma\cdot y_p}$. Apply (E2) to the product and to $e^{i\sigma\cdot y_p}$, use $\|\sigma\cdot v\|_{op}=|v|$, and subtract: $|y_p-dx_p|\le\frac{s_p^2}2e^{s_p}+\frac{|y_p|^2}2e^{|y_p|}\le s_p^2e^{s_p}$.
\end{proof}

We also use three elementary inequalities: $1-\cos x\ge2x^2/\pi^2$ on $[-\pi,\pi]$; $\tfrac12(a-b)_+^2\ge a^2/4-b^2/2$ for $a,b\ge0$ (if $a\ge b$ the difference is $(a/2-b)^2$; otherwise the right side is negative); and $\|u+v\|^2\ge\tfrac12\|u\|^2-\|v\|^2$.

A $G$-connection on a $2$-complex $X$ of the lattice is an assignment of $U_{xy}\in G$ to its links with $U_{yx}=U_{xy}^{-1}$; it is \emph{flat} if $U_p=1$ for every plaquette of $X$, and \emph{pure gauge} if $U_{xy}=g_xg_y^{-1}$ for some $g$. The $2$-complex of a box $R$ consists of all links and plaquettes with all sites in $R$; the shell $2$-complex of all links and plaquettes with all sites in $\partial R$.

\begin{lemma}[flat connections on a box]\label{pf:lem-flat}
Let $G$ be any group, $d'\ge1$, and $R'$ a box in $\Z^{d'}$. Every flat $G$-connection on the $2$-complex of $R'$ is pure gauge. If $d'=4$, $R'=R$, and moreover $U=1$ on the shell links and on $T_{\rm int}$, then $U\equiv1$.
\end{lemma}

\begin{proof}
Order the directions as $1,\dots,d'$ and write $R'=\prod_i[a_i,b_i]$ with $a_i<b_i$. Let $T_{\rm ax}$ consist of the $e_i$-links $(x,x+e_i)$ with $x_l=a_l$ for all $l>i$; it is a spanning tree. Let $g_x$ be the holonomy along $T_{\rm ax}$ from the corner $(a_1,\dots,a_{d'})$ to $x$; then $U^g=1$ on $T_{\rm ax}$, and we show $U^g\equiv1$. Let $(x,x+e_j)$ be a link not in $T_{\rm ax}$, and let $k>j$ be the largest index with $x_k>a_k$. The links $(x-e_k,x)$ and $(x-e_k+e_j,x+e_j)$ are $e_k$-links whose coordinates $x_l$, $l>k$, equal $a_l$, so they lie in $T_{\rm ax}$. Flatness of the $(j,k)$ plaquette at $x-e_k$ gives $U^g_{x,j}=U^g_{x-e_k,j}$. Induct on $x_k$ down to $a_k$, then on the next largest index; the induction ends at a link of $T_{\rm ax}$. So $U^g\equiv1$, that is, $U_{xy}=g_x^{-1}g_y$, which is pure gauge with the gauge function $g^{-1}$.

For the second claim, write $U_{xy}=g_xg_y^{-1}$. $U=1$ on shell links makes $g$ constant on each connected component of the shell graph. The shell graph is connected: each face $F_{i,\pm}:=\partial R\cap\{x_i-z_i=\pm N\}$ is a connected $3$-box, every shell site lies in a face, and $F_{i,s}\cap F_{j,s'}\ne\emptyset$ for $i\ne j$, while $F_{i,+}$ and $F_{i,-}$ both meet $F_{j,+}$. So $g\equiv g_b$ on $\partial R$. Each tree of $T_{\rm int}$ is the time path from an interior site up to the unique shell site above it, and $U=1$ on it makes $g$ constant along it. Hence $g\equiv g_b$ and $U\equiv1$.
\end{proof}

\begin{lemma}[the shell is simply connected]\label{pf:lem-shell}
Let $G$ be any group. Every flat $G$-connection on the shell $2$-complex of $R$ is pure gauge. Consequently there is a finite integer $K_N$ such that, for every $\SU(2)$-connection on the shell and every shell link $l\notin T_s$, the variable of $l$ in the $T_s$-gauge (the gauge with $g_b=1$ and $U=1$ on $T_s$) is a product of at most $K_N$ conjugates of shell plaquette holonomies and their inverses; in particular its class angle is at most $K_N\max_q\alpha_q$, the maximum over shell plaquettes.
\end{lemma}

\begin{proof}
Every shell link and every shell plaquette lies in a face $F_{i,\pm}$. Take a plaquette in directions $(j,k)$ and suppose that no other coordinate satisfies $|x_i-z_i|=N$. The coordinate $x_j-z_j$ takes two consecutive values in $[-N,N]$, and one of them, $s$, has $|s|<N$. The two sites with $x_j-z_j=s$ then need $|x_k-z_k|=N$, but $x_k$ takes two consecutive values on them; this is impossible. The argument for links is the same.

\emph{Gluing.} Let $X=Y\cup Z$ be subcomplexes such that flat connections on $Y$ and on $Z$ are pure gauge and the $1$-skeleton of $Y\cap Z$ is connected. If $U$ is flat on $X$, write $U=g\cdot g^{-1}$ on $Y$ and $U=h\cdot h^{-1}$ on $Z$. On each link $xy$ of $Y\cap Z$, $h_x^{-1}g_x=h_y^{-1}g_y$, so $h^{-1}g\equiv c$ on $Y\cap Z$. Then $k:=g$ on $Y$ and $k:=hc$ on $Z$ agree on $Y\cap Z$ and give $U=k\cdot k^{-1}$ on $X$.

Each face is the $2$-complex of a $3$-box, so Lemma~\ref{pf:lem-flat} applies to it. Add the faces in the order $F_{1,-},F_{2,-},F_{3,-},F_{0,-},F_{1,+},F_{2,+},F_{3,+},F_{0,+}$. The intersection of a new face $F_{i,s}$ with the union of the earlier ones is the union of the $2$-boxes $F_{i,s}\cap F_{j,s'}$ over the earlier faces. For $s=-$ these $2$-boxes pairwise intersect. For $s=+$ they include $F_{i,+}\cap F_{j,-}$ for all $j\ne i$, which pairwise intersect, and each $F_{i,+}\cap F_{j,+}$ meets $F_{i,+}\cap F_{m,-}$ for $m\notin\{i,j\}$. So each intersection is connected, and gluing gives the first claim.

For the second claim, let $\Pi$ be the group with one generator for each oriented shell link, the relations $l^{-1}=$ reversed link, $l=1$ for $l\in T_s$, and $\partial q=1$ for every shell plaquette $q$. The tautological $\Pi$-connection is flat, hence pure gauge, $l=g_xg_y^{-1}$; since $g$ is constant along the spanning tree $T_s$, every generator is trivial in $\Pi$. So in the free group on the links not in $T_s$, each such link is a finite product of conjugates of the plaquette words $\partial q^{\pm1}$ (with tree links deleted). Let $K_N$ be the largest, over the finitely many links, of the least number of factors. Evaluating the identity on an $\SU(2)$-connection in $T_s$-gauge, the plaquette words evaluate to the plaquette holonomies in that gauge, which are conjugate to the gauge-invariant ones; \eqref{pf:eq-alpha} gives the bound.
\end{proof}

$K_N$ is finite by this argument; we do not compute it.

\begin{corollary}[zero set and nondegeneracy]\label{pf:cor-zero}
\leavevmode
\begin{enumerate}
\item Let $S_R(U):=\sum_{p\in P_R}(1-\cos\alpha_p)$. If $U=1$ on the shell links and on $T_{\rm int}$, then $S_R(U)=0$ only for $U\equiv1$ on $\Lambda'_R$.
\item \emph{(Dirichlet Maxwell form.)} For a real $1$-form $\varphi$ on $\Lambda''_R$, extended by $0$ to $T_{\rm int}$ and to the shell links, $Q(\varphi):=\sum_{p\in P_R}|(d\varphi)_p|^2$ is positive definite. Its smallest eigenvalue $\lambda_N$ is positive.
\end{enumerate}
\end{corollary}

\begin{proof}
A plaquette containing a link with an interior endpoint has all its sites in $R$, and every plaquette of $R$ is either in $P_R$ or a shell plaquette. (1) $S_R=0$ makes every $p\in P_R$ flat, and the shell plaquettes are flat because the shell links are $1$. Apply Lemma~\ref{pf:lem-flat}. (2) $Q(\varphi)=0$ means that $\varphi$ is a flat connection with group $(\R,+)$ on the $2$-complex of $R$, equal to $0$ on the shell links and on $T_{\rm int}$. Lemma~\ref{pf:lem-flat} gives $\varphi=0$.
\end{proof}

For orientation only: the computed values are $\lambda_1=0.837722$ and $\lambda_2=0.147122$; the proof uses $\lambda_N>0$ only.

\subsection{The conditional law in a box and the uniform Laplace limit}\label{pf:sec-L3}

Fix $R=z+[-N,N]^4$. Under $\mu\in\cL_\beta$, Lemma~\ref{pf:lem-limit}(3) with $\Lambda=\Lambda'_R$ shows that the conditional law of $U_{\Lambda'_R}$ given the links outside $\Lambda'_R$ is $\propto e^{-\beta S_R}\,dU_{\Lambda'_R}$. Every $p\in P_R$ has its links in $\Lambda'_R$ or among the shell links, so this law depends only on the shell links.

\emph{The gauge.} We fix the gauge in two steps that use no interior information. Let $g^s$ be the gauge transformation on shell sites with $g^s_b=1$ and $U^{g^s}=1$ on $T_s$; it is a function of the shell links only. Let $V:=U^{g^s}$ on the shell links, and let $\psi_l:=\sqrt\beta\log V_l$ for the shell links $l\notin T_s$ (and $\psi_l:=0$ on $T_s$). So $\psi$ is a function of the shell links only. Inside the conditional integral we apply Lemma~\ref{pf:lem-forest} with the forest $T_{\rm int}$ and $B=\partial R$: each tree of $T_{\rm int}$ meets $\partial R$ in exactly one site, and the gauge is $1$ there, so the shell links are unchanged.

Let $w$ be a bounded measurable function of the links of $R$ that is invariant under gauge transformations $g$ on $R$ with $g_b=1$, and let $\Psi_w$ be its conditional expectation. A gauge transformation at interior sites is absorbed by Haar invariance of $dU_{\Lambda'_R}$; at shell sites it acts on the argument of $\Psi_w$. Hence $\Psi_w(\text{shell})=\Psi_w(V)$, a function of $\psi$. With the exponential chart $U_l=e^{i\sigma\cdot\varphi_l/\sqrt\beta}$ on $\Lambda''_R$, and $J(x):=(\sin|x|/|x|)^2/(2\pi^2)$ the Haar density in that chart (normalised: $\int_{|x|<\pi}J=\frac{4\pi}{2\pi^2}\cdot\frac\pi2=1$), we obtain $\Psi_w=K_\beta w$ with
\begin{equation}\label{pf:eq-Kbeta}
K_\beta(d\varphi\,|\,\psi)=Z_\beta(\psi)^{-1}\exp\big(-\beta S_R(\varphi/\sqrt\beta;\psi/\sqrt\beta)\big)\prod_{l\in\Lambda''_R}J(\varphi_l/\sqrt\beta)\,\mathbf 1\{|\varphi_l|<\pi\sqrt\beta\}\,d\varphi,
\end{equation}
where $S_R(x;y)$ is $S_R$ with interior variables $e^{i\sigma\cdot x}$ on $\Lambda''_R$, $1$ on $T_{\rm int}$, and shell variables $e^{i\sigma\cdot y}$ off $T_s$, $1$ on $T_s$. The factors $\beta^{-3\nu_R/2}$ cancel, and $\{|x_l|=\pi\}$ is Haar-null.

\emph{Shell data from shell curvatures.} By Lemma~\ref{pf:lem-shell}, $|\psi_l|=\sqrt\beta\alpha(V_l)\le K_N\max_q\sqrt\beta\alpha_q\le K_N(\pi/\sqrt2)\max_q\sqrt{O_q}$, the maximum over the shell plaquettes $q$. With Corollary~\ref{pf:cor-tail} and a union bound, for every $M$,
\begin{equation}\label{pf:eq-epsM}
\mu(|\psi|>M)\le\varepsilon_R(M)\quad\text{for all }\beta\ge1,\ \mu\in\cL_\beta,\qquad\varepsilon_R(M)\xrightarrow[M\to\infty]{}0.
\end{equation}
Bounded shell data are thus implied by bounded gauge-invariant shell energies. No holonomy variables arise.

Write $Q(\varphi;\psi):=\sum_{p\in P_R}|(d\varphi+d\psi)_p|^2$ (colours summed), where $\varphi$ and $\psi$ are extended by $0$. Then $\tfrac12Q(\varphi;\psi)=\tfrac12\langle\varphi,Q\varphi\rangle+\langle\varphi,B\psi\rangle+\tfrac12\langle\psi,C\psi\rangle$ for fixed matrices $B,C$, and $Q$ acts on each colour as the form of Corollary~\ref{pf:cor-zero}(2). Let $K_\infty(\cdot|\psi)$ be the Gaussian measure $\propto e^{-\frac12Q(\varphi;\psi)}d\varphi$: three independent colours, each $N(-Q^{-1}B\psi,\,Q^{-1})$.

\begin{theorem}[uniform Laplace limit; L3]\label{pf:thm-L3}
Fix a box $R$ and $M<\infty$. There are $c>0$ and $\beta_0(M)<\infty$ such that, for all $\beta\ge\beta_0(M)$, all $|\psi|\le M$ and all $\varphi$ in the chart,
\begin{equation}\label{pf:eq-lower}
\beta S_R(\varphi/\sqrt\beta;\psi/\sqrt\beta)\ge c|\varphi|^2-7M^2 .
\end{equation}
Moreover
\[
\lim_{\beta\to\infty}\ \sup_{|\psi|\le M}\ \big\|K_\beta(\cdot|\psi)-K_\infty(\cdot|\psi)\big\|_{TV}=0,
\]
so $\sup_{|\psi|\le M}|K_\beta W(\psi)-K_\infty W(\psi)|\to0$ for every bounded measurable $W$, also when $W$ depends on $\psi$ or on further parameters with $\|W\|_\infty$ uniformly bounded. Finally, uniformly in $|\psi|\le M$ and $\beta\ge\beta_0(M)$, $K_\beta(|\varphi|>r\,|\,\psi)\le z_M^{-1}e^{7M^2}\int_{|\varphi|>r}e^{-c|\varphi|^2}d\varphi+o(1)$ as $\beta\to\infty$, with $z_M>0$.
The constant $c=\min(\lambda_N/16,\eta(\delta)/(2\nu_R\pi^2))$ contains $\eta(\delta)>0$, which is obtained by compactness and is not explicit.
\end{theorem}

\begin{proof}
\emph{Choice of $\delta$.} Fix $\delta\in(0,1]$ with $48\delta^2e^{8\delta}\le\min(\lambda_N/16,1)$. Write $x=\varphi/\sqrt\beta$ on the interior and $x=\psi/\sqrt\beta$ on the shell, $\tilde\varphi$ for $\varphi$ on the interior and $\psi$ on the shell, $s_p:=\sum_{i\in p}|x_i|$ and $\tilde s_p:=\sqrt\beta s_p$.

\emph{Near region} $\mathcal N:=\{|x_l|\le\delta\ \forall l\in\Lambda''_R\}$, with $\beta\ge M^2/\delta^2$ so that also every shell $|x_l|\le\delta$. For $p\in P_R$, $s_p\le4\delta$. By (E1), (E2), $\|\sigma\cdot v\|_{HS}=\sqrt2|v|$ and $\|A\|_{HS}\le\sqrt2\|A\|_{op}$,
$\|U_p-1\|_{HS}\ge\sqrt2\big(|dx_p|-\tfrac12s_p^2e^{s_p}\big)$. With $\beta(1-\cos\alpha_p)=\frac\beta4\|U_p-1\|^2_{HS}$ this gives $\beta(1-\cos\alpha_p)\ge\tfrac12(a_p-b_p)_+^2$ with $a_p:=|(d\tilde\varphi)_p|$ and $b_p:=\tfrac12\sqrt\beta s_p^2e^{s_p}\le2\delta e^{4\delta}\tilde s_p$. The second elementary inequality and $\tilde s_p^2\le4\sum_{i\in p}|\tilde\varphi_i|^2$ give
$\beta(1-\cos\alpha_p)\ge\tfrac14a_p^2-8\delta^2e^{8\delta}\sum_{i\in p}|\tilde\varphi_i|^2$. Each link lies in at most $6$ plaquettes, so
$\beta S_R\ge\tfrac14Q(\varphi;\psi)-48\delta^2e^{8\delta}(|\varphi|^2+|\psi|^2)$. The third elementary inequality, Corollary~\ref{pf:cor-zero}(2) and $|(d\psi)_p|^2\le4\sum_{i\in p}|\psi_i|^2$ give $Q(\varphi;\psi)\ge\tfrac12\lambda_N|\varphi|^2-24|\psi|^2$. Hence on $\mathcal N$
\[
\beta S_R\ge(\lambda_N/8-48\delta^2e^{8\delta})|\varphi|^2-(6+48\delta^2e^{8\delta})|\psi|^2\ge(\lambda_N/16)|\varphi|^2-7|\psi|^2 .
\]

\emph{Far region.} The set $\mathcal K_\delta:=\{U_{\Lambda''_R}:\max_l\alpha(U_l)\ge\delta\}$ is compact, and $S_R(\cdot\,;\text{shell}=1)$ is continuous and, by Corollary~\ref{pf:cor-zero}(1), positive on it. Let $\eta(\delta)>0$ be its minimum. By the last relation in \eqref{pf:eq-alpha}, $1-\cos$ being $1$-Lipschitz, and each shell link lying in at most $6$ plaquettes,
$|S_R(U;\psi/\sqrt\beta)-S_R(U;0)|\le6\sum_l|\psi_l|/\sqrt\beta\le6\sqrt{\nu_s}M/\sqrt\beta$. For $\beta\ge(12\sqrt{\nu_s}M/\eta)^2$ this is $\le\eta/2$, so $S_R\ge\eta/2$ on $\mathcal N^c$. Since $|\varphi|^2<\nu_R\pi^2\beta$ on the chart, $\beta S_R\ge(\eta/(2\nu_R\pi^2))|\varphi|^2$ there. With $\beta_0(M):=\max(M^2/\delta^2,144\nu_sM^2/\eta^2)$ this proves \eqref{pf:eq-lower}.

\emph{Pointwise limit.} Fix $r$ and let $|\varphi|\le r$, $|\psi|\le M$, $\beta$ large. By (E3), $|\sqrt\beta\alpha_p-|(d\tilde\varphi)_p||\le\sqrt\beta|y_p-dx_p|\le\tilde s_p^2e^{\tilde s_p/\sqrt\beta}/\sqrt\beta$, and $\sqrt\beta\alpha_p,|(d\tilde\varphi)_p|\le\tilde s_p\le4(r+M)$. Also $0\le\beta(\alpha^2/2-(1-\cos\alpha))\le\beta\alpha^4/24$. Hence
\[
\big|\beta S_R-\tfrac12Q(\varphi;\psi)\big|\le C_N(r+M)^3e^{4(r+M)/\sqrt\beta}/\sqrt\beta+C_N(r+M)^4/\beta,
\]
with $C_N$ depending on $N$ only. Moreover $(2\pi^2)^{\nu_R}\prod_lJ(\varphi_l/\sqrt\beta)\to1$ uniformly on $|\varphi|\le r$, and the indicator equals $1$ once $\beta>r^2/\pi^2$.

\emph{Convergence.} Let $f_\beta$ be the unnormalised density in \eqref{pf:eq-Kbeta} times $(2\pi^2)^{\nu_R}$, and $f_\infty:=e^{-\frac12Q(\varphi;\psi)}$. By the pointwise limit, $f_\beta\to f_\infty$ uniformly on $\{|\varphi|\le r,|\psi|\le M\}$. By \eqref{pf:eq-lower} and $J\le1/(2\pi^2)$, $f_\beta\le e^{7M^2}e^{-c|\varphi|^2}$; and $f_\infty\le e^{12M^2}e^{-\lambda_N|\varphi|^2/4}$. Splitting $\int=\int_{|\varphi|\le r}+\int_{|\varphi|>r}$ gives $\sup_{|\psi|\le M}\int|f_\beta-f_\infty|\,d\varphi\to0$. $Z_\infty(\psi):=\int f_\infty$ is continuous and positive, so $z_M:=\min_{|\psi|\le M}Z_\infty>0$. For $|W|\le1$, $|\int Wf_\beta/\int f_\beta-\int Wf_\infty/\int f_\infty|\le2\int|f_\beta-f_\infty|/\int f_\infty$, which gives the total-variation statement. The tail statement follows from the bound on $f_\beta$ and $\int f_\beta\ge z_M-o(1)$.
\end{proof}

\subsection{The limit curvature field}\label{pf:sec-L4}

Fix a base site $o$ and, for every plaquette $p$, a lattice path $\gamma_p$ from $o$ to $x_p$. Define the \emph{lasso field}
\[
F^\beta_p:=\sqrt\beta\,\mathrm{Ad}(U_{\gamma_p})\log U_p\in\R^3 .
\]
Then $|F^\beta_p|=\sqrt\beta\alpha_p$ exactly, and $O_p=\phi_\beta(|F^\beta_p|)$ with $\phi_\beta(t):=\beta(1-\cos(t/\sqrt\beta))$, $|\phi_\beta(t)-t^2/2|\le t^4/(24\beta)$. Under a gauge transformation, $F^\beta\mapsto\mathrm{Ad}(g_o)F^\beta$. No gauge tree is used to define $F^\beta$. We write $F^{\beta,(o)}$ when the base point matters, and $\mathcal P$ for the set of plaquettes of $\Z^4$.

\begin{proposition}[L4]\label{pf:prop-L4}
Let $\beta_k\ge1$, $\beta_k\to\infty$, and let $\mu_k\in\cL_{\beta_k}$ be arbitrary. Then there is a subsequence along which the finite-dimensional laws of $F^{\beta_k}$ under $\mu_k$ converge to the finite-dimensional laws of a probability measure $\mu_\infty$ on $(\R^3)^{\mathcal P}$. Every $\mu_\infty$ obtained in this way, along any subsequence, has the following properties; $F$ denotes the coordinate field.
\begin{enumerate}
\item[(a)] $\mu_\infty$ does not depend on the choice of $o$ and of the paths $\gamma_p$, in the sense that the same subsequence gives the same limit for every other choice.
\item[(b)] \emph{(Bianchi.)} $(dF)_c=0$ for every $3$-cube $c$, $\mu_\infty$-almost surely.
\item[(c)] \emph{(Invariances.)} $\mu_\infty$ is invariant under $F\mapsto\mathrm{Ad}(\mathcal R)F$ for each $\mathcal R\in\SU(2)$, and under lattice translations.
\item[(d)] \emph{(Moments.)} $\mu_\infty(\tfrac12|F_p|^2>t)\le e^{60.2-t/8}$ for all $p,t$, so $F$ has moments of all orders. Along the subsequence, $(O_p)_{p\in\mathcal P}$ under $\mu_k$ converges in finite-dimensional laws to $(\tfrac12|F_p|^2)_{p\in\mathcal P}$, and $\mu_k(O_p)\to\E_\infty\tfrac12|F_p|^2$ and $\mu_k(O_pO_q)\to\E_\infty\tfrac14|F_p|^2|F_q|^2$ for all $p,q$.
\item[(e)] \emph{(Gaussian DLR.)} For every box $R$, let $\cG_R:=\sigma(F_q:q\notin P_R)$, $F_{\rm in}:=(F_p)_{p\in P_R}$ and $F_{\rm sh}:=(F_q)_{q\text{ shell plaquette}}$, which is $\cG_R$-measurable. Let $\mathsf L$ be a fixed linear left inverse of $d$ on the shell $1$-forms vanishing on $T_s$. Then a regular conditional law of $F_{\rm in}$ given $\cG_R$ is the law of $d(A+\Psi)|_{P_R}$, where $\Psi:=\mathsf LF_{\rm sh}$ and $A\sim K_\infty(\cdot|\Psi)$.
\end{enumerate}
\end{proposition}

The map $d$ is injective on shell $1$-forms vanishing on $T_s$: a form in the kernel is a flat $(\R,+)$-connection on the shell, hence $d\lambda$ by Lemma~\ref{pf:lem-shell}, and $\lambda$ is constant along the spanning tree $T_s$, so the form is $0$. Hence $\mathsf L$ exists.

\begin{proof}[Proof of existence and of (a)]
By Corollary~\ref{pf:cor-tail}, $|F^\beta_p|=\sqrt\beta\alpha_p\le(\pi/\sqrt2)\sqrt{O_p}$ is tight uniformly in $\beta\ge1$ and $\mu\in\cL_\beta$. So every finite family $(F^{\beta_k}_p)_{p\in\mathcal F}$ is tight. Prokhorov's theorem and diagonal extraction over the countably many finite families give a subsequence along which all finite-dimensional laws converge; they are consistent, and Kolmogorov's extension theorem gives $\mu_\infty$.

\emph{Paths.} Let $\gamma'_p$ be another path from $o$ to $x_p$ and $F'$ the corresponding field. Then $F'_p=\mathrm{Ad}(V)F_p$, where $V=U_{\gamma'_p}U_{\gamma_p}^{-1}$ is the holonomy of a closed path in a box. By Lemma~\ref{pf:lem-flat} applied to the tautological connection of the box (as in the proof of Lemma~\ref{pf:lem-shell}), the loop word is a product of a fixed number $m$ of conjugates of plaquette words $q_1,\dots,q_m$, so $\alpha(V)\le\sum_i\alpha_{q_i}\le(\pi/\sqrt2)\sum_i\sqrt{O_{q_i}/\beta}$. Corollary~\ref{pf:cor-tail} makes this $O_P(\beta^{-1/2})$ uniformly in $\mu\in\cL_\beta$, and $|F'_p-F_p|\le2\alpha(V)|F_p|\to0$ in probability. By Slutsky's lemma the limits coincide.

\emph{Base point.} Let $b\ne o$, let $W$ be the holonomy of a fixed path from $o$ to $b$, and let $H:=F^{(b)}$ (paths from $b$) on a finite family. Then $F^{(o)}$ with paths through $b$ equals $\mathrm{Ad}(W)H$. Let $g$ be independent Haar variables at the finitely many sites involved (and $g=1$ elsewhere), independent of $U$, and $\tilde U:=U^g$. By Lemma~\ref{pf:lem-limit}(1), $\tilde U\sim\mu$. Now $W(\tilde U)=g_oWg_b^{-1}$ and $H(\tilde U)=\mathrm{Ad}(g_b)H(U)$. Given $(U,g_b)$, $g_oWg_b^{-1}$ is Haar distributed, so $W(\tilde U)$ is Haar and independent of $H(\tilde U)$, whose law is that of $H$. Hence $\mathrm{law}(\mathrm{Ad}(W)H)=\mathrm{law}(\mathrm{Ad}(\mathcal H)H)$ with $\mathcal H$ Haar and independent of $H$. Since $\mathrm{law}(H)$ is invariant under each fixed $\mathrm{Ad}(\mathcal R)$ (constant gauge transformations, Lemma~\ref{pf:lem-limit}(1)), this is $\mathrm{law}(H)$. So at every $\beta$ the field based at $o$ (paths through $b$) and the field based at $b$ have the same law. With the path statement, the limit does not depend on the base point.
\end{proof}

\begin{proof}[Proof of (c)]
Colour invariance holds at every $\beta$ by gauge invariance under constant $g\equiv\mathcal R$. For translations, translation invariance of $\mu_k$ maps the field based at $o$ with paths $\gamma_p$ to the field based at $o+a$ with paths $\gamma_p+a$ on the plaquettes $p+a$. By (a), changing the base point back to $o$ and the paths to $\gamma_{p+a}$ does not change the limit. So $(F_{p+a})_p$ and $(F_p)_p$ have the same limit law.
\end{proof}

\begin{proof}[Proof of (b)]
\emph{The cube word identity.} Let $c=(x;\lambda,\mu,\nu)$ with $\lambda<\mu<\nu$; write $1,2,3$ for $\lambda,\mu,\nu$, $U_{ab}(y)$ for the holonomy of the $(a,b)$ plaquette with corner $y$ (as in \S\ref{pf:sec-setting}), and $u_a:=U_{x,a}$. Then, for every $\SU(2)$-connection on the twelve edges of $c$,
\begin{equation}\label{pf:eq-cube}
U_{23}(x)^{-1}\cdot u_2U_{13}(x+e_2)^{-1}u_2^{-1}\cdot U_{12}(x)^{-1}\cdot u_1U_{23}(x+e_1)u_1^{-1}\cdot U_{13}(x)\cdot u_3U_{12}(x+e_3)u_3^{-1}=1 .
\end{equation}
This is an identity in the free group on the twelve edges: writing each factor as a word in the edges, the product reduces to the empty word. To check it, label the edge from $x+\epsilon$ in direction $a$ by $a_\epsilon$, $\epsilon\in\{0,1\}^3$. The six words are
\begin{align*}
U_{23}(x)^{-1}&=3_{000}2_{001}3_{010}^{-1}2_{000}^{-1}, &
u_2U_{13}(x+e_2)^{-1}u_2^{-1}&=2_{000}3_{010}1_{011}3_{110}^{-1}1_{010}^{-1}2_{000}^{-1},\\
U_{12}(x)^{-1}&=2_{000}1_{010}2_{100}^{-1}1_{000}^{-1}, &
u_1U_{23}(x+e_1)u_1^{-1}&=1_{000}2_{100}3_{110}2_{101}^{-1}3_{100}^{-1}1_{000}^{-1},\\
U_{13}(x)&=1_{000}3_{100}1_{001}^{-1}3_{000}^{-1}, &
u_3U_{12}(x+e_3)u_3^{-1}&=3_{000}1_{001}2_{101}1_{011}^{-1}2_{001}^{-1}3_{000}^{-1}.
\end{align*}
Concatenating and cancelling adjacent inverse pairs from left to right removes all letters. The abelianised identity, with signs $-,-,-,+,+,+$ on $U_{23}(x),U_{13}(x+e_2),U_{12}(x),U_{23}(x+e_1),U_{13}(x),U_{12}(x+e_3)$, is the boundary of $c$ with the incidence numbers of $d$ in \S\ref{pf:sec-setting}.

\emph{The estimate.} Write $y_f:=\log U_f$ for the six faces $f$ (based at their corners), $\varepsilon_f$ for the sign of $f$ in \eqref{pf:eq-cube}, and $W_f\in\{1,u_1,u_2,u_3\}$ for its transport. Each factor of \eqref{pf:eq-cube} is $e^{i\varepsilon_f\sigma\cdot\mathrm{Ad}(W_f)y_f}$ whenever $\alpha_f<\pi$. Since the product is $1$, (E2) gives $|\sum_f\varepsilon_f\mathrm{Ad}(W_f)y_f|\le\frac{s^2}2e^s$ with $s=\sum_f\alpha_f$. Multiply by $\sqrt\beta$ and apply the isometry $\mathrm{Ad}(U_{\gamma_x})$, $\gamma_x$ a fixed path from $o$ to $x$. Then $\mathrm{Ad}(U_{\gamma_x}W_f)\sqrt\beta y_f$ is the lasso field of $f$ with the path $\gamma_x$ followed by the edge of $W_f$; by the path estimate in the proof of (a) it differs from $F^\beta_f$ by at most $2\alpha(V_f)|F^\beta_f|=O_P(\beta^{-1/2})$. Hence, with $S_c:=\sum_f|F^\beta_f|$,
\[
|(dF^\beta)_c|\le\frac{S_c^2}{2\sqrt\beta}e^{S_c/\sqrt\beta}+O_P(\beta^{-1/2}),
\]
on the event $\max_f\alpha_f<\pi$, whose probability tends to $1$ by Corollary~\ref{pf:cor-tail}. $S_c$ is tight, so $(dF^{\beta_k})_c\to0$ in probability. The set $\{|(dF)_c|>\epsilon\}$ is open, so $\mu_\infty(|(dF)_c|>\epsilon)\le\liminf_k\mu_k(|(dF^{\beta_k})_c|>\epsilon)=0$ for every $\epsilon>0$. There are countably many cubes.
\end{proof}

\begin{proof}[Proof of (d)]
$O_p=\phi_{\beta}(|F^\beta_p|)$. For a bounded Lipschitz $h$ of finitely many coordinates $p\in\mathcal F$,
\[
\big|\E_kh(\phi_{\beta_k}(|F^{\beta_k}|))-\E_kh(\tfrac12|F^{\beta_k}|^2)\big|\le\mathrm{Lip}(h)\,|\mathcal F|\sup_{t\le r}\frac{t^4}{24\beta_k}+2\|h\|_\infty\,\mu_k\big(\max_{\mathcal F}|F^{\beta_k}_p|>r\big),
\]
which tends to $0$ as $k\to\infty$ and then $r\to\infty$. Since $F\mapsto h(\tfrac12|F|^2)$ is bounded and continuous, $\E_kh(\tfrac12|F^{\beta_k}|^2)\to\E_\infty h(\tfrac12|F|^2)$; bounded Lipschitz functions determine convergence in law. For the tail bound: the law of $\tfrac12|F_p|^2$ under $\mu_\infty$ is the limit in law of the law of $O_p$ under $\mu_k$, just shown; the set $\{y>t\}$ is open, so $\mu_\infty(\tfrac12|F_p|^2>t)\le\liminf_k\mu_k(O_p>t)\le e^{60.2-t/8}$ by Corollary~\ref{pf:cor-tail}. By Corollary~\ref{pf:cor-tail}, $\sup_k\mu_k(O_p^4)<\infty$, so $O_p$, $O_pO_q$ (with $\mu_k(O_p^2O_q^2)\le\mu_k(O_p^4)^{1/2}\mu_k(O_q^4)^{1/2}$) are uniformly integrable, and their expectations converge.
\end{proof}

\begin{proof}[Proof of (e)]
Fix $R$ and $b\in\partial R$, the root of $T_s$. Let $T:=T_s\cup T_{\rm int}$.

\emph{Step 1: adapted paths.} By (a) we may compute with the base point $b$ and with any paths. For $p\in P_R$ choose $\gamma_p$ inside $T$: along $T_s$ to a shell site, then down the tree of $T_{\rm int}$. For $q\notin P_R$ choose $\gamma_q$ through links not in $\Lambda'_R$: a plaquette $q\notin P_R$ has no interior site, and the graph of non-interior sites with the links not in $\Lambda'_R$ is connected. For a shell plaquette choose $\gamma_q$ inside $T_s$. Then $F^\beta_{\rm out}:=(F^\beta_q)_{q\notin P_R}$ is a function of the links outside $\Lambda'_R$. In the box gauge ($g^s$ on the shell, then the $T_{\rm int}$ gauge, both with $g_b=1$), $F^{\beta,(b)}$ is unchanged and $U_{\gamma_p}=1$ for $p\in P_R$; so $F^\beta_{\rm in}=F_{\rm in}(\varphi,\psi):=(\sqrt\beta\log U_p)_{p\in P_R}$, and by (E3)
\begin{equation}\label{pf:eq-Fin}
|F_{\rm in}(\varphi,\psi)-(d\varphi+d\psi)|_{P_R}|\le C_R(|\varphi|+|\psi|)^2e^{4(|\varphi|+|\psi|)/\sqrt\beta}/\sqrt\beta
\end{equation}
when $4(|\varphi|+|\psi|)<\pi\sqrt\beta$. Likewise $F^\beta_{\rm sh}=(\sqrt\beta\log V_q)_q$, so $|F^\beta_{\rm sh}-d\psi|\le C_R|\psi|^2e^{4|\psi|/\sqrt\beta}/\sqrt\beta$.

\emph{Step 2: finite-$\beta$ DLR.} Let $w_{\rm in}$ be bounded Lipschitz on $(\R^3)^{P_R}$ and $w_{\rm out}$ bounded Lipschitz in finitely many $F_q$, $q\notin P_R$. The function $w_{\rm in}(F^{\beta,(b)}_{\rm in})$ is invariant under gauge transformations with $g_b=1$. By \S\ref{pf:sec-L3} and Lemma~\ref{pf:lem-limit}(3),
$\E_k[w_{\rm in}(F^\beta_{\rm in})w_{\rm out}]=\E_k[(K_\beta W_\beta)(\psi)\,w_{\rm out}]$ with $W_\beta(\varphi,\psi):=w_{\rm in}(F_{\rm in}(\varphi,\psi))$.

\emph{Step 3: Laplace limit.} Let $\Phi(\psi):=K_\infty[w_{\rm in}(d\varphi+d\psi)](\psi)$. For $|\psi|\le M$,
$|K_\beta W_\beta-\Phi|\le|K_\beta W_\beta-K_\beta W_\infty|+|K_\beta W_\infty-K_\infty W_\infty|$ with $W_\infty(\varphi,\psi):=w_{\rm in}(d\varphi+d\psi)$. The second term tends to $0$ uniformly by Theorem~\ref{pf:thm-L3}. The first is at most $\mathrm{Lip}(w_{\rm in})\sup_{|\varphi|\le r}(\text{right side of }\eqref{pf:eq-Fin})+2\|w_{\rm in}\|_\infty K_\beta(|\varphi|>r|\psi)$, which is small uniformly by the tightness statement of Theorem~\ref{pf:thm-L3}. Under $K_\infty(\cdot|\psi)$, $d\varphi+d\psi=D\zeta+E\psi$ with $\zeta\sim N(0,Q^{-1})$ per colour and fixed linear maps $D,E$; so $\Phi$ is bounded and Lipschitz.

\emph{Step 4: tail.} Off $\{|\psi|\le M\}$ the difference is at most $2\|w_{\rm in}\|_\infty\|w_{\rm out}\|_\infty\varepsilon_R(M)$ by \eqref{pf:eq-epsM}, uniformly in $k$.

\emph{Step 5: shell data from curvatures.} $\psi=\mathsf Ld\psi$ exactly. So on $\{|\psi|\le M\}$,
\[|\psi-\mathsf LF^{\beta}_{\rm sh}|\le\|\mathsf L\|\,C_RM^2e^{4M/\sqrt\beta}/\sqrt\beta .\] Since $\Phi$ is Lipschitz, $|\E_k[\Phi(\psi)w_{\rm out}]-\E_k[\Phi(\mathsf LF^\beta_{\rm sh})w_{\rm out}]|\le o(1)+2\|\Phi\|_\infty\|w_{\rm out}\|_\infty\varepsilon_R(M)$.

\emph{Step 6: limit.} The laws of the adapted field under $\mu_k$ converge to $\mu_\infty$ along the subsequence: at each $k$ the law of the field based at $b$ equals that of the field based at $o$ with paths through $b$, which differs from $F^{\beta_k}$ by $o_P(1)$ (proof of (a)); apply Slutsky's lemma. The functions $w_{\rm in}(F_{\rm in})w_{\rm out}$ and $\Phi(\mathsf LF_{\rm sh})w_{\rm out}$ are bounded and continuous in finitely many coordinates. Hence
$|\E_\infty[w_{\rm in}(F_{\rm in})w_{\rm out}]-\E_\infty[\Phi(\mathsf LF_{\rm sh})w_{\rm out}]|\le C\varepsilon_R(M)$ for every $M$, and so equality holds.

\emph{Step 7: conditional law.} The functions $w_{\rm out}$ generate $\cG_R$ and form a class closed under products, so by the monotone class theorem $\E_\infty[w_{\rm in}(F_{\rm in})|\cG_R]=\Phi_{w_{\rm in}}(\mathsf LF_{\rm sh})$ a.s. A regular conditional law of $F_{\rm in}$ given $\cG_R$ exists because $(\R^3)^{P_R}$ is Polish. A countable family of bounded Lipschitz $w_{\rm in}$ determines probability measures on $(\R^3)^{P_R}$, so this regular conditional law equals the law of $d(A+\Psi)|_{P_R}$, $A\sim K_\infty(\cdot|\Psi)$, $\Psi=\mathsf LF_{\rm sh}$, almost surely.
\end{proof}

At finite $\beta$ the conditioning is on links modulo gauge; in the limit it is on curvatures only. Lemma~\ref{pf:lem-shell} and the injectivity of $d$ on $T_s$-gauge shell forms show that nothing is lost: shell curvatures determine shell links modulo gauge, exactly for the group and linearly in the limit.

\subsection{The increasing Dirichlet covariances}\label{pf:sec-D}

In this subsection and the next, $\mu_\infty$ is a limit as in Proposition~\ref{pf:prop-L4}, $\E$ is its expectation, and $F$ is its field. For a box $R$ let $h^R_p:=\E[F_p|\cG_R]$; so $h^R_q=F_q$ for $q\notin P_R$.
Let $P$ be the orthogonal projection in $\ell^2(\mathcal P)$ (one colour) onto the closure of $\{dA:A\text{ a finitely supported real }1\text{-form}\}$.

\begin{definition}\label{pf:def-G}
$G(n):=\langle\delta_{p_0},P\delta_{p_n}\rangle$ for $n\in\Z$, and $c(n):=(G(n)/G(0))^2$.
\end{definition}

\begin{lemma}[Dirichlet covariances; Lemma D]\label{pf:lem-D}
Let $R$ be a box. Given $\cG_R$, the field $\xi^R:=F-h^R$ is centred Gaussian, with three independent colours, each with the deterministic covariance $C^R$ given by
\begin{equation}\label{pf:eq-CR}
\langle f,C^Rf\rangle=\sup\Big\{\frac{\langle f,dA\rangle^2}{|dA|^2}:A\text{ a real }1\text{-form supported on }\Lambda'_R,\ dA\ne0\Big\}
\end{equation}
for finitely supported $f\in\ell^2(\mathcal P)$. Moreover $(d^*h^R)_l=0$ for every $l\in\Lambda'_R$. For $R_N=[-N,N]^4$, $\langle f,C^{R_N}f\rangle$ is nondecreasing in $N$ and converges to $|Pf|^2$.
\end{lemma}

\begin{proof}
By Proposition~\ref{pf:prop-L4}(e), given $\cG_R$, $F_{\rm in}=d(A+\Psi)|_{P_R}$ with $A\sim K_\infty(\cdot|\Psi)$: per colour $A=a(\Psi)+\zeta$ with $a(\Psi)=-Q^{-1}B\Psi$ and $\zeta\sim N(0,Q^{-1})$ independent of $\cG_R$. So $h^R|_{P_R}=d(a(\Psi)+\Psi)$, and $\xi^R=D\zeta$ with $D$ the restriction of $d$ to $1$-forms on $\Lambda''_R$. Its covariance per colour is $C^R=DQ^{-1}D^{\mathsf T}$, where $Q=D^{\mathsf T}D$. For $f$ supported on $P_R$, Cauchy--Schwarz in the inner product $\langle\cdot,Q\cdot\rangle$ gives
$\langle f,DQ^{-1}D^{\mathsf T}f\rangle=\sup_{\varphi\ne0}\langle f,D\varphi\rangle^2/\langle\varphi,Q\varphi\rangle$. A $1$-form $A$ on $\Lambda'_R$ can be changed to vanish on $T_{\rm int}$ by adding $d\lambda$ with $\lambda$ supported on interior sites (solve along each tree of $T_{\rm int}$ from its shell site); this does not change $dA$. So the supremum over $\varphi$ on $\Lambda''_R$ equals the supremum in \eqref{pf:eq-CR}. Since $dA$ is supported on $P_R$, \eqref{pf:eq-CR} holds for all finitely supported $f$ (and $\xi^R_q=0$ for $q\notin P_R$).

The mean $a(\Psi)+\Psi$ minimises $\varphi\mapsto\sum_{p\in P_R}|(d(\varphi+\Psi))_p|^2$ over $\varphi$ on $\Lambda''_R$. By the same gauge remark it also minimises over $\varphi$ on all of $\Lambda'_R$. Every plaquette containing $l\in\Lambda'_R$ lies in $P_R$, so the stationarity condition in the variable $\varphi_l$ is $(d^*h^R)_l=0$.

$R_N\subset R_{N+1}$ gives $\Lambda'_{R_N}\subset\Lambda'_{R_{N+1}}$, so \eqref{pf:eq-CR} is nondecreasing. The union of the $\Lambda'_{R_N}$ is all links, and the supremum of $\langle f,v\rangle^2/|v|^2$ over a dense subset of a closed subspace is $|Pf|^2$.
\end{proof}

For orientation only: $C^{R_N}_{p_0p_0}=7/24=0.29167$, $0.49267$, $0.49916$ for $N=1,2,3$, where $R_N=[-N,N]^4$ has side $2N$; the limit is $G(0)=\tfrac12$.

\subsection{Decomposition of a limit point}\label{pf:sec-R}

\begin{lemma}[free field plus constant background; Lemma R]\label{pf:lem-R}
Under $\mu_\infty$, $F=\xi+f$, where:
$\xi$ is centred Gaussian with three independent colours, each with covariance $P$;
$f$ is a random constant, that is, $f_p=f_{\mu\nu}\in\R^3$ for every plaquette $p$ of type $(\mu,\nu)$;
$\xi$ and $f$ are independent; and $\E|f_{\mu\nu}|^4<\infty$.
The colours of $f$ may be correlated.
\end{lemma}

\begin{proof}
\emph{Martingale limit.} For $N'\ge N$, $P_{R_N}\subset P_{R_{N'}}$, so $\cG_{R_{N'}}\subset\cG_{R_N}$. Let $\cG_\infty:=\bigcap_N\cG_{R_N}$ and $h^N:=h^{R_N}$. Since $\E|F_p|^2<\infty$ (Proposition~\ref{pf:prop-L4}(d)), L\'evy's downward theorem gives $h^N_p\to h_p:=\E[F_p|\cG_\infty]$ almost surely and in $L^2$. By conditional Jensen, $\E|h_p|^4\le\E|F_p|^4<\infty$.

\emph{Gaussian part.} Let $\xi:=F-h$. For finitely supported $f'$ (all colours) and bounded $\cG_\infty$-measurable $Y$, Lemma~\ref{pf:lem-D} gives $\E[e^{i\langle f',\xi^{R_N}\rangle}Y]=e^{-\frac12\sum_a\langle f'^a,C^{R_N}f'^a\rangle}\E Y$. Let $N\to\infty$: $\xi^{R_N}\to\xi$ in $L^2$ and $\langle f'^a,C^{R_N}f'^a\rangle\to|Pf'^a|^2$. So $\E[e^{i\langle f',\xi\rangle}Y]=e^{-\frac12\sum_a|Pf'^a|^2}\E Y$. Hence $\xi$ is centred Gaussian with independent colours of covariance $P$, and independent of $\cG_\infty$, hence of $h$.

\emph{Closed and co-closed.} $(dh)_c=\E[(dF)_c|\cG_\infty]=0$ by Proposition~\ref{pf:prop-L4}(b). For a link $l$, $(d^*h^{N})_l=0$ once $l\in\Lambda'_{R_N}$ (Lemma~\ref{pf:lem-D}); the $L^2$ limit gives $(d^*h)_l=0$.

\emph{Stationarity.} Since $\xi$ is centred and independent of $h$, $\E[h^a_ph^a_q]=\E[F^a_pF^a_q]-P_{pq}$. Both terms are invariant under lattice translations (Proposition~\ref{pf:prop-L4}(c), and $P$ commutes with translations). So $h$ is second-order stationary.

\emph{Constancy.} The following argument for constancy was supplied by an independent check of this lemma. On $\Z^4$ the Hodge Laplacian $dd^*+d^*d$ on $2$-forms acts on each component $\omega_{\mu\nu}$ as the scalar operator $8-\sum_{e}\mathsf S_e$, the sum over the eight unit vectors $e=\pm e_\lambda$ and $\mathsf S_e\omega(x)=\omega(x+e)$. (Both sides are finite-range, translation-invariant operators; on finitely supported forms their Fourier symbols agree, because $\hat d=\sum_\lambda\hat e_\lambda\,e^\lambda\wedge$ and $\hat d^*=\sum_\lambda\overline{\hat e_\lambda}\,\iota_\lambda$ with $\hat e_\lambda=e^{ik_\lambda}-1$, and $e^\lambda\wedge\iota_\kappa+\iota_\kappa e^\lambda\wedge=\delta_{\lambda\kappa}$, so the symbol is $\sum_\lambda|\hat e_\lambda|^2=\sum_\lambda(2-2\cos k_\lambda)$.) Since $dh=0$ and $d^*h=0$ almost surely, every component is harmonic: $\sum_eh(x+e)=8h(x)$ for each fixed type $(\mu,\nu)$ and colour, writing $h(x)$ for that component at the plaquette with corner $x$. Fix $x$ and put $u:=h(x)$, $v_e:=h(x+e)$ in real $L^2(\mu_\infty)$. By stationarity, $\|v_e\|=\|u\|$ for all $e$. Hence
\[
\sum_e\|v_e-u\|^2=\sum_e\|v_e\|^2-2\Big\langle\sum_ev_e,u\Big\rangle+8\|u\|^2=16\|u\|^2-16\|u\|^2=0 .
\]
So $h(x+e)=h(x)$ almost surely for every $x$ and $e$. There are countably many pairs, so almost surely $h$ is constant in $x$ for each type and colour. Put $f:=h$.
\end{proof}

A constant background matters: for a background with $f_{p_n}=-f_{p_0}$ the covariance ratio can be negative.

\subsection{The free correlation}\label{pf:sec-G}

Write $\hat k_\lambda^2:=4\sin^2(k_\lambda/2)$ and $\hat k^2:=\sum_{\lambda=0}^3\hat k_\lambda^2$, and for $\vec k=(k_1,k_2,k_3)\in[-\pi,\pi]^3$ let $a(\vec k):=\sum_{i=1}^3\hat k_i^2$ and $m(\vec k)>0$ with $\cosh m=1+a/2$.

\begin{lemma}[the free correlation; Lemma G]\label{pf:lem-G}
\leavevmode
\begin{enumerate}
\item $G(n)=\displaystyle\int_{[-\pi,\pi]^4}\frac{\hat k_1^2+\hat k_2^2}{\hat k^2}\,e^{ik_0n}\,\frac{dk}{(2\pi)^4}=\int_{[-\pi,\pi]^3}\frac{\hat k_1^2+\hat k_2^2}{2\sinh m}\,e^{-m|n|}\,\frac{d\vec k}{(2\pi)^3}$.
\item $G(0)=\tfrac12$, and $0<G(n+1)<G(n)\le\tfrac12$ for $n\ge0$.
\item $G(n)\ge\frac4{\pi^2}\cdot\frac{16r^5}3\cdot\frac{e^{-1}}{2\sinh1}\cdot\frac1{(2\pi)^3}$ with $r=1/(\sqrt3n)$, for $n\ge1$.
\item For all $n\ge1$, with $m_1:=\mathrm{arccosh}(1+2/\pi^2)$,
\[\Big|G(n)-\frac1{\pi^2n^4}\Big|\le\frac{14400}{\pi^2n^6}+\frac12e^{-m_1n}+\frac1{6\pi^2}\int_1^\infty r^3e^{-nr}\,dr ;\]
in particular $G(n)=1/(\pi^2n^4)+O(n^{-6})$. We do not convert the two exponentially small terms into a numerical multiple of $n^{-6}$.
\end{enumerate}
Consequently $c(n)=4G(n)^2=\frac4{\pi^4n^8}\big(1+O(n^{-2})\big)$, and $-\log c(n)/n\to0$.
\end{lemma}

\begin{proof}
(1) In Fourier, with $\hat d$ as above, $\hat d\hat d^*+\hat d^*\hat d=\hat k^2$ on forms of every degree. For $k\ne0$, the range of $\hat d$ on $2$-forms equals the kernel of $\hat d$, and $\hat d\hat d^*/\hat k^2$ is the identity on it and vanishes on the range of $\hat d^*$; so $\hat P(k)=\hat d\hat d^*/\hat k^2$. The closure of the range of $d$ is the set of $\omega\in\ell^2$ with $\hat\omega(k)\in\mathrm{ran}\,\hat d(k)$ for almost every $k$: this set is closed and contains the range; conversely, for such $\omega$ and $\varepsilon>0$, $\mathbf 1\{|k|>\varepsilon\}\hat\omega=\hat d\hat A_\varepsilon$ with $\hat A_\varepsilon:=\mathbf 1\{|k|>\varepsilon\}\hat d^*\hat\omega/\hat k^2\in L^2$, and $A_\varepsilon$ is an $\ell^2$-limit of finitely supported forms, on which $d$ is bounded. So $P$ is multiplication by $\hat P(k)$. Its diagonal entry at the type $(1,2)$ is $|\hat e_1|^2+|\hat e_2|^2=\hat k_1^2+\hat k_2^2$. Since $p_n=p_0+ne_0$, Plancherel gives the first formula. For $a>0$, $\int_{-\pi}^\pi e^{ik_0n}/(2+a-2\cos k_0)\,dk_0/(2\pi)=e^{-m|n|}/(2\sinh m)$ (residues), which gives the second; the integrand is integrable because $\hat k_1^2+\hat k_2^2\le a$ and $\sinh m\ge m\ge c\sqrt a$ near $0$.

(2) By the symmetry $k_1\leftrightarrow k_3$ and so on, $\int(\hat k_1^2+\hat k_2^2)/\hat k^2=2\cdot\frac14$. The integrand of the second formula is positive almost everywhere and strictly decreasing in $|n|$.

(3) On the cube $|k_i|\le r$ we have $a\le|\vec k|^2\le1/n^2$. Since $a=4\sinh^2(m/2)\ge m^2$, $m\le1/n$, so $e^{-mn}\ge e^{-1}$ and $\sinh m\le\sinh1$. Also $\hat k_i^2\ge4k_i^2/\pi^2$, and $\int_{\text{cube}}(k_1^2+k_2^2)\,d\vec k=16r^5/3$. Drop the rest of the (positive) integrand.

(4) Let $g:=(\hat k_1^2+\hat k_2^2)/(2\sinh m)$ and $g_0:=(k_1^2+k_2^2)/(2|\vec k|)\le|\vec k|/2$. A direct computation in polar coordinates gives $\int_{\R^3}g_0e^{-n|\vec k|}d\vec k/(2\pi)^3=\frac{8\pi/3}{2(2\pi)^3}\int_0^\infty r^3e^{-nr}dr=\frac1{\pi^2n^4}$. For $|\vec k|\le1$ we use: $0\le x^2-\hat x^2\le x^4/12$; $|\vec k|-|\vec k|^3/8\le m\le|\vec k|$ (from $m=2\,\mathrm{arcsinh}(\sqrt a/2)$, $y-y^3/6\le\mathrm{arcsinh}\,y\le y$, and $|\vec k|^2-|\vec k|^4/12\le a\le|\vec k|^2$); $m\le\sinh m\le m+m^3$. These give $|g-g_0|\le3|\vec k|^3$ and $|e^{-nm}-e^{-n|\vec k|}|\le\frac n8|\vec k|^3e^{-n|\vec k|/2}$, so
\[
\int_{|\vec k|\le1}|ge^{-nm}-g_0e^{-n|\vec k|}|\,\frac{d\vec k}{(2\pi)^3}\le\frac1{(2\pi)^3}\int_{\R^3}\Big(3|\vec k|^3+\frac n{16}|\vec k|^4\Big)e^{-n|\vec k|/2}d\vec k=\frac{14400}{\pi^2n^6}.
\]
For $|\vec k|>1$ in the cube, $a\ge4|\vec k|^2/\pi^2>4/\pi^2$, so $m\ge m_1>0$ and that part of $G(n)$ is at most $\tfrac12e^{-nm_1}$; the part of the $g_0$-integral over $|\vec k|>1$ is at most $\frac1{6\pi^2}\int_1^\infty r^3e^{-nr}dr$. This gives the bound in (4). Finally $c(n)=(G(n)/G(0))^2=4G(n)^2$, and (4) gives the asymptotics; (3) or (4) gives $-\log c(n)/n\to0$.
\end{proof}

For orientation only: $G(1)=0.0688747$, and $\pi^2n^4G(n)=1.092$ at $n=8$.

\subsection{The Wick inequality}\label{pf:sec-wick}

\begin{lemma}[Wick inequality]\label{pf:lem-wick}
Let $F=\xi+f$ be as in Lemma~\ref{pf:lem-R}, $n\ge1$, and $X_m:=\tfrac12|F_{p_m}|^2$. Then
\begin{align*}
\Cov(X_0,X_n)&=\tfrac32G(n)^2+G(n)\,\E|f_{12}|^2+\tfrac14\Var|f_{12}|^2,\\
\Var X_0&=\tfrac32G(0)^2+G(0)\,\E|f_{12}|^2+\tfrac14\Var|f_{12}|^2,
\end{align*}
so $\Var X_0\ge\tfrac38$ and $\Cov(X_0,X_n)\ge c(n)\Var X_0$.
\end{lemma}

\begin{proof}
$p_0$ and $p_n$ have the same type, so $F_{p_m}=\xi_{p_m}+f_{12}$ with the same $f_{12}$. Condition on $f_{12}$, which is independent of $\xi$. Then $\E[X_m|f_{12}]=\tfrac32G(0)+\tfrac12|f_{12}|^2$, whose covariance at $m=0,n$ is $\tfrac14\Var|f_{12}|^2$. Given $f_{12}$, the conditional covariance of $\tfrac12|\xi_{p_0}|^2+\langle\xi_{p_0},f_{12}\rangle$ and $\tfrac12|\xi_{p_n}|^2+\langle\xi_{p_n},f_{12}\rangle$ has three parts: odd Gaussian moments vanish; $\Cov(\tfrac12|\xi_{p_0}|^2,\tfrac12|\xi_{p_n}|^2)=\tfrac14\cdot3\cdot2G(n)^2$ by Isserlis' theorem and independence of colours; and $\Cov(\langle\xi_{p_0},f_{12}\rangle,\langle\xi_{p_n},f_{12}\rangle\,|\,f_{12})=G(n)|f_{12}|^2$. Taking expectations gives the formula; $n=0$ gives the variance. All moments are finite by Lemma~\ref{pf:lem-R}. By Lemma~\ref{pf:lem-G}(2), $0<G(n)/G(0)\le1$, so each of the three terms of the covariance is at least $c(n)=(G(n)/G(0))^2$ times the corresponding term of the variance; the terms are nonnegative, and summing gives the inequality. $\Var X_0\ge\frac32\cdot\frac14$.
\end{proof}

Constancy of the background is used here: $p_0$ and $p_n$ see the same $f_{12}$.

\subsection{Assembly}\label{pf:sec-assembly}

\begin{proof}[Proof of Theorem~\ref{thm:main}]
Fix $n\ge1$. Suppose that the first claim in \eqref{pf:eq-main} fails at $n$. Then there are $\epsilon>0$, $\beta_k\to\infty$ with $\beta_k\ge1$, and $\mu_k\in\cL_{\beta_k}$ with
$\Cov_{\mu_k}(O_{p_0},O_{p_n})\le(c(n)-\epsilon)\Var_{\mu_k}(O_{p_0})$, where $\Var_{\mu_k}(O_{p_0})>0$ by Proposition~\ref{pf:prop-spectral}. By Proposition~\ref{pf:prop-L4}, along a subsequence the laws converge to some $\mu_\infty$, and by part (d), with $X_m=\tfrac12|F_{p_m}|^2$,
\[
\Cov_{\mu_k}(O_{p_0},O_{p_n})\to\Cov_{\mu_\infty}(X_0,X_n),\qquad\Var_{\mu_k}(O_{p_0})\to\Var_{\mu_\infty}(X_0)\ge\tfrac38 .
\]
Lemmas~\ref{pf:lem-R} and~\ref{pf:lem-wick} apply to $\mu_\infty$ and give $\Cov_{\mu_\infty}(X_0,X_n)\ge c(n)\Var_{\mu_\infty}(X_0)$. The limit of the assumed inequality gives $\Cov_{\mu_\infty}(X_0,X_n)\le(c(n)-\epsilon)\Var_{\mu_\infty}(X_0)$, a contradiction since $\Var_{\mu_\infty}(X_0)>0$. This proves the first claim for every $n\ge1$.

By Lemma~\ref{pf:lem-G}, $G(0)=\tfrac12$, $G(n)>0$, $G(n)=1/(\pi^2n^4)+O(n^{-6})$, and $c(n)>0$ with $-\log c(n)/n\to0$. Proposition~\ref{pf:prop-spectral} gives $\sup_{\mu\in\cL_\beta}m_\mu(\beta)\to0$.
\end{proof}

The bound is through the single vector $v$ of Proposition~\ref{pf:prop-spectral}; this suffices for an upper bound on $m_\mu$. The lower bound $c(n)$ on the ratio decays like $4/(\pi^4n^8)$, and no rate in $\beta$ is obtained.

\subsection{Standard facts used}\label{pf:sec-standard}

The proof uses the following textbook facts without proof.
\begin{enumerate}
\item \emph{Prokhorov's theorem} (tight families in $\R^m$ are relatively compact): existence in Proposition~\ref{pf:prop-L4}.
\item \emph{Diagonal extraction}: existence in Proposition~\ref{pf:prop-L4}, over the countably many finite families of plaquettes.
\item \emph{Portmanteau theorem}, including that bounded Lipschitz functions determine convergence in law, and the continuous-mapping theorem: Corollary~\ref{pf:cor-tail} (open sets); Proposition~\ref{pf:prop-L4}(b), (d), (e).
\item \emph{Uniform integrability and convergence of moments}: Proposition~\ref{pf:prop-L4}(d).
\item \emph{Monotone class ($\pi$--$\lambda$) theorem}: Lemma~\ref{pf:lem-limit}(3) (from continuous local to bounded measurable functions); Proposition~\ref{pf:prop-L4}(e), Step~7.
\item \emph{Slutsky's lemma}: Proposition~\ref{pf:prop-L4}(a) (paths and base points) and (e), Step~6.
\item \emph{Regular conditional distributions} on Polish spaces: Proposition~\ref{pf:prop-L4}(e), Step~7, and Lemma~\ref{pf:lem-D}.
\item \emph{Kolmogorov's extension theorem}: existence of $\mu_\infty$ in Proposition~\ref{pf:prop-L4}.
\item \emph{Scheff\'e's lemma}: the total-variation convergence in Theorem~\ref{pf:thm-L3} is of this type; the proof there bounds the $L^1$ distance of the densities directly.
\item \emph{L\'evy's downward martingale theorem} (a.s.\ and $L^2$): Lemma~\ref{pf:lem-R}.
\end{enumerate}
We also use Weyl's integration formula for $\SU(2)$ (\S\ref{pf:sec-setting}), Plancherel's theorem on $\Z^4$ and a residue computation (Lemma~\ref{pf:lem-G}), and Isserlis' theorem for Gaussian moments (Lemma~\ref{pf:lem-wick}).

% sec_negative.tex -- negative-results sections. \input into a document using preamble.tex.
% Label prefix: neg:. Sources: worker reports and chains listed in report.md of the writer.

\section{Routes to Hypothesis (A) that fail}\label{sec:routes}

Five routes to Hypothesis~(A) were written out as chains of steps. Each step was tagged
proved, cited, claimed or missing, and the first unproved step was tested. This section
states, for each route, what is proved and where the route fails. All programs were written
by workers of one AI model family; no program was run by anyone outside it.

We use the following status words for evidence. \emph{Computed exactly}: exact arithmetic or
a closed form. \emph{Computed in floating point}: a deterministic computation in double
precision. \emph{Monte Carlo by two separate programs}: two codes written independently,
neither copied from the other. \emph{Monte Carlo by one program}: one code, possibly rerun.
\emph{Cited, as stated in its abstract}: the source was read at abstract or search-snippet
level only. Numerical evidence is never proof. Where the sources do not give seeds, we say so.

Two discriminators are used throughout. Compact $\U(1)$ in four dimensions has no gap at weak
coupling \cite{Guth80,FS82} (both as stated in their abstracts). Finite gauge groups in four
dimensions have exponential decay of correlations at weak coupling \cite{AC22}. A route to (A)
that applies verbatim to $\U(1)$ proves something false. A route that cannot reach the finite
groups is incomplete, not wrong.

Throughout, $\beta$ is the Wilson coupling: the weight is
$\exp\bigl(\tfrac{\beta}{2}\sum_p \operatorname{Re}\Tr U_p\bigr)$ against Haar measure on
$\SU(2)^E$. Boxes are $\Lambda$; boundary conditions are $\tau$.

\subsection{Positive curvature}\label{neg:sub-curv}

\emph{Mechanism.} The Wilson measure is $e^{-V}\,d\mathrm{Haar}$ on the Riemannian manifold
$\SU(2)^E$, with $V=-\frac{\beta}{2}\sum_p\operatorname{Re}\Tr Q_p$. If the Bakry--\'Emery tensor
$\mathrm{Ric}+\operatorname{Hess}V$ is bounded below by $K>0$, a Poincar\'e inequality holds with
constant $1/K$, uniformly in $\Lambda$ and $\tau$. Uniform Poincar\'e inequalities give decay of
correlations at a rate independent of the supports, and such decay gives a transfer-matrix gap.

Fix the metric $\langle X,Y\rangle=-\Tr(XY)$ on $\mathfrak{su}(2)$. Then $\mathrm{Ric}=1$
\cite{SZZ23}. Let $B(U)=1+\operatorname{Hess}V(U)$ on $T_U\SU(2)^E$. Let $H_U$ be the orthogonal
complement of the tangent space of the gauge orbit through $U$. Put
$\kappa_H(U)=\min\{B(U)(v,v):v\in H_U,\ |v|=1\}$.

\begin{lemma}\label{neg:lem-horizontal}
Let $f$ be a smooth gauge-invariant function on $\SU(2)^{E(\Lambda)}$. Then $\nabla f(U)\in H_U$
for every $U$, and the Langevin semigroup $P_t$ of $\mu_{\Lambda,\tau}$ maps gauge-invariant
functions to gauge-invariant functions. Consequently, if $\kappa_H(U)\ge K>0$ for every $U$,
then $\Var_{\mu_{\Lambda,\tau}}(f)\le K^{-1}\int|\nabla f|^2\,d\mu_{\Lambda,\tau}$ for every
gauge-invariant $f$, with $K$ independent of $\Lambda$ and $\tau$.
\end{lemma}

\begin{proof}
The gauge group acts by left and right translations in each factor. The metric is
bi-invariant, and $V$ and Haar measure are invariant. So the generator commutes with the
action, and $P_t$ preserves invariant functions. The differential of an invariant $f$
vanishes on orbit directions, so $\nabla f\in H_U$. Bochner's formula gives
$\Gamma_2(f)=\|\nabla^2 f\|^2+B(\nabla f,\nabla f)\ge\kappa_H\Gamma(f)$ along $P_tf$, and the
Bakry--\'Emery argument \cite[Thm.~1.4 and (4.7)]{SZZ23} gives the inequality.
\end{proof}

The gain over the full tensor is not small. On the periodic $3^4$ lattice the minimum of $B$
over all directions was $-6.4$, $-9.0$, $-12.2$ at $\beta=12,24,48$, against $\kappa_H=-1.6$,
$-2.0$, $-1.5$ (Monte Carlo by one program). The lemma was rechecked by a second worker by
argument only.

\begin{proposition}\label{neg:prop-flat}
\textup{(i)} At every flat configuration ($Q_p=1$ for all $p$), $\kappa_H=1$ on the periodic torus.
\textup{(ii)} For one plaquette with $Q=-1$, $B$ has eigenvalues $1-2\beta$ with multiplicity $3$,
on horizontal vectors, and $1$ with multiplicity $9$. Hence a pointwise bound
$\kappa_H\ge K>0$ fails for the one-plaquette system when $\beta\ge 1/2$.
\textup{(iii)} For compact $\U(1)$ the constant in \textup{(i)} is $0$.
\end{proposition}

\begin{proof}
(i) At $Q_p=1$, $\operatorname{Hess}V(t,t)=\frac{\beta}{2}\sum_p|Y_p(t)|^2\ge0$, where $Y_p(t)$
is the signed sum of the transported generators along $p$. So $B\ge1$, and $B=1$ exactly on
horizontal zero modes of the curl form. On the torus these exist (the twelve constant modes).
(ii) At $Q=-1$ the sign of the second-order term reverses:
$\operatorname{Hess}V(t,t)=-\frac{\beta}{2}|Y(t)|^2$. On $\mathfrak{su}(2)^4$ the form $|Y|^2$ has
eigenvalue $4$ with multiplicity $3$ and $0$ with multiplicity $9$. The gauge orbit through a
configuration with $Q=-1$ central has dimension $9$ and lies in the kernel. So the three
eigenvalues $1-2\beta$ are horizontal.
(iii) $\U(1)$ is flat, so $\mathrm{Ric}=0$.
\end{proof}

Part (ii) was computed exactly and confirmed by finite differences in three separate programs
(builder, challenger, and the builder's child). Part (i) is the one place in the route where
non-abelian structure enters. The cited criterion $\kappa_H\ge1-12\beta$ \cite[(4.2), (4.8),
Thm.~1.4]{SZZ23} gives (A) for $\beta<1/12$. On the four-dimensional periodic $3^4$ lattice,
the configuration with all links $1$ except one link $-1$ gives $\kappa_H=1-4.31\beta$ for
$\beta=1,\dots,6$ (computed in floating point by one program).

The route needs a step that carries the curvature to large $\beta$. The step written was:

\begin{hypothesis}[Step S6]\label{neg:hyp-S6}
For large $\beta$, $\kappa_H(U)\ge 1/2$ except on configurations containing plaquettes with
$\operatorname{Re}\Tr Q_p<0$; such configurations have density $e^{-c\beta}$, so a local
perturbation argument applies.
\end{hypothesis}

\begin{observation}\label{neg:obs-curv}
Step S6 is false on every periodic box tested. On the periodic lattices $3^4$ and $4^4$ at
$\beta\in\{4,6,12,24,48\}$, $\kappa_H$ of equilibrium configurations lies in
$[-2.30,-1.12]$ (Table~\ref{neg:tab-curv}). No sample of either program has $\kappa_H\ge1/2$:
$0$ of $32$ (builder) and $0$ of $24$ (challenger). At $\beta=4$ and $6$ the fraction of
plaquettes with $\operatorname{Re}\Tr Q_p<0$ is $0.002$ and $0.000$. Status: Monte Carlo by two
separate programs. The sources do not give the seeds. Thermalization was short ($30+4k$
sweeps in the builder's program; cold start plus $300$ sweeps in the challenger's).
\end{observation}

\begin{table}[ht]
\centering
\small
\begin{tabular}{@{}lllll@{}}
\toprule
program & box & $\beta$ & samples & $\kappa_H$ \\
\midrule
builder & $3^4$ periodic & $4$, $6$ & $8$ each & $-2.2$ to $-1.5$ \\
builder & $3^4$ periodic & $12$, $24$, $48$ & $4$ each & $-2.30$ to $-1.17$ \\
builder & $4^4$ periodic & $4$ & $2$ & $-1.60$, $-1.88$ \\
builder & $4^4$ periodic & $6$ & $2$ & $-1.53$, $-1.58$ \\
challenger & $3^4$ periodic & $4$ & & $-1.79$ to $-1.46$ \\
challenger & $3^4$ periodic & $6$ & & $-1.96$ to $-1.39$ \\
challenger & $3^4$ periodic & $12$ & & $-1.52$ to $-1.12$ \\
challenger & $3^4$ periodic & $24$ & & $-1.93$ to $-1.12$ \\
challenger & $3^4$ periodic & $48$ & & $-2.04$ to $-1.55$ \\
challenger & $4^4$ periodic & $6$ & $2$ & $-1.77$, $-1.26$ \\
challenger & $4^4$ periodic & $24$ & $2$ & $-1.15$, $-1.67$ \\
\midrule
challenger & $4^4$ open & $4$ & & $-2.40$ to $-1.67$ \\
challenger & $4^4$ open & $6$ & $3$ & $+0.16$, $-0.32$, $-0.04$ \\
challenger & $4^4$ open & $12$, $24$, $48$ & & $+3$, $+9.6$, $+23$ \\
\bottomrule
\end{tabular}
\caption{Horizontal curvature $\kappa_H$ of equilibrium configurations, $\SU(2)$. The
challenger ran $20$ periodic samples on $3^4$ in total. The builder's own summary row gives
$[-2.3,-1.2]$ and the director's ledger gives $-1.1$ to $-2.3$; the union of the per-run values
of both programs is $[-2.30,-1.12]$.}
\label{neg:tab-curv}
\end{table}

The challenger found the mechanism, measured only. On the torus the negative directions are
the constant (toron) modes: the lowest eigenvector has zero-momentum weight $0.96$ to $0.99$ in
Landau gauge. On a one-loop ensemble without toron fluctuations, $\kappa_H\to\mathrm{Ric}-k\,c$
with $c=\beta\varepsilon^2$ and $c=1$ at equilibrium; $k=2.62$ and $2.51$ for $\SU(2)$,
constant for $\beta$ from $16$ to $1024$, so the threshold $c_0=\mathrm{Ric}/k$ is $0.38$ to
$0.40$. For $\SU(3)$, $k=3.42$ and $c_0=0.44$ (one seed). In the open $4^4$ box at $\beta\ge12$
the typical curvature is positive (Table~\ref{neg:tab-curv}). This does not rescue the route.
The criterion is pointwise, the open-box positivity grows like $\beta p_{\min}^2$, and every box
of side $L\gtrsim\sqrt\beta$ or with a periodic direction has soft modes; the last statement is
an extrapolation, not a measurement.

\begin{observation}\label{neg:obs-curv-missing}
The route needs, in place of S6, the following statement, which is missing. There are
$\ell,C<\infty$ and $m>0$ such that for every cube $B$ of side at least $\ell$, every pair of
boundary conditions $\tau,\tau'$ and every gauge-invariant $f$ supported in $B$ with
$\|f\|_\infty\le1$,
\[
|\mu_{B,\tau}(f)-\mu_{B,\tau'}(f)|\le C\sum_{e:\,\tau_e\ne\tau'_e}e^{-m\,\mathrm{dist}(\operatorname{supp}f,e)}.
\]
This is strong mixing, essentially the strong mass gap of \cite[Def.~2.3]{Ch21}. It is (A)
in a stronger form, so the route is circular. The two steps after it (strong mixing gives a
uniform log-Sobolev inequality; a uniform Poincar\'e inequality gives decay with a rate
independent of the supports) are claimed from memory, not read.
\end{observation}

\emph{Discriminators.} For $\U(1)$ the route fails at its first step, since $\mathrm{Ric}=0$;
for finite groups there is no Riemannian structure and the route does not apply.

\subsection{Center vortices}\label{neg:sub-vort}

\emph{Mechanism.} A twist of the coupling on a coclosed stack $\mathcal V$ of $L^2$ plaquettes
in the periodic box $(\Z/L\Z)^4$ defines the vortex free energy $F_L(\beta)=-\ln(Z^-/Z)$.
Reflection positivity turns a bound on $F_L$ into a bound on Wilson loops, so decay of $F_L$
like $e^{-\sigma L^2}$ gives an area law. The route then needs an area law to give a gap.

\begin{lemma}\label{neg:lem-Zminus}
For every $\beta>0$ and even $L$, $0<Z^-\le Z$; that is, $F_L\ge0$.
\end{lemma}

\begin{proof}
$Z^->0$ because its integrand is positive. Let $O[\mathcal V]$ be the ratio of twisted to
untwisted weights and $P=\frac12(1-O[\mathcal V])$. By the change of variables of
\cite[(4), (7)]{KT02}, $\langle P\rangle=\langle P\cdot\theta P\rangle$, where $\theta$ is the
reflection in a plane between sites that does not meet $\mathcal V$. Reflection positivity
\cite{OS78} (as stated in its abstract) gives $\langle P\cdot\theta P\rangle\ge0$. So
$\langle O[\mathcal V]\rangle\le1$.
\end{proof}

\begin{proposition}\label{neg:prop-arealaw}
Let $C$ be an $R\times T$ rectangle in the $(1,2)$ plane, with $L/R$ and $L/T$ powers of $2$
and the reflection planes placed as in \cite{KT02} (through sites, not bisecting $\mathcal V$).
Then $\langle W[C]\rangle_{\Lambda_L}\le2\bigl(\tfrac12(1-e^{-F_L})\bigr)^{RT/L^2}$
\cite[(11)]{KT02}. Consequently, if $F_L(\beta)\le KL^2e^{-\sigma L^2}$ for all $L=2^kL_0$, then
every limit point of the periodic states along $L=2^kL_0$ satisfies
$\langle W[C]\rangle\le2e^{-\sigma RT}$.
\end{proposition}

\begin{proof}
The inequality is \cite[(11)]{KT02}, read in full text by the worker; \cite{KT02} attribute it
to \cite{TY85}, which was not read. Since $F_L\ge0$, $\frac12(1-e^{-F_L})\le F_L/2$, so
$\ln\langle W\rangle_{\Lambda_L}\le\ln2+\frac{RT}{L^2}\bigl(\ln(KL^2/2)-\sigma L^2\bigr)\to\ln2-\sigma RT$.
$W$ is local, so the bound passes to weak limits.
\end{proof}

\begin{proposition}\label{neg:prop-noarea}
An area law does not imply a gap. Let $\mu$ be any lattice gauge measure with an area law and
let $\nu$ be the free massless lattice scalar field. Under $\mu\otimes\nu$ every Wilson loop has
the same expectation as under $\mu$, and the transfer matrix has no gap.
\end{proposition}

\begin{proof}
Wilson loops depend only on the gauge factor. The transfer matrix of the product is
$T_\mu\otimes T_\nu$, and the spectrum of its generator contains $0+\operatorname{spec}H_\nu$,
which accumulates at $0$ for the massless field.
\end{proof}

\begin{observation}\label{neg:obs-vort}
The route fails at two independent steps.
\textup{(1)} The decay of $F_L$ at every $\beta$, $\liminf_L(-\ln F_L)/L^2\ge\sigma(\beta)>0$, is
missing; it is the confinement problem. The step meant to approach it, $F_{L'}\le F_L$ for
$L\mid L'$, is claimed by a renormalization-group comparison that is exactly the step
\cite{IS07,IS08} object to in \cite{To07}, and is not closed. It is false for $\Z_2$ at
$\beta=0.6$ (Table~\ref{neg:tab-vort}). For $\SU(2)$ the data neither support nor refute it:
the errors are as large as the differences, and at $\beta=2.4$,
$F_4-F_3=+1.26\pm1.05$, the wrong sign at $1.2$ standard errors.
\textup{(2)} The implication ``area law $\Rightarrow$ gap in the zero-flux sector'' is missing
and false as an abstract implication (Proposition~\ref{neg:prop-noarea}); the theorem of
\cite{Ch21} goes the other way. The flux free energy controls the sectors of nonzero electric
flux, not the neutral sector.
Status: $\SU(2)$ values, Monte Carlo by one program (Metropolis with a step size adapted during
the run, so not strictly detailed balance); $\Z_2$ values, Monte Carlo by one program, rerun bit
for bit by a second worker; the three-dimensional $\Z_2$ check on $L=2$ was computed exactly by
two separate programs ($F(0.3)=0.0325242683$, $F(0.6)=0.6912790067$, agreement to $10^{-10}$).
The sources do not give the seeds.
\end{observation}

\begin{table}[ht]
\centering
\small
\begin{tabular}{@{}llllll@{}}
\toprule
group & $\beta$ & $F_2$ & $F_3$ & $F_4$ & status \\
\midrule
$\SU(2)$ & $1.8$ & $0.69\pm0.24$ & $-0.26\pm0.28$ & $0.49\pm0.40$ & MC, one program \\
$\SU(2)$ & $2.4$ & $3.06\pm0.42$ & $1.55\pm0.45$ & $2.81\pm0.95$ & MC, one program \\
$\Z_2$ & $0.3$ & $0.0774$ & $0.0041$ & $0.0061$ & MC, one program, rerun \\
$\Z_2$ & $0.6$ & $3.205$ & $8.244$ & $15.84$ & MC, one program, rerun \\
\bottomrule
\end{tabular}
\caption{Vortex free energy $F_L$ on $L^4$, $d=4$. $\Z_2$ errors: $0.002$ to $0.005$ at
$\beta=0.3$, $\pm0.25$ at $\beta=0.6$.}
\label{neg:tab-vort}
\end{table}

No proved or cited step of this chain uses non-abelian structure; Lemma~\ref{neg:lem-Zminus}
and Proposition~\ref{neg:prop-arealaw} hold verbatim for $\U(1)$ and $\Z_2$. The non-abelian
input would enter only through asymptotic freedom in the monotonicity step, which is claimed.
An area law at all couplings for a non-abelian group is known only in two dimensions
\cite[p.~8]{Ch18}.

\emph{Discriminators.} For $\U(1)$ the first failing step is the decay of $F_L$ (perimeter law
\cite{Guth80,FS82}, as stated in their abstracts); for $\Z_2$ at weak coupling the monotonicity
step fails while (A) is true there, so the route cannot reach the finite groups.

\subsection{A finite-size criterion}\label{neg:sub-fs}

\emph{Mechanism.} A Dobrushin--Shlosman-type condition $K(L,\beta,\gamma)$ at one scale $L$
gives uniqueness and decay at rate $-\ln\gamma/(2L)$. A finite-volume coupling
$z(L,\beta)=L\,E_f(L,\beta)$, with $E_f$ the lowest energy in one unit of 't~Hooft electric flux
in the periodic box, is claimed to grow with $L$ for $\SU(2)$. If $z$ became large at some scale
and large $z$ implied $K$, the gap would follow.

\begin{lemma}\label{neg:lem-decaygap}
Let $T=e^{-H}$ be the transfer matrix of a reflection-positive state, $0\le T\le1$. Suppose
that for a set of local observables $F$ whose vectors $\psi_F=(F-\langle F\rangle)\Omega$ are
dense in $\Omega^\perp$, $0\le\langle\psi_F,T^n\psi_F\rangle\le K_Fe^{-mn}$ for all $n\ge0$, with
$m$ independent of $F$. Then $\operatorname{spec}(H|_{\Omega^\perp})\subset[m,\infty)$.
\end{lemma}

\begin{proof}
Let $\nu_F$ be the spectral measure of $\psi_F$. For $m'<m$,
$e^{-nm'}\nu_F([0,m'))\le K_Fe^{-mn}$, so $\nu_F([0,m'))=0$. The spectral projection of
$[0,m')$ is bounded and vanishes on a dense set, hence on $\Omega^\perp$.
\end{proof}

The criterion itself (decay from $K$) is claimed: the Dobrushin--Shlosman paper was not read,
and its hypotheses for link variables were not checked.

\begin{observation}\label{neg:obs-fs}
The route fails at two missing steps.
\textup{(1)} For every $\beta>0$, $\sup_Lz(L,\beta)=\infty$. This follows from a positive string
tension at every $\beta$, which is open \cite{Ch21}; it is a confinement statement.
\textup{(2)} There are $z^*<\infty$, $k\in\N$ and $\gamma<1$ such that for every $\beta$ and $L$,
$z(L,\beta)\ge z^*$ implies $K(kL,\beta,\gamma)$. No inequality of any kind is known connecting
a periodic-box, one-sector number to a condition uniform over exterior configurations, even at
strong coupling.
The claimed monotonicity of $z$ rests on the sign of the one-loop coefficient
$b_0=11\cdot2/(48\pi^2)$ and on Table~\ref{neg:tab-fs}. Status: Monte Carlo by one program, one
$\beta$, three sizes, effective mass at $t=1$ (for $\SU(2)$, $m_{\mathrm{eff}}(2)<m_{\mathrm{eff}}(1)$,
so $z$ is overestimated). The sources do not give the seeds. The predicate ``$z_{\U(1)}$ within a
factor $2$ of $1/(2\beta)=0.217$ at $L=4,6$'' failed; that was the worker's expectation, which
ignored the finite temporal extent. A rotor model fitted afterwards reproduces the drift with
residuals of $1.5$ to $2.6$ standard errors.
\end{observation}

\begin{table}[ht]
\centering
\small
\begin{tabular}{@{}llll@{}}
\toprule
$L$ & $z$, $\SU(2)$ & $z$, $\U(1)$ & rotor model at $T=12$ (post hoc) \\
\midrule
$2$ & $0.562\pm0.016$ & $0.335\pm0.024$ & $0.324$ \\
$4$ & $2.12\pm0.06$ & $0.553\pm0.078$ & $0.422$ \\
$6$ & $8.02\pm0.40$ & $0.684\pm0.068$ & $0.500$ \\
\bottomrule
\end{tabular}
\caption{Flux coupling $z=L\,m_{\mathrm{eff}}(1)$ on $L^3\times12$ at $\beta=2.3$, from the
correlator of the zero-momentum spatial Polyakov loop. Free-photon value $1/(2\beta)=0.217$.}
\label{neg:tab-fs}
\end{table}

The table shows a steep rise for $\SU(2)$ and a slow drift for $\U(1)$ that is probably thermal.
It shows confinement-scale physics at $\beta=2.3$, not running from weak coupling, which would
need $L$ of order $\exp(\mathrm{const}\cdot\beta)$.

\emph{Discriminators.} For $\U(1)$ the route stops at the growth of $z$, as it must; for finite
groups $z\to0$ in the deconfined phase while (A) holds, so the route cannot reach them.

\subsection{Positivity of the character expansion}\label{neg:sub-dual}

\emph{Mechanism.} Expand each plaquette weight in characters and integrate out the links. If
every resulting amplitude were nonnegative, $\SU(2)$ would be a positive random-surface model.
Percolation-type sharpness arguments could then give the gap without any appeal to a running
coupling.

\begin{lemma}\label{neg:lem-cj}
For $\beta>0$ and $j\in\frac12\N_0$,
$e^{\frac\beta2\Tr U}=\sum_jc_j(\beta)\chi_j(U)$ with
$c_j(\beta)=I_{2j}(\beta)-I_{2j+2}(\beta)=\frac{2(2j+1)}{\beta}I_{2j+1}(\beta)>0$.
\end{lemma}

\begin{proof}
For $U$ with eigenvalues $e^{\pm i\theta}$, $\frac12\Tr U=\cos\theta$ and
$\chi_j=\sin((2j+1)\theta)/\sin\theta$. From $e^{\beta\cos\theta}=\sum_{n\in\Z}I_n(\beta)e^{in\theta}$,
the coefficient of $\sin((2j+1)\theta)$ in $e^{\beta\cos\theta}\sin\theta$ is $I_{2j}-I_{2j+2}$.
The recurrence $I_{\nu-1}-I_{\nu+1}=(2\nu/\beta)I_\nu$ gives the second form, and $I_n(\beta)>0$
because its power series has positive terms.
\end{proof}

The identity was also checked in floating point for $2j\le10$ and five values of $\beta$
(relative difference below $10^{-15}$). For a configuration $\{j_p\}$ of plaquette spins put
$A(\{j_p\})=\int\prod_p\chi_{j_p}(U_p)\prod_e dU_e$. Then
$Z=\sum_{\{j_p\}}\prod_pc_{j_p}(\beta)\,A(\{j_p\})$. For an abelian group $\prod_p\chi_{j_p}(U_p)$
is a character of the group of link variables, so $A\in\{0,1\}$.

\begin{observation}\label{neg:obs-dual}
Not every $\SU(2)$ amplitude is nonnegative. On the periodic $2\times1\times1$ lattice in three
dimensions, $4$ of the $729$ spin configurations considered have $A<0$, $141$ have $A>0$ and
$584$ have $A=0$. The minimum is $A=-0.003086419753$, at spins $(j_{01},j_{02},j_{12})=(0,1,1)$ on
site $0$ and $(1,0,\frac12)$ on site $1$. A separate Monte Carlo program gives
$-0.002892\pm0.000237$ (seed $7$, $4\cdot10^7$ samples), $12$ standard errors below zero.
Setting every plaquette that touches a fourth direction to spin $0$ embeds the configuration in
four dimensions with the same amplitude (proved by the worker: those links occur in no character).
Status: computed in floating point by one program (intertwiner contraction), confirmed by a
second program by a different method. The value agrees with $-1/324$ to the digits printed;
the rational value is not proved. On the one-site three-torus all $27$ amplitudes are positive,
so the smallest lattice does not show the sign.
\end{observation}

The literature knows this as the sign problem of the dual \cite{CCK07}; that source was read in
search snippets only. The test refutes positivity of the plaquette-character dual only. It does
not exclude a positive representation obtained by resummation over intertwiners, or any other
representation.

\emph{Discriminators.} For $\U(1)$ and $\Z_2$ the amplitudes are indicators and positivity holds;
since $\U(1)$ in four dimensions has no gap at weak coupling, positivity alone could never have
carried the route, and for $\Z_2$ the surface model has a transition at large $\beta$ while (A)
holds.

\subsection{Exclusion}\label{neg:sub-excl}

\emph{Mechanism.} $\SU(2)$ has no continuous global or one-form symmetry, only its center
$\Z_2$. If ``no continuous symmetry to break'' implied a gap except at isolated couplings, (A)
would reduce to excluding a transition at weak coupling. That exclusion would follow if the
zero set of $m$ were closed, the correlation length diverged at its infimum, and the free energy
were non-analytic there.

\begin{lemma}\label{neg:lem-usc}
Assume \textup{(H1)} a unique infinite-volume Gibbs state $\mu_\beta$ for every $\beta$;
\textup{(H2)} the specification is continuous in $\beta$ and in the boundary condition;
\textup{(H3)} reflection positivity, so that $C_F(n)=\langle\psi_F,T^n\psi_F\rangle$;
\textup{(H4)} vectors of local positive-time observables are dense. Then
$m(\beta)=\inf_F\inf_k\bigl(-\tfrac1{2k}\log(C_F(2k)/C_F(0))\bigr)$ is upper semicontinuous in
$\beta$. In particular $\{\beta:m(\beta)\ge c\}$ is closed for every $c$.
\end{lemma}

\begin{proof}
(H1), (H2) and compactness give weak continuity of $\beta\mapsto\mu_\beta$, so each $C_F(n)$ is
continuous. By the Lyapunov inequality, $-\frac1{2k}\log(C_F(2k)/C_F(0))$ is nonincreasing in $k$
with limit $-\log\|T|_{\overline{\mathrm{span}}}\|$. By (H4) the infimum over $F$ is $m$. An
infimum of continuous functions is upper semicontinuous.
\end{proof}

(H1) is the standing assumption of the paper and is not known at large $\beta$; (H2) is claimed
as standard, with no source read.

\begin{observation}\label{neg:obs-usc}
Under \textup{(H1)--(H4)} nothing beyond upper semicontinuity follows from continuity of the
correlations. An abstract continuous family of spectral measures has zero set
$\{m=0\}=(1,\infty)$, which is not closed, and a correlation length equal to $1$ up to the
transition and infinite after it (checked by one program). It is not a lattice model. No theorem
of the form ``the gap vanishes $\Rightarrow$ the free energy is non-analytic'' was found.
\end{observation}

The principle itself is false; the counterexamples are Observation~\ref{neg:obs-excl}.

\emph{Discriminators.} For $\U(1)$ the route is consistent but uninformative (it has a continuous
one-form symmetry and no gap); the finite groups $\Z_N$ with large $N$ are the counterexample to
the principle.

\section{A shared-properties observation}\label{sec:shared}

\begin{observation}\label{neg:obs-shared}
The following six properties hold for $\SU(2)$, $\U(1)$ and $\Z_2$ Wilson lattice gauge theory
in four dimensions.
\begin{enumerate}
\item[P1] Reflection positivity at every $\beta$. Cited \cite{OS78}, as stated in its abstract.
\item[P2] A strictly positive self-adjoint transfer matrix with a Perron--Frobenius ground
state. Positivity of the transfer matrix is cited \cite{OS78,Lu77}, as stated in their
abstracts; the Perron--Frobenius ground state is argued in the worker report from the
Hamiltonian limit, not cited.
\item[P3] Positive character coefficients of the Wilson weight at every $\beta>0$. Proved here
for $\SU(2)$ (Lemma~\ref{neg:lem-cj}); for $\U(1)$ the coefficients are $I_n(\beta)>0$; for
$\Z_2$ they are $\cosh\beta$ and $\sinh\beta$.
\item[P4] A gap and analyticity at small $\beta$. Cited \cite{OS78} for the gap and area law at
strong coupling for every compact group, as stated in its abstract; analyticity by cluster
expansion is claimed as standard.
\item[P5] A positive gap in every finite periodic volume at every $\beta$. Stated in the worker
reports, not written out there; it follows from P2 by the Perron--Frobenius theorem for the
compact finite-volume transfer operator.
\item[P6] A nontrivial center: $\Z_2\subset\SU(2)$; $\U(1)$ and $\Z_2$ are their own centers.
Elementary.
\end{enumerate}
Compact $\U(1)$ in four dimensions has P1--P6 and no gap at weak coupling \cite{Guth80,FS82}
(as stated in their abstracts). Consequently no argument that uses only P1--P6 proves (A).
\end{observation}

The observation does not exclude an argument that uses a property outside the list. It does
not exclude an argument that uses the dimension in a way that fails for $\U(1)$ in four
dimensions, and it does not exclude an argument that uses a property $\U(1)$ lacks. Three such
properties are known with proof or computation:
\begin{enumerate}
\item $\SU(2)$ is perfect: $[\SU(2),\SU(2)]=\SU(2)$, so it has no nontrivial homomorphism to
$\U(1)$. Elementary.
\item Its Ricci curvature is the positive constant $1$, against $0$ for $\U(1)$
(Proposition~\ref{neg:prop-flat}). Proved.
\item Its plaquette-character amplitudes can be negative, while those of $\U(1)$ and $\Z_2$ are
indicators (Observation~\ref{neg:obs-dual}). Computed in floating point and confirmed by Monte
Carlo; not proved in exact arithmetic.
\end{enumerate}
All three are local. None has a known global consequence for the gap. The official problem
statement records the same absence: there is no known idea for a gap uniform in the volume
\cite{JW}. In the routes of
Section~\ref{sec:routes} the second enters only at flat configurations and is outweighed on
typical ones (Observation~\ref{neg:obs-curv}); the third enters as cancellation, which removes
the positivity tools. Three other candidates meet the same wall. Analyticity in $\beta$ is
proved only at small $\beta$, and at all $\beta$ it would exclude a transition that $\U(1)$ has.
A uniform gap of the Langevin generator is the strong mixing of
Observation~\ref{neg:obs-curv-missing}. No inequality monotone in the dimension is known, and
$\U(1)$ has a gap in three dimensions \cite{GM82} (as stated in its abstract) and none in four.

\begin{observation}\label{neg:obs-excl}
The principle ``a lattice gauge theory with no continuous global or one-form symmetry has a
gap except at isolated couplings'' fails.
\begin{enumerate}
\item Four-dimensional $\Z_N$ lattice gauge theory with $N$ large has an intermediate massless
phase, on an interval of couplings, in which local $\U(1)$ invariance is restored dynamically
\cite{FS82}. Status: cited, as stated in its abstract. The action, the exact form of the
massless statement and the interval were not read. $\Z_N$ has only a $\Z_N$ one-form symmetry.
\item Five-dimensional $\SU(2)$ Wilson theory has a first-order bulk transition at
$\beta_c\approx1.64$ into a phase with zero string tension \cite[\S3 and Fig.~1]{FiveD11},
credited there to \cite{Cr79} (as stated in its abstract). Status: numerical. That the
weak-coupling phase is massless is perturbative expectation; it was not found stated in the
text read, and it is not proved. A Monte Carlo by one program on $4^5$ (200 Metropolis sweeps,
one seed per start, written by a sub-worker and rerun bit-identical by the worker) gives a
hot--cold hysteresis gap of $0.090\pm0.003$ at $\beta=1.6$, $0.018$ at $\beta=1.7$ (still
drifting), and below $0.0015$ at $\beta=1.3$, $1.5$, $1.9$ and $2.2$ (the same report's summary
table says below $0.002$). Its weak-coupling plaquette approaches the one-loop value
$1-3/(5\beta)$: $0.6925$ against $0.7273$ at $\beta=2.2$ and $0.630$ against $0.684$ at
$\beta=1.9$. This is consistent with a perturbative phase and is not a measurement of a massless
correlator.
\end{enumerate}
$\SU(2)$ in five dimensions is compact, connected, simple and perfect, has positive Ricci
curvature and signed amplitudes, and its reflection positivity in site planes does not depend on
the dimension (claimed, re-derived for four dimensions only). So an exclusion argument for (A)
must use $d=4$ in a way that fails at $d=5$, and must forbid an emergent continuous one-form
symmetry, as in the $\Z_N$ example. Both are dynamical statements about the running of the
coupling.
\end{observation}

\section{Hypotheses (C) and (E)}\label{sec:CE}

\subsection{Hypothesis (C)}

Let $T=e^{-H}$ be the transfer matrix at $\beta$, $\psi\perp\Omega$ the vector of a local
observable, $\nu$ its spectral measure, $C(n)=\langle\psi,T^n\psi\rangle=\int e^{-n\varepsilon}\,d\nu(\varepsilon)$
for real $n\ge0$, $m=m(\beta)$ and $\xi=1/m$. We write (C) in the form
$C(\xi/2)\le D\,C(\xi)$ with $D$ independent of $\beta\ge\beta_1$. The effective mass
$u(n)=-\frac{d}{dn}\log C(n)$ is the mean energy of $e^{-nH/2}\psi$; (C) is equivalent, up to
constants and a change of time from $\xi/2$ to $\xi$, to $u(\xi/2)\le K\,m$.

\begin{proposition}\label{neg:prop-Cprime}
Assume \textup{(C$'$)}: $1-C_\beta(1)/C_\beta(0)\le K'm(\beta)$ for all $\beta\ge\beta_1$. Then for
every $j\ge0$ with $\xi/2^j\ge1$ and $\xi\ge2K'$,
$C(\xi/2^j)\le e^{2K'}C(\xi)$. In particular (C) holds with $D=e^{2K'}$.
\end{proposition}

\begin{proof}
$(\log C)''$ is the variance of $\varepsilon$ under $e^{-n\varepsilon}d\nu/C(n)$, so $\log C$ is
convex. Put $q=C(1)/C(0)$. For $a\ge1$ the slope of $-\log C$ on $[a,\infty)$ is at most its
chord slope on $[0,1]$, which is $\log(1/q)$. Hence
$\log C(\xi/2^j)-\log C(\xi)\le\xi\log(1/q)\le2\xi K'm=2K'$, using $-\log(1-x)\le2x$ for
$x\le\frac12$.
\end{proof}

The worker compiled the integer form of this chain (Lean, declaration \texttt{A3C.step\_chain})
with exit code $0$ and the axioms propext, Classical.choice, Quot.sound; it was not recompiled
for this section. The real-time convexity argument is not formalized.

\begin{observation}\label{neg:obs-Cprime-free}
(C$'$) fails for lattice-local observables in the free massive lattice field, while (C) holds
there. For $\varphi(0)^2$, $\xi(1-C(1)/C(0))=0.963\,\xi$ for $\xi$ from $4$ to $256$, and $C(1)/C(0)$
tends to the lattice constant $0.0373$. For $\varphi(0)$, $R(\xi)=C(\xi/2)/C(\xi)$ is $6.584$,
$5.921$, $5.580$, $5.522$, $5.508$, $5.505$, $5.504$ at $\xi=4,8,\dots,256$, against the continuum
value $2K_1(1/2)/K_1(1)=5.503975$; for $\varphi(0)^2$ it tends to $30.2937$. Status: computed in
floating point by one program (Schwinger representation); the values of $R$ for $\varphi(0)$ were
recomputed by a second program (momentum integral), agreeing to $10^{-6}$ relative. So (C$'$)
can be asked only of an observable smeared on the physical scale. For the slice-flowed free
field at $s=0.1$, $\xi(1-C(1)/C(0))=3.56$, $4.13$, $4.46$, $4.64$, $4.73$, $4.78$ at
$\xi=8,16,32,64,128,256$ (one program): bounded, as (C$'$) requires.
\end{observation}

What (C) needs for $\SU(2)$ is therefore missing: (C$'$) for the slice-flowed magnetic density at
flow time $s\xi(\beta)^2$, in infinite volume, uniformly in $\beta$.

\begin{proposition}\label{neg:prop-smear}
Smearing in Euclidean time by a fixed number of correlation lengths does not make (C)
automatic. Let $k\ge1$, $\psi=T^{k\xi/2}\varphi$, and let $\varphi$ have spectral measure
$\nu_\varphi=w\,\delta_m+\delta_1$ with $m=1/\xi<1$ and $w=e^{-(k+1)(\xi-1)}$. Then
$C_\psi(\xi/2)/C_\psi(\xi)\ge\frac12e^{\xi/2}$.
\end{proposition}

\begin{proof}
$C_\psi(n)=C_\varphi(n+k\xi)$ and $C_\varphi(n)=we^{-n/\xi}+e^{-n}$. So
$C_\varphi((k+1)\xi)=2e^{-(k+1)\xi}$ and $C_\varphi((k+\frac12)\xi)\ge e^{-(k+\frac12)\xi}$.
\end{proof}

Smearing by a $\beta$-dependent time $n_\beta$ chosen so that $u_\varphi(n_\beta)\le(1+\delta)m$
does give (C), trivially. If $n_\beta/\xi$ is unbounded, the vector is no longer a sharp-time
local field; if it is bounded, that boundedness is already (C)-type information. By explicit
examples, (C) fails when the spectral weight at the bottom is exponentially small in $\xi$, a
hole in the spectral measure alone does not break it, and (C) does not exclude a limit that is a
single atom.

\begin{observation}\label{neg:obs-matched}
In a three-spacing $\SU(2)$ measurement of the slice-flowed magnetic density
(Table~\ref{neg:tab-matched}), $\xi(1-C(1)/C(0))$ is $0.53$, $0.68$, $0.86$ at $\xi=0.72$, $1.01$,
$1.40$. It increases, roughly like $\xi^{0.75}$, over the factor $1.9$ available. Status: Monte
Carlo by one set of programs, measured only. Caveats: the correlation lengths are about one lattice
step, the worst place to ask the question; $\xi$ is taken from a published glueball mass, not
measured; each of the two added seeds at $N=12$ repeats three configurations (IDs $40$--$42$)
after a restart, found by another director's audit, and the table gives the original
row-based estimates (according to the ledger, with the repeats removed the estimates move
within their errors); the boxes are matched in physical size only to $3$ to $7$ per cent.
\end{observation}

\begin{table}[ht]
\centering
\small
\begin{tabular}{@{}lllllllll@{}}
\toprule
$N^4$ & $\beta$ & $s/a^2$ & configs & $a\sqrt\sigma$ & $R(0)\cdot10^2$ & $R(1)\cdot10^2$ & $\xi$ & $\xi(1-q)$ \\
\midrule
$6^4$ & $2.3$ & $0.5$ & $2990$ & $0.3690(30)$ & $3.751(42)$ & $1.012(32)$ & $0.72$ & $0.53$ \\
$8^4$ & $2.4$ & $0.889$ & $1500$ & $0.2660(20)$ & $3.615(60)$ & $1.166(51)$ & $1.01$ & $0.68$ \\
$12^4$ & $2.5$ & $2.0$ & $840$ ($834$) & $0.1905(8)$ & $3.233(61)$ & $1.245(51)$ & $1.40$ & $0.86$ \\
\bottomrule
\end{tabular}
\caption{Slice-flowed magnetic density, flow to $s=\ell^2/72$ with $\ell$ the box side.
$R(\tau)$ is the connected correlator of the slice average over $\langle\bar E\rangle^2$, so
$q=R(1)/R(0)=C(1)/C(0)$: $0.270$, $0.323$, $0.385$. Scale $a\sqrt\sigma$ from \cite{FHK93}
(Table~3); $\xi=1/(3.74\,a\sqrt\sigma)$ with the scalar glueball mass $3.74(12)\sqrt\sigma$ of
\cite{Te98} (\cite{LTW04} gives $3.78(7)$). At $N=12$: three seeds, $840$ rows, $834$ distinct
configurations. The last three columns are derived; errors are not propagated because the
correlation of $R(0)$ and $R(1)$ is not given. The sources do not give the production seeds.}
\label{neg:tab-matched}
\end{table}

The measurement supports one sentence: on these lattices $\xi(1-C(1)/C(0))$ grows with $\xi$,
which is the direction that (C$'$) for the slice-flowed density does not predict. It refutes
neither (C$'$) nor (C).

\subsection{Hypothesis (E)}

(E) asks that the scaling limit be interacting and asymptotically free at short distance. A
generalized free field (GFF) is a Gaussian field with an arbitrary K\"all\'en--Lehmann measure;
a free algebra is the Wick-polynomial algebra of finitely many GFFs.

\begin{proposition}\label{neg:prop-GFF}
\textup{(i)} Every Euclidean-invariant, reflection-positive two-point function is the two-point
function of a GFF. Hence no statement about a single two-point function excludes a GFF.
\textup{(ii)} Let $O=\frac12\sum_{a=1}^k\varphi_a^2$ with $\varphi_a$ independent copies of one
GFF. Then the normalized fourth cumulant of $O$ at one point is $\kappa_4/\kappa_2^2=12/k$, which
is $4$ for $k=3$. Hence a nonvanishing truncated four-point function of a composite field does
not exclude a free algebra.
\end{proposition}

\begin{proof}
(i) Take the Gaussian measure with that covariance; a Gaussian measure is reflection positive
when its covariance is. (ii) $2O$ is $\sigma^2\chi^2_k$, whose cumulants are
$\kappa_n=2^{n-1}(n-1)!\,k$ in units of $\sigma^{2n}$; so $\kappa_2=2k$, $\kappa_4=48k$.
\end{proof}

Part (ii) applies to the rescaled plaquette energy used in the program's proof of (B),
Theorem~\ref{thm:main}. At fixed lattice distance and
$\beta\to\infty$ the limit of $\beta(1-\cos\alpha_p)$ is $\frac12\sum_{a=1}^3(F^a_p)^2$ (a claimed
step); if the three colour components are independent with equal variance (not checked), it
has $\kappa_4/\kappa_2^2=4$. A test ``$\kappa_4\ne0\Rightarrow$ interacting'' applied to it would
declare the free Maxwell limit interacting. A second mass, or any mass ratio, does not
distinguish either: a GFF with measure $\delta_{m_1^2}+\delta_{m_2^2}$ has any ratio.

A formulation that does exclude every free algebra is the two-particle energy shift. Assume
that the lightest mass $m$ is an isolated eigenvalue of the mass operator. Let $E_2(L)$ be the
lowest two-particle energy at zero total momentum in a periodic box of side $L=z/m$, and
$\delta(z)=\lim_{\beta\to\infty}(E_2(L)-2m)/m$. Then (E$^*$): $\delta(z)\ne0$ for some $z$ in some
symmetry channel. In the Fock space of a free algebra the Hamiltonian is a sum of one-particle
Hamiltonians, so $E_2(L)=2m$ exactly; hence (E$^*$) excludes every free algebra. It needs the
isolated one-particle shell, which a gap does not give. Its relation to the scattering length
through \cite{LW86} is claimed from memory; the source was not read.

\begin{table}[ht]
\centering
\small
\begin{tabular}{@{}lll@{}}
\toprule
model & quantity & values \\
\midrule
free field, $d=4$, $L=8$, $m=0.25$ & $U_4$; $g_R$; block $\kappa_4/\kappa_2^2$ & $0.007\pm0.025$; $0.34\pm1.20$; $-0.19\pm0.07$ \\
Ising, $d=4$, $\xi_2/L=0.25$ & $g_R$ at $L=6,8,12$ & $67\pm9$, $46\pm8$, $32\pm7$ \\
Ising, $d=4$, $\xi_2/L\approx0.33$ & $g_R$ at $L=6,8,12,16$ & $36\pm2$, $\approx29$, $22\pm2$, $20\pm3$ \\
Ising, $d=2$, $K_c$ & $U_4$ at $L=16,32$ & $0.6111\pm0.0005$, $0.6098\pm0.0007$ \\
\bottomrule
\end{tabular}
\caption{Calibration on known answers. Free field: $4000$ exact Gaussian samples, $\xi_2=4.023\pm0.046$
against the exact $4$. Four-dimensional Ising: Wolff cluster algorithm, $17$ runs, seeds
$2001$--$2017$, $15000$ measurements each. Two-dimensional Ising: seeds $31$, $32$; a separate
program (the worker's child) gives $0.6115$, $0.6107$. Monte Carlo, measured only. The
renormalized coupling is $g_R=3U_4(L/\xi_2)^4$ with $U_4=1-\langle M^4\rangle/(3\langle M^2\rangle^2)$.}
\label{neg:tab-ising}
\end{table}

\begin{observation}\label{neg:obs-ising}
On the calibration of Table~\ref{neg:tab-ising}, the free field gives $U_4$ and $g_R$ compatible
with zero (the block cumulant sits $3$ standard errors from zero and was not rerun with a second
seed). In four-dimensional Ising, $g_R$ at fixed $\xi_2/L$ decreases with $L$, consistent with
the Gaussian scaling limit of \cite{ADC21} (as stated in its abstract); a logarithmic law was not
tested, and the predicate ``$g_R(12)$ exceeds zero by five standard errors at $\xi_2/L=0.25$''
failed at $4.6$. In two-dimensional Ising at criticality $U_4\approx0.61$, far from zero. The
critical $U_4$ on a periodic torus is nonzero also for models with Gaussian limits, because of
the zero mode, so only $g_R$ at fixed $L/\xi_2$ calibrates the test.
\end{observation}

What is not known. No rigorous non-Gaussian scaling limit of a lattice gauge theory in three or
more dimensions is known to the workers; no search establishing absence was made. The constant
in the short-distance statement of (E) for the trace of the stress tensor is not derived. Whether
(E$^*$) can be proved from lattice information is open: it lives at the correlation length,
where control is the gap problem itself.

\section{Provenance and formal checks}\label{sec:provenance}

\subsection{Who did what}
The author directed the work, set its aim (an attempt at the whole problem, written to be disproved), and is accountable for the paper. He did not write or check the proofs.

The statements, proofs, programs and text were produced by sessions of one model family, Claude (Anthropic). One session acted as director. It wrote the reduction of Section~\ref{sec:reduction} in its first form, chose the routes of Section~\ref{sec:routes}, and assembled the paper. Subordinate sessions, started by the director as separate processes and each with subordinate sessions of its own, did the following: \emph{builders} wrote the proofs of Theorems~\ref{thm:main}, \ref{thm:reduction} and~\ref{thm:conditional} and the arguments of Sections~\ref{sec:routes}--\ref{sec:CE}; \emph{challengers}, given only the text under attack, tried to break each of them with their own programs; \emph{formalizers} wrote the Lean files of Section~\ref{sec:lean}; \emph{readers} located statements in the literature. Every session shares its model family with the director. Their agreement reduces the exposure to error. It is not independent review.

\subsection{What was attacked, and what the attacks found}
Theorem~\ref{thm:reduction}: one challenge. The proof stood. Eight statements around it were wrong and were corrected, among them the claim that Hypothesis~\ref{hyp:C} concerns locality alone.

Theorem~\ref{thm:main}: three challenges, of successive versions. The first found a circular gauge fixing, which was repaired, and could not check a classification of Gibbs states, which the second version avoids. The second found a mislabelled test and a quantifier missing from one statement. The third, of the two steps written last (Propositions on the Laplace limit and the limit field in Section~\ref{sec:proof}), found no failing line.

Theorem~\ref{thm:conditional}: one challenge. The statement stood. One step of the proof used positivity of the transfer matrix, which the hypotheses do not give; the proof in Section~\ref{sec:reduction} uses the repaired step.

Sections~\ref{sec:routes}--\ref{sec:CE}: the refutation in the curvature route was reproduced by a second session with separate code. The other numerical statements rest on one program each, except where the text says otherwise. A review of the working summaries against the underlying reports found six wrong and ten overstated sentences, which were corrected before this paper was written. The paper itself has had the corresponding review (Section~\ref{sec:paperreview}).

\subsection{Checks by a second model lineage}
Four other working groups of the same program are sessions of OpenAI Codex. Two of them checked Lemma~\ref{pf:lem-R} (the decomposition of a limit point into a free field and an independent constant background) as it stood at commit \repo{34f97d2} of~\cite{Repo}, each with its own argument for the constancy of the background, and one of them derived the asymptotics $G(n)=\pi^{-2}n^{-4}+O(n^{-6})$ and recomputed the constant in the tail bound of Section~\ref{sec:proof}. Their stated scope is that lemma, the Wick inequality that follows it, and the free correlation; they did not check the identification of the limit of the Wilson measures (the Laplace limit and the limit field). These checks are recorded in~\cite{RepoThread}.

\subsection{What has not been done}
No mathematician has read the paper. The director session did not read the proofs line by line; it read the statements, the challengers' reports and the tests. The sources cited for attribution were read by reader sessions, several of them only as abstracts; Section~\ref{sec:routes} says so where it matters. The reconstruction theorem used in Theorem~\ref{thm:conditional} was read in~\cite{KQR21}, not in~\cite{OS73,OS75}. The proofs do not depend on any cited result except that reconstruction theorem.

\subsection{Formal checks}\label{sec:lean}
Lean 4.19.0 with Mathlib at tag \repo{v4.19.0} (commit \repo{c44e0c8}). Each file compiles with exit code $0$, contains no \repo{sorry}, and \repo{\#print axioms} reports only \repo{propext}, \repo{Classical.choice}, \repo{Quot.sound} for every declaration listed. Paths are relative to \repo{yang-mills/z4/attempt/reviews} in~\cite{Repo}. The label \emph{formally checked} below means exactly this; it is not a statement about the analytic steps, none of which is formalized.

\begin{center}
\begin{tabular}{@{}p{0.36\textwidth}p{0.27\textwidth}p{0.29\textwidth}@{}}
\toprule
Printed statement & File & Declarations\\
\midrule
Finite-sum inequalities in the proof of Theorem~\ref{thm:reduction} (tail bound, gap bound, Jensen bound, mean energy, sharp constant $e$) & \repo{a2-leanR/A1R.lean} & \repo{A1R.L1}--\repo{L7}, \repo{L6\_sharp'}, \repo{D\_ge}, \repo{mutation\_false}\\
Reduction step in Section~\ref{sec:proof}: a lower bound on a normalized correlation bounds the gap & \repo{a3-B/S2.lean} & \repo{s2\_gap\_le}, \repo{s2\_power}\\
Scalar inequalities used in the Laplace limit: $1-\cos x\ge 2x^2/\pi^2$ on $[-\pi,\pi]$, and three quadratic and exponential bounds & \repo{a8-BL34/L34.lean} & \repo{one\_sub\_cos\_ge}, \repo{half\_pos\_part\_sq}, \repo{norm\_add\_sq\_ge}, \repo{exp\_rem}\\
\bottomrule
\end{tabular}
\end{center}

The formal statements are for finite sums of atoms and for real numbers. Where the text uses a measure in place of a finite sum, the printed statement is more general than the formal one and is proved in the text.

\subsection{Records}
The working record of the attempt, including every brief given to a subordinate session, every program, result file and report, is in~\cite{Repo}. Each session's runs are recorded in an append-only hash-chained register (PCP OS; the Rust host built from commit \repo{92623fa} of its source verifies each chain), and the hierarchy of sessions for the second half of the work is recorded in that host's task graph. These records show what was run and in what order. They say nothing about whether a result is right.

\subsection{Review of this paper}\label{sec:paperreview}
(to be filled after the review of the assembled paper)



\begin{thebibliography}{99}
% refs.tex -- bibliography entries, p4-refs (task T22). Every field comes from a record
% fetched on 23 November 2026; see report.md for the record of each field.
% Order: keys sorted case-insensitively (the four Repo* entries are merged into the order).

\bibitem{AC22}
A.~Adhikari and S.~Cao, \emph{Correlation decay for finite lattice gauge theories at weak coupling},
Ann. Probab. \textbf{53} (2025), no.~1, 140--174.
\url{https://doi.org/10.1214/24-AOP1702}; arXiv:2202.10375.

\bibitem{ADC21}
M.~Aizenman and H.~Duminil-Copin, \emph{Marginal triviality of the scaling limits of critical 4D Ising and $\phi_4^4$ models},
Ann. of Math. \textbf{194} (2021), no.~1, 163--235.
\url{https://doi.org/10.4007/annals.2021.194.1.3}; arXiv:1912.07973.

\bibitem{CCK07}
J.~W.~Cherrington, J.~D.~Christensen, and I.~Khavkine, \emph{Dual computations of non-Abelian Yang--Mills theories on the lattice},
Phys. Rev. D \textbf{76} (2007), 094503.
\url{https://doi.org/10.1103/PhysRevD.76.094503}; arXiv:0705.2629.

\bibitem{Ch18}
S.~Chatterjee, \emph{Yang--Mills for probabilists}, preprint (2018), arXiv:1803.01950.

\bibitem{Ch21}
S.~Chatterjee, \emph{A probabilistic mechanism for quark confinement},
Comm. Math. Phys. \textbf{385} (2021), no.~2, 1007--1039.
\url{https://doi.org/10.1007/s00220-021-04086-y}; arXiv:2006.16229.

\bibitem{Cr79}
M.~Creutz, \emph{Confinement and the critical dimensionality of space-time},
Phys. Rev. Lett. \textbf{43} (1979), 553--556; Erratum, Phys. Rev. Lett. \textbf{43} (1979), 890.
\url{https://doi.org/10.1103/PhysRevLett.43.553}.

\bibitem{FHK93}
J.~Fingberg, U.~M.~Heller, and F.~Karsch, \emph{Scaling and asymptotic scaling in the SU(2) gauge theory},
Nuclear Phys. B \textbf{392} (1993), 493--517.
\url{https://doi.org/10.1016/0550-3213(93)90682-F}; arXiv:hep-lat/9208012.

\bibitem{FILS78}
J.~Fr\"ohlich, R.~Israel, E.~H.~Lieb, and B.~Simon, \emph{Phase transitions and reflection positivity. I. General theory and long range lattice models},
Comm. Math. Phys. \textbf{62} (1978), 1--34.
\url{https://doi.org/10.1007/BF01940327}.

\bibitem{FiveD11}
F.~Knechtli, M.~Luz, and A.~Rago, \emph{On the phase structure of five-dimensional SU(2) gauge theories with anisotropic couplings},
Nuclear Phys. B \textbf{856} (2012), 74--94.
\url{https://doi.org/10.1016/j.nuclphysb.2011.11.001}; arXiv:1110.4210.

\bibitem{FS82}
J.~Fr\"ohlich and T.~Spencer, \emph{Massless phases and symmetry restoration in abelian gauge theories and spin systems},
Comm. Math. Phys. \textbf{83} (1982), 411--454.
\url{https://doi.org/10.1007/BF01213610}.

\bibitem{GM82}
M.~G\"opfert and G.~Mack, \emph{Proof of confinement of static quarks in 3-dimensional U(1) lattice gauge theory for all values of the coupling constant},
Comm. Math. Phys. \textbf{82} (1982), no.~4, 545--606.
\url{https://doi.org/10.1007/BF01961240}.

\bibitem{Guth80}
A.~H.~Guth, \emph{Existence proof of a nonconfining phase in four-dimensional U(1) lattice gauge theory},
Phys. Rev. D \textbf{21} (1980), no.~8, 2291--2307.
\url{https://doi.org/10.1103/PhysRevD.21.2291}.

\bibitem{IS07}
K.~R.~Ito and E.~Seiler, \emph{On the recent paper on quark confinement by Tomboulis},
preprint (2007), arXiv:0711.4930.

\bibitem{IS08}
K.~R.~Ito and E.~Seiler, \emph{Further discussion of Tomboulis' approach to the confinement problem},
preprint (2008), arXiv:0803.3019.

\bibitem{JW}
A.~M.~Jaffe and E.~Witten, \emph{Quantum Yang--Mills theory}, official problem description,
Clay Mathematics Institute, \url{https://www.claymath.org/wp-content/uploads/2022/06/yangmills.pdf}
(accessed 23 November 2026).

\bibitem{KQR21}
P.~Kravchuk, J.~Qiao, and S.~Rychkov, \emph{Distributions in CFT. Part II. Minkowski space},
J. High Energy Phys. \textbf{08} (2021), 094.
\url{https://doi.org/10.1007/JHEP08(2021)094}; arXiv:2104.02090.

\bibitem{KT02}
T.~G.~Kov\'acs and E.~T.~Tomboulis, \emph{Vortex waistlines and long range fluctuations},
Phys. Rev. D \textbf{65} (2002), 074501.
\url{https://doi.org/10.1103/PhysRevD.65.074501}; arXiv:hep-lat/0108017.

\bibitem{LTW04}
B.~Lucini, M.~Teper, and U.~Wenger, \emph{Glueballs and $k$-strings in SU($N$) gauge theories: calculations with improved operators},
J. High Energy Phys. \textbf{06} (2004), 012.
\url{https://doi.org/10.1088/1126-6708/2004/06/012}; arXiv:hep-lat/0404008.

\bibitem{Lu77}
M.~L\"uscher, \emph{Construction of a selfadjoint, strictly positive transfer matrix for Euclidean lattice gauge theories},
Comm. Math. Phys. \textbf{54} (1977), 283--292.
\url{https://doi.org/10.1007/BF01614090}.

\bibitem{LW86}
M.~L\"uscher, \emph{Volume dependence of the energy spectrum in massive quantum field theories. II. Scattering states},
Comm. Math. Phys. \textbf{105} (1986), no.~2, 153--188.
\url{https://doi.org/10.1007/BF01211097}.

\bibitem{OS73}
K.~Osterwalder and R.~Schrader, \emph{Axioms for Euclidean Green's functions},
Comm. Math. Phys. \textbf{31} (1973), no.~2, 83--112.
\url{https://doi.org/10.1007/BF01645738}.

\bibitem{OS75}
K.~Osterwalder and R.~Schrader, \emph{Axioms for Euclidean Green's functions II},
Comm. Math. Phys. \textbf{42} (1975), no.~3, 281--305.
\url{https://doi.org/10.1007/BF01608978}.

\bibitem{OS78}
K.~Osterwalder and E.~Seiler, \emph{Gauge field theories on a lattice},
Ann. Physics \textbf{110} (1978), no.~2, 440--471.
\url{https://doi.org/10.1016/0003-4916(78)90039-8}.

\bibitem{Repo} K. Justus, \emph{Millennium-Prize-Problems}, research repository, \url{https://github.com/KirtusJ/Millennium-Prize-Problems}, folder \texttt{yang-mills/z4/attempt} at commit \texttt{d801091} (23 November 2026).

\bibitem{RepoB4} Same repository, folder \texttt{yang-mills/transfer-operator/periodic-time/packet-bands} at commit \texttt{4e29b71}.

\bibitem{RepoC4} Same repository, folder \texttt{yang-mills/central-quotients/tangent-curvature/periodic-limit/block} at commit \texttt{16175fe}.

\bibitem{RepoThread} Same repository, issue 1 (coordination record), comments of 23 November 2026.

\bibitem{SZZ23}
H.~Shen, R.~Zhu, and X.~Zhu, \emph{A stochastic analysis approach to lattice Yang--Mills at strong coupling},
Comm. Math. Phys. \textbf{400} (2023), no.~2, 805--851.
\url{https://doi.org/10.1007/s00220-022-04609-1}; arXiv:2204.12737.

\bibitem{Te98}
M.~Teper, \emph{Glueball masses and other physical properties of SU(N) gauge theories in D=3+1: a review of lattice results for theorists},
preprint (1998), arXiv:hep-th/9812187.

\bibitem{To07}
E.~T.~Tomboulis, \emph{Confinement for all values of the coupling in four-dimensional SU(2) gauge theory},
preprint (2007), arXiv:0707.2179.

\bibitem{TY85}
E.~T.~Tomboulis and L.~G.~Yaffe, \emph{Finite temperature SU(2) lattice gauge theory},
Comm. Math. Phys. \textbf{100} (1985), 313--341.
\url{https://doi.org/10.1007/BF01206134}.

\end{thebibliography}
\end{document}
